%% filedir.tex
%%
%% Copyright (C) 2000-2018 Simone Piccardi.  Permission is granted to
%% copy, distribute and/or modify this document under the terms of the GNU Free
%% Documentation License, Version 1.1 or any later version published by the
%% Free Software Foundation; with the Invariant Sections being "Un preambolo",
%% with no Front-Cover Texts, and with no Back-Cover Texts.  A copy of the
%% license is included in the section entitled "GNU Free Documentation
%% License".
%%

\chapter{La gestione di file e directory}
\label{cha:files_and_dirs}

In questo capitolo tratteremo in dettaglio le modalità con cui si gestiscono
file e directory, iniziando da un approfondimento dell'architettura del
sistema illustrata a grandi linee in sez.~\ref{sec:file_arch_overview} ed
illustrando le principali caratteristiche di un filesystem e le interfacce
che consentono di controllarne il montaggio e lo smontaggio. 

Esamineremo poi le funzioni di libreria che si usano per copiare, spostare e
cambiare i nomi di file e directory e l'interfaccia che permette la
manipolazione dei loro attributi. Tratteremo inoltre la struttura di base del
sistema delle protezioni e del controllo dell'accesso ai file e le successive
estensioni (\textit{Extended Attributes}, ACL, quote disco). Tutto quello che
riguarda invece la gestione dell'I/O sui file è lasciato al capitolo
successivo.



\section{L'architettura della gestione dei file}
\label{sec:file_arch_func}

In questa sezione tratteremo con maggiori dettagli rispetto a quanto visto in
sez.~\ref{sec:file_arch_overview} il \textit{Virtual File System} di Linux e
come il kernel può gestire diversi tipi di filesystem, descrivendo prima le
caratteristiche generali di un filesystem di un sistema unix-like, per poi
fare una panoramica sul filesystem tradizionalmente più usato con Linux,
l'\acr{ext2} ed i suoi successori.


\subsection{Il funzionamento del \textit{Virtual File System} di Linux}
\label{sec:file_vfs_work}

% NOTE articolo interessante:
% http://www.ibm.com/developerworks/linux/library/l-virtual-filesystem-switch/index.html?ca=dgr-lnxw97Linux-VFSdth-LXdW&S_TACT=105AGX59&S_CMP=GRlnxw97

\itindbeg{Virtual~File~System~(VFS)}

Come illustrato brevemente in sez.~\ref{sec:file_arch_overview} in Linux il
concetto di \textit{everything is a file} è stato implementato attraverso il
\textit{Virtual File System}, la cui struttura generale è illustrata in
fig.~\ref{fig:file_VFS_scheme}.  Il VFS definisce un insieme di funzioni che
tutti i filesystem devono implementare per l'accesso ai file che contengono e
l'interfaccia che consente di eseguire l'I/O sui file, che questi siano di
dati o dispositivi. 

\itindbeg{inode}

L'interfaccia fornita dal VFS comprende in sostanza tutte le funzioni che
riguardano i file, le operazioni implementate dal VFS sono realizzate con una
astrazione che prevede quattro tipi di oggetti strettamente correlati: i
filesystem, le \textit{dentry}, gli \textit{inode} ed i file. A questi oggetti
corrispondono una serie di apposite strutture definite dal kernel che
contengono come campi le funzioni di gestione e realizzano l'infrastruttura
del VFS. L'interfaccia è molto complessa, ne faremo pertanto una trattazione
estremamente semplificata che consenta di comprenderne i principi
di funzionamento.

Il VFS usa una tabella mantenuta dal kernel che contiene il nome di ciascun
filesystem supportato, quando si vuole inserire il supporto di un nuovo
filesystem tutto quello che occorre è chiamare la funzione
\code{register\_filesystem} passando come argomento la struttura
\kstruct{file\_system\_type} (la cui definizione è riportata in
fig.~\ref{fig:kstruct_file_system_type}) relativa a quel filesystem. Questa
verrà inserita nella tabella, ed il nuovo filesystem comparirà in
\procfile{/proc/filesystems}.

\begin{figure}[!htb]
  \footnotesize \centering
  \begin{minipage}[c]{0.80\textwidth}
    \includestruct{listati/file_system_type.h}
  \end{minipage}
  \normalsize 
  \caption{Estratto della struttura \kstructd{file\_system\_type} usata dal
    VFS (da \texttt{include/linux/fs.h}).}
  \label{fig:kstruct_file_system_type}
\end{figure}

La struttura \kstruct{file\_system\_type}, oltre ad una serie di dati interni,
come il nome del tipo di filesystem nel campo \var{name},\footnote{quello che
  viene riportato in \procfile{/proc/filesystems} e che viene usato come
  valore del parametro dell'opzione \texttt{-t} del comando \texttt{mount} che
  indica il tipo di filesystem.}  contiene i riferimenti alle funzioni di base
che consentono l'utilizzo di quel filesystem. In particolare la funzione
\code{mount} del quarto campo è quella che verrà invocata tutte le volte che
si dovrà effettuare il montaggio di un filesystem di quel tipo. Per ogni nuovo
filesystem si dovrà allocare una di queste strutture ed inizializzare i
relativi campi con i dati specifici di quel filesystem, ed in particolare si
dovrà creare anche la relativa versione della funzione \code{mount}.

\itindbeg{pathname}
\itindbeg{pathname~resolution}

Come illustrato in fig.~\ref{fig:kstruct_file_system_type} questa funzione
restituisce una \textit{dentry}, abbreviazione che sta per \textit{directory
  entry}. Le \textit{dentry} sono gli oggetti che il kernel usa per eseguire
la \textit{pathname resolution}, ciascuna di esse corrisponde ad un
\textit{pathname} e contiene il riferimento ad un \textit{inode}, che come
vedremo a breve è l'oggetto usato dal kernel per identificare un
file.\footnote{in questo caso si parla di file come di un qualunque oggetto
  generico che sta sul filesystem e non dell'oggetto file del VFS cui
  accennavamo prima.} La \textit{dentry} ottenuta dalla chiamata alla funzione
\code{mount} sarà inserita in corrispondenza al \textit{pathname} della
directory in cui il filesystem è stato montato.

% NOTA: struct dentry è dichiarata in include/linux/dcache.h

Le \textit{dentry} sono oggetti del VFS che vivono esclusivamente in memoria,
nella cosiddetta \textit{directory entry cache} (spesso chiamata in breve
\textit{dcache}). Ogni volta che una \textit{system call} specifica un
\textit{pathname} viene effettuata una ricerca nella \textit{dcache} per
ottenere immediatamente la \textit{dentry} corrispondente,\footnote{il buon
  funzionamento della \textit{dcache} è in effetti di una delle parti più
  critiche per le prestazioni del sistema.} che a sua volta ci darà, tramite
l'\textit{inode}, il riferimento al file.

Dato che normalmente non è possibile mantenere nella \textit{dcache} le
informazioni relative a tutto l'albero dei file la procedura della
\textit{pathname resolution} richiede un meccanismo con cui riempire gli
eventuali vuoti. Il meccanismo prevede che tutte le volte che si arriva ad una
\textit{dentry} mancante venga invocata la funzione \texttt{lookup}
dell'\textit{inode} associato alla \textit{dentry} precedente nella
risoluzione del \textit{pathname},\footnote{che a questo punto è una
  directory, per cui si può cercare al suo interno il nome di un file.} il cui
scopo è risolvere il nome mancante e fornire la sua \textit{dentry} che a
questo punto verrà inserita nella cache.

Dato che tutte le volte che si monta un filesystem la funzione \texttt{mount}
(vedi sez.~\ref{sec:filesystem_mounting}) della corrispondente
\kstruct{file\_system\_type} inserisce la \textit{dentry} iniziale nel
\textit{mount point} dello stesso, si avrà comunque un punto di
partenza. Inoltre essendo questa \textit{dentry} relativa a quel tipo di
filesystem essa farà riferimento ad un \textit{inode} di quel filesystem, e
come vedremo questo farà sì che venga eseguita una \texttt{lookup} adatta per
effettuare la risoluzione dei nomi per quel filesystem.

\itindend{pathname}
\itindend{pathname~resolution}

% Un secondo effetto della chiamata funzione \texttt{mount} di
% \kstruct{file\_system\_type} è quello di allocare una struttura
% \kstruct{super\_block} per ciascuna istanza montata, che contiene le
% informazioni generali di un qualunque filesystem montato, come le opzioni di
% montaggio, le dimensioni dei blocchi, quando il filesystem è stato montato
% ecc. Fra queste però viene pure inserta, nel campo \var{s\_op}, una ulteriore
% struttura \kstruct{super\_operations}, il cui contenuto sono i puntatori
% alle funzioni di gestione di un filesystem, anche inizializzata in modo da
% utilizzare le versioni specifiche di quel filesystem.

L'oggetto più importante per il funzionamento del VFS è probabilmente
l'\textit{inode}, ma con questo nome si può fare riferimento a due cose
diverse.  La prima è la struttura su disco (su cui torneremo anche in
sez.~\ref{sec:file_filesystem}) che fa parte della organizzazione dei dati
realizzata dal filesystem e che contiene le informazioni relative alle
proprietà (i cosiddetti \textsl{metadati}) di ogni oggetto presente su di esso
(si intende al solito uno qualunque dei tipi di file di
tab.~\ref{tab:file_file_types}).

La seconda è la corrispondente struttura \kstruct{inode}, della cui
definizione si è riportato un estratto in
fig.~\ref{fig:kstruct_inode}.\footnote{l'estratto fa riferimento alla versione
  del kernel 2.6.37.} Questa struttura viene mantenuta in memoria ed è a
questa che facevamo riferimento quando parlavamo dell'\textit{inode} associato
a ciascuna \textit{dentry}. Nella struttura in memoria sono presenti gli
stessi \textsl{metadati} memorizzati su disco, che vengono letti quando questa
struttura viene allocata e trascritti all'indietro se modificati.

\begin{figure}[!htb]
  \footnotesize \centering
  \begin{minipage}[c]{0.8\textwidth}
    \includestruct{listati/inode.h}
  \end{minipage}
  \normalsize 
  \caption{Estratto della struttura \kstructd{inode} del kernel (da
    \texttt{include/linux/fs.h}).}
  \label{fig:kstruct_inode}
\end{figure}

Il fatto che la struttura \kstruct{inode} sia mantenuta in memoria,
direttamente associata ad una \textit{dentry}, rende sostanzialmente immediate
le operazioni che devono semplicemente effettuare un accesso ai dati in essa
contenuti: è così ad esempio che viene realizzata la \textit{system call}
\func{stat} che vedremo in sez.~\ref{sec:file_stat}. Rispetto ai dati salvati
sul disco questa struttura contiene però anche quanto necessario alla
implementazione del VFS, ed in particolare è importante il campo \var{i\_op}
che, come illustrato in fig.~\ref{fig:kstruct_inode}, contiene il puntatore ad
una struttura di tipo \kstruct{inode\_operation}, la cui definizione si può
trovare nel file \texttt{include/kernel/fs.h} dei sorgenti del kernel.

Questa struttura non è altro che una tabella di funzioni, ogni suo membro cioè
è un puntatore ad una funzione e, come suggerisce il nome della struttura
stessa, queste funzioni sono quelle che definiscono le operazioni che il VFS
può compiere su un \textit{inode}. Si sono riportate in
tab.~\ref{tab:file_inode_operations} le più rilevanti.

\begin{table}[htb]
  \centering
  \footnotesize
  \begin{tabular}[c]{|l|l|}
    \hline
    \textbf{Funzione} & \textbf{Operazione} \\
    \hline
    \hline
    \textsl{\code{create}} & Chiamata per creare un nuovo file (vedi
                             sez.~\ref{sec:file_open_close}).\\ 
    \textsl{\code{link}}   & Crea un \textit{hard link} (vedi
                             sez.~\ref{sec:link_symlink_rename}).\\
    \textsl{\code{unlink}} & Cancella un \textit{hard link} (vedi
                             sez.~\ref{sec:link_symlink_rename}).\\
    \textsl{\code{symlink}}& Crea un collegamento simbolico (vedi
                             sez.~\ref{sec:link_symlink_rename}).\\
    \textsl{\code{mkdir}}  & Crea una directory (vedi
                             sez.~\ref{sec:file_dir_creat_rem}).\\
    \textsl{\code{rmdir}}  & Rimuove una directory (vedi
                             sez.~\ref{sec:file_dir_creat_rem}).\\
    \textsl{\code{mknod}}  & Crea un file speciale (vedi
                             sez.~\ref{sec:file_mknod}).\\
    \textsl{\code{rename}} & Cambia il nome di un file (vedi
                             sez.~\ref{sec:link_symlink_rename}).\\
    \textsl{\code{lookup}}&  Risolve il nome di un file.\\
    \hline
  \end{tabular}
  \caption{Le principali operazioni sugli \textit{inode} definite tramite
    \kstructd{inode\_operation}.} 
  \label{tab:file_inode_operations}
\end{table}

Possiamo notare come molte di queste funzioni abbiano nomi sostanzialmente
identici alle varie \textit{system call} con le quali si gestiscono file e
directory, che tratteremo nel resto del capitolo. Quello che succede è che
tutte le volte che deve essere eseguita una \textit{system call}, o una
qualunque altra operazione su un \textit{inode} (come \texttt{lookup}) il VFS
andrà ad utilizzare la funzione corrispondente attraverso il puntatore
\var{i\_op}.

Sarà allora sufficiente che nella realizzazione di un filesystem si crei una
implementazione di queste funzioni per quel filesystem e si allochi una
opportuna istanza di \kstruct{inode\_operation} contenente i puntatori a dette
funzioni. A quel punto le strutture \kstruct{inode} usate per gli oggetti di
quel filesystem otterranno il puntatore alla relativa istanza di
\kstruct{inode\_operation} e verranno automaticamente usate le funzioni
corrette.

Si noti però come in tab.~\ref{tab:file_inode_operations} non sia presente la
funzione \texttt{open} che invece è citata in
tab.~\ref{tab:file_file_operations}.\footnote{essa può essere comunque
  invocata dato che nella struttura \kstruct{inode} è presente anche il
  puntatore \var{i\_fop} alla struttura \kstruct{file\_operation} che fornisce
  detta funzione.} Questo avviene perché su Linux l'apertura di un file
richiede comunque un'altra operazione che mette in gioco l'omonimo oggetto del
VFS: l'allocazione di una struttura di tipo \kstruct{file} che viene associata
ad ogni file aperto nel sistema.  I motivi per cui viene usata una struttura a
parte sono diversi, anzitutto, come illustrato in sez.~\ref{sec:file_fd},
questa è necessaria per le operazioni eseguite dai processi con l'interfaccia
dei file descriptor. Ogni processo infatti mantiene il riferimento ad una
struttura \kstruct{file} per ogni file che ha aperto, ed è tramite essa che
esegue le operazioni di I/O. Inoltre il kernel mantiene un elenco di tutti i
file aperti nella \textit{file table} (torneremo su questo in
sez.~\ref{sec:file_fd}).

Inoltre se le operazioni relative agli \textit{inode} fanno riferimento ad
oggetti posti all'interno di un filesystem e vi si applicano quindi le
funzioni fornite nell'implementazione di quest'ultimo, quando si apre un file
questo può essere anche un file di dispositivo, ed in questo caso il VFS
invece di usare le operazioni fornite dal filesystem (come farebbe per un file
di dati) dovrà invece ricorrere a quelle fornite dal driver del dispositivo.

\itindend{inode}

\begin{figure}[!htb]
  \footnotesize \centering
  \begin{minipage}[c]{0.8\textwidth}
    \includestruct{listati/file.h}
  \end{minipage}
  \normalsize 
  \caption{Estratto della struttura \kstructd{file} del kernel (da
    \texttt{include/linux/fs.h}).}
  \label{fig:kstruct_file}
\end{figure}

Come si può notare dall'estratto di fig.~\ref{fig:kstruct_file}, la struttura
\kstruct{file} contiene, oltre ad alcune informazioni usate dall'interfaccia
dei file descriptor il cui significato emergerà più avanti (vedi
sez.~\ref{sec:file_unix_interface}), il puntatore \var{f\_op} ad una struttura
\kstruct{file\_operation}. Questa è l'analoga per i file di
\kstruct{inode\_operation}, e definisce le operazioni generiche fornite dal
VFS per i file. Si sono riportate in tab.~\ref{tab:file_file_operations} le
più significative.

\begin{table}[htb]
  \centering
  \footnotesize
  \begin{tabular}[c]{|l|p{8cm}|}
    \hline
    \textbf{Funzione} & \textbf{Operazione} \\
    \hline
    \hline
    \textsl{\code{open}}   & Apre il file (vedi
                             sez.~\ref{sec:file_open_close}).\\ 
    \textsl{\code{read}}   & Legge dal file (vedi sez.~\ref{sec:file_read}).\\
    \textsl{\code{write}}  & Scrive sul file (vedi 
                             sez.~\ref{sec:file_write}).\\
    \textsl{\code{llseek}} & Sposta la posizione corrente sul file (vedi
                             sez.~\ref{sec:file_lseek}).\\
    \textsl{\code{ioctl}}  & Accede alle operazioni di controllo 
                             (vedi sez.~\ref{sec:file_fcntl_ioctl}).\\
    \textsl{\code{readdir}}& Legge il contenuto di una directory (vedi 
                             sez.~\ref{sec:file_dir_read}).\\
    \textsl{\code{poll}}   & Usata nell'I/O multiplexing (vedi
                             sez.~\ref{sec:file_multiplexing}).\\
    \textsl{\code{mmap}}   & Mappa il file in memoria (vedi 
                             sez.~\ref{sec:file_memory_map}).\\
    \textsl{\code{release}}& Chiamata quando l'ultimo riferimento a un file 
                             aperto è chiuso.\\
    \textsl{\code{fsync}}  & Sincronizza il contenuto del file (vedi
                             sez.~\ref{sec:file_sync}).\\
    \textsl{\code{fasync}} & Abilita l'I/O asincrono (vedi
                             sez.~\ref{sec:file_asyncronous_io}) sul file.\\
    \hline
  \end{tabular}
  \caption{Operazioni sui file definite tramite \kstructd{file\_operation}.}
  \label{tab:file_file_operations}
\end{table}

Anche in questo caso tutte le volte che deve essere eseguita una
\textit{system call} o una qualunque altra operazione sul file il VFS andrà ad
utilizzare la funzione corrispondente attraverso il puntatore
\var{f\_op}. Dato che è cura del VFS quando crea la struttura all'apertura del
file assegnare a \var{f\_op} il puntatore alla versione di
\kstruct{file\_operation} corretta per quel file, sarà possibile scrivere allo
stesso modo sulla porta seriale come su un normale file di dati, e lavorare
sui file allo stesso modo indipendentemente dal filesystem.

Il VFS realizza la quasi totalità delle operazioni relative ai file grazie
alle funzioni presenti nelle due strutture \kstruct{inode\_operation} e
\kstruct{file\_operation}.  Ovviamente non è detto che tutte le operazioni
possibili siano poi disponibili in tutti i casi, ad esempio \code{llseek} non
sarà presente per un dispositivo come la porta seriale o per una
\textit{fifo}, mentre sui file del filesystem \texttt{vfat} non saranno
disponibili i permessi, ma resta il fatto che grazie al VFS le \textit{system
  call} per le operazioni sui file possono restare sempre le stesse nonostante
le enormi differenze che possono esserci negli oggetti a cui si applicano.
 

\itindend{Virtual~File~System~(VFS)}

% NOTE: documentazione interessante:
%       * sorgenti del kernel: Documentation/filesystems/vfs.txt
%       * http://thecoffeedesk.com/geocities/rkfs.html
%       * http://www.linux.it/~rubini/docs/vfs/vfs.html



\subsection{Il funzionamento di un filesystem Unix}
\label{sec:file_filesystem}

Come già accennato in sez.~\ref{sec:file_arch_overview} Linux (ed ogni sistema
unix-like) organizza i dati che tiene su disco attraverso l'uso di un
filesystem. Una delle caratteristiche di Linux rispetto agli altri Unix è
quella di poter supportare, grazie al VFS, una enorme quantità di filesystem
diversi, ognuno dei quali avrà una sua particolare struttura e funzionalità
proprie.  Per questo non entreremo nei dettagli di un filesystem specifico, ma
daremo una descrizione a grandi linee che si adatta alle caratteristiche
comuni di qualunque filesystem di un sistema unix-like.

\itindbeg{superblock}

Una possibile strutturazione dell'informazione su un disco è riportata in
fig.~\ref{fig:file_disk_filesys}, dove si hanno tre filesystem su tre
partizioni. In essa per semplicità si è fatto riferimento alla struttura del
filesystem \acr{ext2}, che prevede una suddivisione dei dati in \textit{block
  group}.  All'interno di ciascun \textit{block group} viene anzitutto
replicato il cosiddetto \textit{superblock}, (la struttura che contiene
l'indice iniziale del filesystem e che consente di accedere a tutti i dati
sottostanti) e creata una opportuna suddivisione dei dati e delle informazioni
per accedere agli stessi.  Sulle caratteristiche di \acr{ext2} e derivati
torneremo in sez.~\ref{sec:file_ext2}.

\itindend{superblock}
\itindbeg{inode}

È comunque caratteristica comune di tutti i filesystem per Unix,
indipendentemente da come poi viene strutturata nei dettagli questa
informazione, prevedere la presenza di due tipi di risorse: gli
\textit{inode}, cui abbiamo già accennato in sez.~\ref{sec:file_vfs_work}, che
sono le strutture che identificano i singoli oggetti sul filesystem, e i
blocchi, che invece attengono allo spazio disco che viene messo a disposizione
per i dati in essi contenuti.

\begin{figure}[!htb]
  \centering
  \includegraphics[width=11cm]{img/disk_struct}
  \caption{Organizzazione dello spazio su un disco in partizioni e
  filesystem.}
  \label{fig:file_disk_filesys}
\end{figure}

Se si va ad esaminare con maggiore dettaglio la strutturazione
dell'informazione all'interno del filesystem \acr{ext2}, tralasciando i
dettagli relativi al funzionamento del filesystem stesso come la
strutturazione in gruppi dei blocchi, il \textit{superblock} e tutti i dati di
gestione possiamo esemplificare la situazione con uno schema come quello
esposto in fig.~\ref{fig:file_filesys_detail}.

\begin{figure}[!htb]
  \centering
  \includegraphics[width=11cm]{img/filesys_struct}
  \caption{Strutturazione dei dati all'interno di un filesystem.}
  \label{fig:file_filesys_detail}
\end{figure}

Da fig.~\ref{fig:file_filesys_detail} si evidenziano alcune delle
caratteristiche di base di un filesystem, che restano le stesse anche su
filesystem la cui organizzazione dei dati è totalmente diversa da quella
illustrata, e sulle quali è bene porre attenzione visto che sono fondamentali
per capire il funzionamento delle funzioni che manipolano i file e le
directory che tratteremo nel prosieguo del capitolo. In particolare è
opportuno tenere sempre presente che:


\begin{enumerate*}
  
\item L'\textit{inode} contiene i cosiddetti \textsl{metadati}, vale dire le
  informazioni riguardanti le proprietà del file come oggetto del filesystem:
  il tipo di file, i permessi di accesso, le dimensioni, i puntatori ai
  blocchi fisici che contengono i dati e così via. Le informazioni che la
  funzione \func{stat} (vedi sez.~\ref{sec:file_stat}) fornisce provengono
  dall'\textit{inode}.  Dentro una directory si troverà solo il nome del file
  e il numero dell'\textit{inode} ad esso associato; il nome non è una
  proprietà del file e non viene mantenuto nell'\textit{inode}. Da da qui in
  poi chiameremo il nome del file contenuto in una directory
  ``\textsl{voce}'', come traduzione della nomenclatura inglese
  \textit{directory entry} che non useremo per evitare confusione con le
  \textit{dentry} del kernel viste in sez.~\ref{sec:file_vfs_work}.
  
\item Come mostrato in fig.~\ref{fig:file_filesys_detail} per i file
  \texttt{macro.tex} e \texttt{gapil\_macro.tex}, ci possono avere più voci
  che fanno riferimento allo stesso \textit{inode}. Fra le proprietà di un
  file mantenute nell'\textit{inode} c'è anche il contatore con il numero di
  riferimenti che sono stati fatti ad esso, il cosiddetto \textit{link
    count}.\footnote{mantenuto anche nel campo \var{i\_nlink} della struttura
    \kstruct{inode} di fig.~\ref{fig:kstruct_inode}.}  Solo quando questo
  contatore si annulla i dati del file possono essere effettivamente rimossi
  dal disco. Per questo la funzione per cancellare un file si chiama
  \func{unlink} (vedi sez.~\ref{sec:link_symlink_rename}), ed in realtà non
  cancella affatto i dati del file, ma si limita ad eliminare la relativa voce
  da una directory e decrementare il numero di riferimenti
  nell'\textit{inode}.
  
\item All'interno di ogni filesystem ogni \textit{inode} è identificato da un
  numero univoco. Il numero di \textit{inode} associato ad una voce in una
  directory si riferisce ad questo numero e non ci può essere una directory
  che contiene riferimenti ad \textit{inode} relativi ad altri filesystem.
  Questa è la ragione che limita l'uso del comando \cmd{ln}, che crea una
  nuova voce per un file esistente con la funzione \func{link} (vedi
  sez.~\ref{sec:link_symlink_rename}), a operare su file nel filesystem
  corrente.
  
\item Quando si cambia nome ad un file senza cambiare filesystem il contenuto
  del file non viene spostato fisicamente, viene semplicemente creata una
  nuova voce per l'\textit{inode} in questione e rimossa la precedente, questa
  è la modalità in cui opera normalmente il comando \cmd{mv} attraverso la
  funzione \func{rename} (vedi sez.~\ref{sec:link_symlink_rename}). Questa
  operazione non modifica minimamente neanche l'\textit{inode} del file, dato
  che non si opera sul file ma sulla directory che lo contiene.

\item Gli \textit{inode} dei file, che contengono i \textsl{metadati}, ed i
  blocchi di spazio disco, che contengono i dati, sono risorse indipendenti ed
  in genere vengono gestite come tali anche dai diversi filesystem; è pertanto
  possibile esaurire sia lo spazio disco (il caso più comune) che lo spazio
  per gli \textit{inode}. Nel primo caso non sarà possibile allocare ulteriore
  spazio, ma si potranno creare file (vuoti), nel secondo non si potranno
  creare nuovi file, ma si potranno estendere quelli che ci
  sono.\footnote{questo comportamento non è generale, alcuni filesystem più
    sofisticati possono evitare il problema dell'esaurimento degli
    \textit{inode} riallocando lo spazio disco libero per i blocchi.}

\end{enumerate*}

\begin{figure}[!htb]
  \centering 
  \includegraphics[width=12cm]{img/dir_links}
  \caption{Organizzazione dei \textit{link} per le directory.}
  \label{fig:file_dirs_link}
\end{figure}

Infine tenga presente che, essendo file pure loro, il numero di riferimenti
esiste anche per le directory. Per questo se a partire dalla situazione
mostrata in fig.~\ref{fig:file_filesys_detail} creiamo una nuova directory
\file{img} nella directory \file{gapil}, avremo una situazione come quella
illustrata in fig.~\ref{fig:file_dirs_link}.

La nuova directory avrà un numero di riferimenti pari a due, in quanto è
referenziata dalla directory da cui si era partiti (in cui è inserita la nuova
voce che fa riferimento a \texttt{img}) e dalla voce interna ``\texttt{.}''
che è presente in ogni directory.  Questo è il valore che si troverà sempre
per ogni directory che non contenga a sua volta altre directory. Al contempo,
la directory da cui si era partiti avrà un numero di riferimenti di almeno
tre, in quanto adesso sarà referenziata anche dalla voce ``\texttt{..}'' di
\texttt{img}. L'aggiunta di una sottodirectory fa cioè crescere di uno il
\textit{link count} della directory genitrice.

\itindend{inode}


\subsection{Alcuni dettagli sul filesystem \acr{ext2} e successori}
\label{sec:file_ext2}

Benché non esista ``il'' filesystem di Linux, dato che esiste un supporto
nativo di diversi filesystem che sono in uso da anni, quello che gli avvicina
di più è la famiglia di filesystem evolutasi a partire dal \textit{second
  extended filesystem}, o \acr{ext2}. Il filesystem \acr{ext2} ha subito un
grande sviluppo e diverse evoluzioni, fra cui l'aggiunta del
\textit{journaling} con il passaggio ad \acr{ext3}, che probabilmente è ancora
il filesystem più diffuso, ed una serie di ulteriori miglioramenti con il
successivo \acr{ext4}. In futuro è previsto che questo debba essere sostituito
da un filesystem completamente diverso, \acr{Btrfs}, che dovrebbe diventare il
filesystem standard di Linux, ma questo al momento è ancora in fase di
sviluppo.\footnote{si fa riferimento al momento dell'ultima revisione di
  questo paragrafo, l'inizio del 2012.}

Il filesystem \acr{ext2} nasce come filesystem nativo per Linux a partire
dalle prime versioni del kernel e supporta tutte le caratteristiche di un
filesystem standard Unix: è in grado di gestire nomi di file lunghi (256
caratteri, estensibili a 1012) e supporta una dimensione massima dei file fino
a 4~Tb. I successivi filesystem \acr{ext3} ed \acr{ext4} sono evoluzioni di
questo filesystem, e sia pure con molti miglioramenti ed estensioni
significative ne mantengono le caratteristiche fondamentali.

Oltre alle caratteristiche standard, \acr{ext2} fornisce alcune estensioni che
non sono presenti su un classico filesystem di tipo Unix; le principali sono
le seguenti:
\begin{itemize*}
\item gli attributi estesi (vedi sez.~\ref{sec:file_xattr}) che consentono di
  estendere le informazioni salvabili come metadati e le ACL (vedi
  sez.~\ref{sec:file_ACL}) che consentono di estendere il modello tradizionale
  dei permessi sui file.
\item sono supportate entrambe le semantiche di BSD e SVr4 come opzioni di
  montaggio. La semantica BSD comporta che i file in una directory sono creati
  con lo stesso identificatore di gruppo della directory che li contiene. La
  semantica SVr4 comporta che i file vengono creati con l'identificatore del
  gruppo primario del processo, eccetto il caso in cui la directory ha il bit
  di \acr{sgid} impostato (per una descrizione dettagliata del significato di
  questi termini si veda sez.~\ref{sec:file_access_control}), nel qual caso
  file e subdirectory ereditano sia il \ids{GID} che lo \acr{sgid}.
\item l'amministratore può scegliere la dimensione dei blocchi del filesystem
  in fase di creazione, a seconda delle sue esigenze: blocchi più grandi
  permettono un accesso più veloce, ma sprecano più spazio disco.
\item il filesystem implementa collegamenti simbolici veloci, in cui il nome
  del file non è salvato su un blocco, ma tenuto all'interno
  dell'\textit{inode} (evitando letture multiple e spreco di spazio), non
  tutti i nomi però possono essere gestiti così per limiti di spazio (il
  limite è 60 caratteri).
\item vengono supportati i cosiddetti \textit{file attributes} (vedi
  sez.~\ref{sec:file_perm_overview}) che attivano comportamenti specifici per
  i file su cui vengono attivati come marcarli come immutabili (che possono
  cioè essere soltanto letti) per la protezione di file di configurazione
  sensibili, o come \textit{append-only} (che possono essere aperti in
  scrittura solo per aggiungere dati) per la protezione dei file di log.
\end{itemize*}

La struttura di \acr{ext2} è stata ispirata a quella del filesystem di BSD: un
filesystem è composto da un insieme di blocchi, la struttura generale è quella
riportata in fig.~\ref{fig:file_filesys_detail}, in cui la partizione è divisa
in gruppi di blocchi.

Ciascun gruppo di blocchi contiene una copia delle informazioni essenziali del
filesystem (i \textit{superblock} sono quindi ridondati) per una maggiore
affidabilità e possibilità di recupero in caso di corruzione del
\textit{superblock} principale. L'utilizzo di raggruppamenti di blocchi ha
inoltre degli effetti positivi nelle prestazioni dato che viene ridotta la
distanza fra i dati e la tabella degli \textit{inode}.

\begin{figure}[!htb]
  \centering
  \includegraphics[width=9cm]{img/dir_struct}  
  \caption{Struttura delle directory nel \textit{second extended filesystem}.}
  \label{fig:file_ext2_dirs}
\end{figure}


Le directory sono implementate come una \textit{linked list} con voci di
dimensione variabile. Ciascuna voce della lista contiene il numero di
\textit{inode}, la sua lunghezza, il nome del file e la sua lunghezza, secondo
lo schema in fig.~\ref{fig:file_ext2_dirs}; in questo modo è possibile
implementare nomi per i file anche molto lunghi (fino a 1024 caratteri) senza
sprecare spazio disco.

Con l'introduzione del filesystem \acr{ext3} sono state introdotte diverse
modifiche strutturali, la principale di queste è quella che \acr{ext3} è un
filesystem \textit{journaled}, è cioè in grado di eseguire una registrazione
delle operazioni di scrittura su un giornale (uno speciale file interno) in
modo da poter garantire il ripristino della coerenza dei dati del
filesystem\footnote{si noti bene che si è parlato di dati \textsl{del}
  filesystem, non di dati \textsl{nel} filesystem, quello di cui viene
  garantito un veloce ripristino è relativo ai dati della struttura interna
  del filesystem, non di eventuali dati contenuti nei file che potrebbero
  essere stati persi.} in brevissimo tempo in caso di interruzione improvvisa
della corrente o di crollo del sistema che abbia causato una interruzione
della scrittura dei dati sul disco.

Oltre a questo \acr{ext3} introduce ulteriori modifiche volte a migliorare
sia le prestazioni che la semplicità di gestione del filesystem, in
particolare per le directory si è passato all'uso di alberi binari con
indicizzazione tramite \textit{hash} al posto delle \textit{linked list} che
abbiamo illustrato, ottenendo un forte guadagno di prestazioni in caso di
directory contenenti un gran numero di file.

% TODO (bassa priorità) portare a ext3, ext4 e btrfs ed illustrare le
% problematiche che si possono incontrare (in particolare quelle relative alla
% perdita di contenuti in caso di crash del sistema)
% TODO (media priorità) trattare btrfs quando sarà usato come stabile


\subsection{La gestione dell'uso dei filesystem}
\label{sec:filesystem_mounting}

Come accennato in sez.~\ref{sec:file_arch_overview} per poter accedere ai file
occorre rendere disponibile al sistema il filesystem su cui essi sono
memorizzati. L'operazione di attivazione del filesystem è chiamata
\textsl{montaggio} e per far questo in Linux si usa la funzione di sistema
\funcd{mount}, il cui prototipo è:\footnote{la funzione è una versione
  specifica di Linux che usa la omonima \textit{system call} e non è
  portabile.}

\begin{funcproto}{
\fhead{sys/mount.h} 
\fdecl{mount(const char *source, const char *target, const char
  *filesystemtype, \\ 
\phantom{mount(}unsigned long mountflags, const void *data)}
\fdesc{Monta un filesystem.} 
}
{La funzione ritorna $0$ in caso di successo e $-1$ per un errore, nel qual
  caso \var{errno} assumerà uno dei valori:
  \begin{errlist}
  \item[\errcode{EACCES}] non si ha il permesso di accesso su uno dei
    componenti del \textit{pathname}, o si è cercato di montare un filesystem
    disponibile in sola lettura senza aver specificato \const{MS\_RDONLY} o il
    device \param{source} è su un filesystem montato con l'opzione
    \const{MS\_NODEV}.
  \item[\errcode{EBUSY}] \param{source} è già montato, o non può essere
    rimontato in sola lettura perché ci sono ancora file aperti in scrittura,
    o non può essere montato su \param{target} perché la directory è ancora in
    uso.
  \item[\errcode{EINVAL}] il dispositivo \param{source} presenta un
    \textit{superblock} non valido, o si è cercato di rimontare un filesystem
    non ancora montato, o di montarlo senza che \param{target} sia un
    \textit{mount point} o di spostarlo quando \param{target} non è un
    \textit{mount point} o è la radice o si è usato un valore di
    \param{mountflags} non valido.
  \item[\errcode{ELOOP}] si è cercato di spostare un \textit{mount point} su
    una sottodirectory di \param{source} o si sono incontrati troppi
    collegamenti simbolici nella risoluzione di un nome.
  \item[\errcode{EMFILE}] in caso di filesystem virtuale, la tabella dei
    dispositivi fittizi (chiamati \textit{dummy} nella documentazione inglese)
    è piena.
  \item[\errcode{ENODEV}] il tipo \param{filesystemtype} non esiste o non è
    configurato nel kernel.
  \item[\errcode{ENOTBLK}] non si è usato un \textit{block device} per
    \param{source} quando era richiesto.
  \item[\errcode{ENXIO}] il \textit{major number} del
    dispositivo \param{source} è sbagliato.
  \item[\errcode{EPERM}] il processo non ha i privilegi di amministratore.
  \end{errlist} 
  ed inoltre \errval{EFAULT}, \errval{ENOMEM}, \errval{ENAMETOOLONG},
  \errval{ENOENT}, \errval{ENOTDIR} nel loro significato generico.}
\end{funcproto}

\itindbeg{mount~point}

L'uso più comune della funzione è quello di montare sulla directory indicata
da \param{target}, detta \textit{mount point}, il filesystem contenuto nel
file di dispositivo indicato da \param{source}. In entrambi i casi, come
daremo per assunto da qui in avanti tutte le volte che si parla di directory o
file nel passaggio di un argomento di una funzione, si intende che questi
devono essere indicati con la stringa contenente il loro \textit{pathname}.

Normalmente un filesystem è contenuto su un disco o una partizione, ma come
illustrato in sez.~\ref{sec:file_vfs_work} la struttura del \textit{Virtual
  File System} è estremamente flessibile e può essere usata anche per oggetti
diversi da un disco. Ad esempio usando il \textit{loop device} si può montare
un file qualunque (come l'immagine di un CD-ROM o di un floppy) che contiene
l'immagine di un filesystem, inoltre alcuni tipi di filesystem, come
\texttt{proc} o \texttt{sysfs} sono virtuali e non hanno un supporto che ne
contenga i dati che sono generati al volo dal kernel ad ogni lettura, e
inviati al kernel ad ogni scrittura (costituiscono quindi un meccanismo di
comunicazione, attraverso l'ordinaria interfaccia dei file, con il kernel).

Il tipo di filesystem che si vuole montare è specificato
dall'argomento \param{filesystemtype}, che deve essere una delle stringhe
riportate nel file \procfilem{/proc/filesystems} che, come accennato in
sez.~\ref{sec:file_vfs_work}, contiene l'elenco dei filesystem supportati dal
kernel. Nel caso si sia indicato un filesystem virtuale, che non è associato a
nessun file di dispositivo, il contenuto di \param{source} viene ignorato.

L'argomento \param{data} viene usato per passare le impostazioni relative alle
caratteristiche specifiche di ciascun filesystem. Si tratta di una stringa di
parole chiave (separate da virgole e senza spazi) che indicano le cosiddette
``\textsl{opzioni}'' del filesystem che devono essere impostate; in genere
viene usato direttamente il contenuto del parametro dell'opzione \texttt{-o}
del comando \texttt{mount}. I valori utilizzabili dipendono dal tipo di
filesystem e ciascuno ha i suoi, pertanto si rimanda alla documentazione della
pagina di manuale di questo comando e dei singoli filesystem.

Dopo l'esecuzione della funzione il contenuto del filesystem viene reso
disponibile nella directory specificata come \textit{mount point} ed il
precedente contenuto di detta directory viene mascherato dal contenuto della
directory radice del filesystem montato. Fino ai kernel della serie 2.2.x non
era possibile montare un filesystem se un \textit{mount point} era già in uso,
coi kernel successivi è possibile montare più filesystem sullo stesso
\textit{mount point} impilandoli l'uno sull'altro, anche in questo caso vale
quanto appena detto, e solo il contenuto dell'ultimo filesystem montato sarà
visibile, mascherando quelli sottostanti.

In realtà quella di montare un filesystem è solo una delle operazioni che si
possono effettuare con \func{mount}, la funzione infatti è dedicata a tutte le
operazioni relative alla gestione del montaggio dei filesystem e dei
\textit{mount point}. Ad esempio fin dalle sue origini poteva essere
utilizzata per effettuare il rimontaggio di un filesystem con opzioni diverse,
ed a partire dal kernel 2.4.x è divenuto possibile usarla per spostare
atomicamente un \textit{mount point} da una directory ad un'altra, per montare
lo stesso filesystem in diversi \textit{mount point}, per montare una
directory su un'altra (il cosiddetto \textit{bind mount}).

\itindend{mount~point}

Il tipo di operazione compiuto da \func{mount} viene stabilito in base al
valore dell'argomento \param{mountflags}, che oltre alla selezione del tipo di
operazione da compiere, consente anche di indicare alcune opzioni generiche
valide per qualunque filesystem.\footnote{benché queste siano espresse nel
  comando \cmd{mount} con l'opzione \texttt{-o} esse non vengono impostate nei
  valori di \param{data}, che serve solo per le opzioni specifiche di ogni
  filesystem.}  Il valore dell'argomento deve essere espresso come maschera
binaria e i vari bit che lo compongono, detti anche \textit{mount flags},
devono essere impostati con un OR aritmetico dei valori dalle opportune
costanti che illustreremo a breve.

In Linux \param{mountflags} deve essere un intero a 32 bit;
fino ai kernel della serie 2.2.x i 16 più significativi avevano un valore
riservato che doveva essere specificato obbligatoriamente,\footnote{il valore
  era il \textit{magic number} \code{0xC0ED}, si può usare la costante
  \constd{MS\_MGC\_MSK} per ottenere la parte di \param{mountflags} riservata
  al \textit{magic number}, mentre per specificarlo si può dare un OR
  aritmetico con la costante \constd{MS\_MGC\_VAL}.} e si potevano usare solo
i 16 meno significativi. Oggi invece, con un numero di opzioni superiore, sono
utilizzati tutti e 32 i bit, ma qualora nei 16 più significativi sia presente
detto valore, che non esprime una combinazione valida, esso viene ignorato.

Come accennato il tipo di operazione eseguito da \func{mount} viene stabilito
in base al contenuto di \param{mountflags}, la scelta viene effettuata
controllando nell'ordine:
\begin{enumerate*}
\item se è presente il flag \const{MS\_REMOUNT}, nel qual caso verrà eseguito
  il rimontaggio del filesystem, con le nuove opzioni indicate da \param{data}
  e dagli altri flag di \param{mountflags};
\item se è presente il flag \const{MS\_BIND}, nel qual caso verrà eseguito un
  \textit{bind mount} (argomento che tratteremo più avanti);
\item se è presente uno fra \const{MS\_SHARED}, \const{MS\_PRIVATE},
  \const{MS\_SLAVE}, \const{MS\_UNBINDABLE}, nel qual caso verrà cambiata la
  modalità di propagazione del montaggio (detti valori sono mutualmente
  esclusivi).
\item se è presente \const{MS\_MOVE}, nel qual caso verrà effettuato uno
  spostamento del \textit{mount point};
\item se nessuno dei precedenti è presente si tratta di una ordinaria
  operazione di montaggio di un filesystem.
\end{enumerate*}

Il fatto che questi valori vengano controllati in quest'ordine significa che
l'effetto di alcuni di questi flag possono cambiare se usati in combinazione
con gli altri che vengono prima nella sequenza (è quanto avviene ad esempio
per \const{MS\_BIND} usato con \const{MS\_REMOUNT}). Tratteremo questi
\textit{mount flags} speciali per primi, nell'ordine appena illustrato,
tornando sugli altri più avanti.

Usando il flag \constd{MS\_REMOUNT} si richiede a \func{mount} di rimontare un
filesystem già montato cambiandone le opzioni di montaggio in maniera atomica
(non è cioè necessario smontare e rimontare il filesystem per effettuare il
cambiamento). Questa operazione consente di modificare le opzioni del
filesystem anche se questo è in uso. Gli argomenti \param{source} e
\param{target} devono essere gli stessi usati per il montaggio originale,
mentre sia \param{data} che \param{mountflags} conterranno le nuove opzioni,
\param{filesystemtype} viene ignorato.  Perché l'operazione abbia successo
occorre comunque che il cambiamento sia possibile (ad esempio non sarà
possibile rimontare in sola lettura un filesystem su cui sono aperti file per
la lettura/scrittura).

Qualunque opzione specifica del filesystem indicata con \param{data} può
essere modificata (ma si dovranno rielencare tutte quelle volute), mentre con
\param{mountflags} possono essere modificate solo alcune opzioni generiche:
\const{MS\_LAZYTIME}, \const{MS\_MANDLOCK}, \const{MS\_NOATIME},
\const{MS\_NODEV}, \const{MS\_NODIRATIME}, \const{MS\_NOEXEC},
\const{MS\_NOSUID}, \const{MS\_RELATIME}, \const{MS\_RDONLY},
\const{MS\_STRICTATIME} e \const{MS\_SYNCHRONOUS}. Inoltre dal kernel 3.17 il
comportamento relativo alle opzioni che operano sui tempi di ultimo accesso
dei file (vedi sez.~\ref{sec:file_file_times}) è cambiato e se non si è
indicato nessuno dei vari \texttt{MS\_*ATIME} vengono mantenute le
impostazioni esistenti anziché forzare l'uso di \const{MS\_RELATIME}.

\itindbeg{bind~mount}

Usando il flag \constd{MS\_BIND} si richiede a \func{mount} di effettuare un
cosiddetto \textit{bind mount}, l'operazione che consente di montare una
directory di un filesystem in un'altra directory. L'opzione è disponibile a
partire dai kernel della serie 2.4. In questo caso verranno presi in
considerazione solo gli argomenti \param{source}, che stavolta indicherà la
directory che si vuole montare e non un file di dispositivo, e \param{target}
che indicherà la directory su cui verrà effettuato il \textit{bind mount}. Gli
argomenti \param{filesystemtype} e \param{data} vengono ignorati.

Quello che avviene con questa operazione è che in corrispondenza del
\textit{pathname} indicato da \param{target} viene montato l'\textit{inode} di
\param{source}, così che la porzione di albero dei file presente sotto
\param{source} diventi visibile allo stesso modo sotto
\param{target}. Trattandosi esattamente dei dati dello stesso filesystem, ogni
modifica fatta in uno qualunque dei due rami di albero sarà visibile
nell'altro, visto che entrambi faranno riferimento agli stessi \textit{inode}.

Dal punto di vista del VFS l'operazione è analoga al montaggio di un
filesystem proprio nel fatto che anche in questo caso si inserisce in
corrispondenza della \textit{dentry} di \texttt{target} un diverso
\textit{inode}, che stavolta, invece di essere quello della radice del
filesystem indicato da un file di dispositivo, è quello di una directory già
montata.

Si tenga presente che proprio per questo sotto \param{target} comparirà il
contenuto che è presente sotto \param{source} all'interno del filesystem in
cui quest'ultima è contenuta. Questo potrebbe non corrispondere alla porzione
di albero che sta sotto \param{source} qualora in una sottodirectory di
quest'ultima si fosse effettuato un altro montaggio. In tal caso infatti nella
porzione di albero sotto \param{source} si troverebbe il contenuto del nuovo
filesystem (o di un altro \textit{bind mount}) mentre sotto \param{target} ci
sarebbe il contenuto presente nel filesystem originale.

L'unico altro \textit{mount flag} usabile direttamente con \const{MS\_BIND} è
\const{MS\_REC} che consente di eseguire una operazione di \textit{bind mount}
ricorsiva, in cui sotto \param{target} vengono montati ricorsivamente anche
tutti gli eventuali ulteriori \textit{bind mount} già presenti sotto
\param{source}.

E' però possibile, a partire dal kernel 2.6.26, usare questo flag insieme a
\const{MS\_REMOUNT}, nel qual caso consente di effettuare una modifica delle
opzioni di montaggio del \textit{bind mount} ed in particolare effettuare il
cosiddetto \textit{read-only bind mount} in cui viene onorata anche la
presenza aggiuntiva del flag \const{MS\_RDONLY}. In questo modo si ottiene che
l'accesso ai file sotto \param{target} sia effettuabile esclusivamente in sola
lettura, mantenendo il normale accesso in lettura/scrittura sotto
\param{source}.

Il supporto per il \textit{bind mount} consente di superare i limiti presenti
per gli \textit{hard link} (di cui parleremo in
sez.~\ref{sec:link_symlink_rename}) con la possibilità di fare riferimento
alla porzione dell'albero dei file di un filesystem presente a partire da una
certa directory utilizzando una qualunque altra directory, anche se questa sta
su un filesystem diverso.\footnote{e non c'è neanche il problema di non esser
  più in grado di cancellare un \textit{hard link} ad una directory sullo
  stesso filesystem (vedi sez.~\ref{sec:link_symlink_rename}), per cui su
  Linux questi non sono possibili, dato che in questo caso per la rimozione
  del collegamento basta smontare \param{target}.} Si può così fornire una
alternativa all'uso dei collegamenti simbolici (di cui parleremo in
sez.~\ref{sec:link_symlink_rename}) che funziona correttamente anche
all'intero di un \textit{chroot} (argomento su cui torneremo in
sez.~\ref{sec:file_chroot}).

\itindend{bind~mount}
\itindbeg{shared~subtree}

I quattro flag \const{MS\_PRIVATE}, \const{MS\_SHARED}, \const{MS\_SLAVE} e
\const{MS\_UNBINDABLE} sono stati introdotti a partire dal kernel 2.6.15 per
realizzare l'infrastruttura dei cosiddetti \textit{shared subtree}, che
estendono le funzionalità dei \textit{bind mount}.  La funzionalità nasce
dalle esigenze di poter utilizzare a pieno le funzionalità di isolamento
fornite dal kernel per i processi (i \textit{namespace}, che tratteremo in
sez.~\ref{sec:process_namespaces}) in particolare per quanto riguarda la
possibilità di far avere ad un processo una visione ristretta dei filesystem
montati (il \textit{mount namespace}), ma l'applicazione è comunque rilevante
anche con un classico \textit{chroot} (vedi sez.~\ref{sec:file_chroot}).

\itindbeg{submount}

Abbiamo visto come nella modalità ordinaria in cui si esegue un
\textit{bind mount} sotto \param{target} compaia lo stesso ramo di albero dei
file presente sotto \param{source}, ma limitato a quanto presente nel
filesystem di \param{source}; i risultati di un eventuale
``\textit{submount}'' effettuato all'interno di \param{source} non saranno
visibili. Ed anche se quelli presenti al momento dell'uso di \const{MS\_BIND}
possono essere riottenuti usando \const{MS\_REC}, ogni eventuale
``\textit{submount}'' successivo (che avvenga sotto \param{source} o sotto
\param{target}) resterà ``\textsl{privato}'' al ramo di albero su cui è
avvenuto.

\itindend{submount}
\itindbeg{mount peer group}

Ci sono casi però in cui può risultare utile che eventuali
``\textit{submount}'' siano visibili sui rami di albero presenti al di sotto
di tutte le directory coinvolte in un \textit{bind mount}, anche se effettuati
in un secondo tempo. Per poter ottenere questa funzionalità i
\textit{bind mount} sono stati estesi introducendo i \textit{mount peer
  group}, che consentono di raggrupparli in modo da poter inviare a ciascuno
di essi tutti gli eventi relativi a montaggi o smontaggi effettuati al loro
interno ed avere sempre una propagazione degli stessi che li renda coerenti.

Quando si effettua un montaggio ordinario, o si esegue un \textit{bind mount},
di default non viene utilizzato nessun \textit{mount peer group} ed il
\textit{mount point} viene classificato come ``\textsl{privato}'', nel senso
che abbiamo appena visto.  Si può però marcare un \textit{mount point} come
``\textsl{condiviso}'', ed in questo modo esso verrà associato ad un
\textit{mount peer group} insieme a tutti gli altri ulteriori \textit{mount
  point} per i quali sia stato eseguito un \textit{bind mount}. Questo fa sì
che tutte le volte che si effettua un montaggio o uno smontaggio all'interno
di uno qualunque dei \textit{mount point} del gruppo, questo venga propagato
anche su tutti gli altri e sotto tutti sia visibile sempre lo stesso ramo di
albero dei file.

A completare l'infrastruttura degli \textit{shared subtree} sono state
previste due ulteriori funzionalità: la prima è quella di marcare un
\textit{mount point} come ``\textit{slave}'', in tal caso le operazioni di
montaggio e smontaggio effettuate al suo interno non verranno più propagate
agli altri membri del \textit{mount peer group} di cui fa parte, ma continuerà
a ricevere quelle eseguite negli altri membri.

La seconda funzionalità è quella di marcare un \textit{mount point} come
``\textit{unbindable}''; questo anzitutto impedirà che possa essere usato come
sorgente di un \textit{bind mount} ed inoltre lo renderà privato, con la
conseguenza che quando è presente all'interno di altri \textit{bind mount},
all'interno di questi si vedrà solo il contenuto originale e non quello
risultante da eventuali ulteriori montaggi effettuati al suo interno.

\itindend{mount peer group}

I \textit{mount flag} che controllano le operazioni relative agli
\textit{shared subtree} sono descritti nella lista seguente. Si ricordi che
sono mutuamente esclusivi, e compatibili solo con l'uso degli ulteriori flag
\const{MS\_REC} (che applica ricorsivamente l'operazione a tutti gli eventuali
\textit{mount point} sottostanti) e \const{MS\_SILENT}; in tutti gli altri
casi \func{mount} fallirà con un errore di \errval{EINVAL}. L'unico altro
argomento che deve essere specificato quando li si usano è \param{target};
\param{source}, \param{data} e \param{filesystem} sono ignorati.

\begin{basedescript}{\desclabelwidth{1.9cm}\desclabelstyle{\nextlinelabel}}

\item[\constd{MS\_PRIVATE}] Marca un \textit{mount point} come \textit{private
    mount}.  Di default, finché non lo si marca altrimenti con una delle altre
  opzioni dell'interfaccia, ogni \textit{mount point} è privato. Ogni
  \textit{bind mount} ottenuto da un \textit{mount point} privato si comporta
  come descritto nella trattazione di \const{MS\_BIND}. Si usa questo flag
  principalmente per revocare gli effetti delle altre opzioni e riportare il
  comportamento a quello ordinario.

\item[\constd{MS\_SHARED}] Marca un \textit{mount point} come \textit{shared
    mount}.  Lo scopo dell'opzione è ottenere che tutti i successivi
  \textit{bind mount} ottenuti da un \textit{mount point} così marcato siano
  di tipo \textit{shared} e vengano inseriti nello stesso \textit{mount peer
    group} in modo da ``\textsl{condividere}'' ogni ulteriore operazione di
  montaggio o smontaggio. Con questa opzione le operazioni di montaggio e
  smontaggio effettuate al di sotto di uno \textit{shared mount} vengono
  automaticamente ``\textsl{propagate}'' a tutti gli altri membri del
  \textit{mount peer group} di cui fa parte, in modo che la sezione di albero
  dei file visibile al di sotto di ciascuno di essi sia sempre la stessa.

\item[\constd{MS\_SLAVE}] Marca un \textit{mount point} come \textit{slave
    mount}. Se il \textit{mount point} è parte di un \textit{mount peer group}
  esso diventerà di tipo \textit{slave}: le operazioni di montaggio e
  smontaggio al suo interno non verranno più propagate agli altri membri del
  gruppo, ma continuerà a ricevere quelle eseguite negli altri membri. Se non
  esistono altri membri nel gruppo il \textit{mount point} diventerà privato,
  negli altri casi non subirà nessun cambiamento.

\item[\constd{MS\_UNBINDABLE}] Marca un \textit{mount point} come
  \textit{unbindable mount}.  Un \textit{mount point} marcato in questo modo
  non può essere usato per un \textit{bind mount} del suo contenuto. Si
  comporta cioè come allo stesso modo di un \textit{mount point} ordinario di
  tipo \textit{private} con in più la restrizione che nessuna sua
  sottodirectory (anche se relativa ad un ulteriore montaggio) possa essere
  utilizzata come sorgente di un \textit{bind mount}.
  
\end{basedescript}
\itindend{shared~subtree}

L'ultimo \textit{mount flag} che controlla una modalità operativa di
\func{mount} è \constd{MS\_MOVE}, che consente di effettuare lo spostamento
del \textit{mount point} di un filesystem. La directory del \textit{mount
  point} originale deve essere indicata nell'argomento \param{source}, e la
sua nuova posizione nell'argomento \param{target}. Tutti gli altri argomenti
della funzione vengono ignorati.

Lo spostamento avviene atomicamente, ed il ramo di albero presente sotto
\param{source} sarà immediatamente visibile sotto \param{target}. Non esiste
cioè nessun momento in cui il filesystem non risulti montato in una o
nell'altra directory e pertanto è garantito che la risoluzione di
\textit{pathname} relativi all'interno del filesystem non possa fallire.

Elenchiamo infine i restanti \textit{mount flag}, il cui utilizzo non attiene
alle operazioni di \func{mount}, ma soltanto l'impostazione di opzioni
generiche relative al funzionamento di un filesystem e che vengono per lo più
utilizzati solo in fase di montaggio:

\begin{basedescript}{\desclabelwidth{2.0cm}\desclabelstyle{\nextlinelabel}}
\item[\constd{MS\_DIRSYNC}] Richiede che ogni modifica al contenuto di una
  directory venga immediatamente registrata su disco in maniera sincrona
  (introdotta a partire dai kernel della serie 2.6). L'opzione si applica a
  tutte le directory del filesystem, ma su alcuni filesystem è possibile
  impostarla a livello di singole directory o per i sotto-rami di una directory
  con il comando \cmd{chattr}.\footnote{questo avviene tramite delle opportune
    \texttt{ioctl} (vedi sez.~\ref{sec:file_fcntl_ioctl}).}

  Questo consente di ridurre al minimo il rischio di perdita dei dati delle
  directory in caso di crollo improvviso del sistema, al costo di una certa
  perdita di prestazioni dato che le funzioni di scrittura relative ad
  operazioni sulle directory non saranno più bufferizzate e si bloccheranno
  fino all'arrivo dei dati sul disco prima che un programma possa proseguire.

\item[\constd{MS\_LAZYTIME}] Modifica la modalità di registrazione di tempi
  dei file (vedi sez.~\ref{sec:file_file_times}) per ridurre al massimo gli
  accessi a disco (particolarmente utile per i portatili). Attivandolo i tempi
  dei file vengono mantenuti in memoria e vengono salvati su disco solo in
  quattro casi: quando c'è da eseguire un aggiornamento dei dati
  dell'\textit{inode} per altri motivi, se viene usata una delle funzioni di
  sincronizzazione dei dati su disco (vedi sez.~\ref{sec:file_sync}), se
  l'\textit{inode} viene rimosso dalla memoria, o se è passato un giorno
  dall'ultima registrazione. Introdotto a partire dal kernel 4.0.

  In questo modo si possono ridurre significativamente le scritture su disco
  mantenendo tutte le informazioni riguardo ai tempi dei file, riducendo anche
  l'impatto dell'uso di \const{MS\_STRICTATIME}. Il costo da pagare è il
  rischio, in caso di crash del sistema, di avere dati vecchi fino a 24 ore
  per quel che riguarda i tempi dei file.
  
\item[\constd{MS\_MANDLOCK}] Consente l'uso del \textit{mandatory locking}
  (vedi sez.~\ref{sec:file_mand_locking}) sui file del filesystem. Per poterlo
  utilizzare effettivamente però esso dovrà essere comunque attivato
  esplicitamente per i singoli file impostando i permessi come illustrato in
  sez.~\ref{sec:file_mand_locking}.

\item[\constd{MS\_NOATIME}] Viene disabilitato sul filesystem l'aggiornamento
  dell'\textit{access time} (vedi sez.~\ref{sec:file_file_times}) per
  qualunque tipo di file. Dato che l'aggiornamento dell'\textit{access time} è
  una funzionalità la cui utilità è spesso irrilevante ma comporta un costo
  elevato visto che una qualunque lettura comporta comunque una scrittura su
  disco, questa opzione consente di disabilitarla completamente. La soluzione
  può risultare troppo drastica dato che l'informazione viene comunque
  utilizzata da alcuni programmi, per cui nello sviluppo del kernel sono state
  introdotte altre opzioni che forniscono soluzioni più appropriate e meno
  radicali.

\item[\constd{MS\_NODEV}] Viene disabilitato sul filesystem l'accesso ai file
  di dispositivo eventualmente presenti su di esso. L'opzione viene usata come
  misura di precauzione per rendere inutile la presenza di eventuali file di
  dispositivo su filesystem che non dovrebbero contenerne.\footnote{si ricordi
    che le convenzioni del \textit{Linux Filesystem Hierarchy Standard}
    richiedono che questi siano mantenuti esclusivamente sotto \texttt{/dev}.}

  Viene utilizzata, assieme a \const{MS\_NOEXEC} e \const{MS\_NOSUID}, per
  fornire un accesso più controllato a quei filesystem di cui gli utenti hanno
  il controllo dei contenuti, in particolar modo quelli posti su dispositivi
  rimuovibili. In questo modo si evitano alla radice possibili situazioni in
  cui un utente malizioso inserisce su uno di questi filesystem dei file di
  dispositivo con permessi ``opportunamente'' ampliati che gli consentirebbero
  di accedere anche a risorse cui non dovrebbe.

\item[\constd{MS\_NODIRATIME}] Viene disabilitato sul filesystem
  l'aggiornamento dell'\textit{access time} (vedi
  sez.~\ref{sec:file_file_times}) ma soltanto per le directory. Costituisce
  una alternativa per \const{MS\_NOATIME}, che elimina l'informazione per le
  directory, che in pratica che non viene mai utilizzata, mantenendola per i
  file in cui invece ha un impiego, sia pur limitato.

\item[\constd{MS\_NOEXEC}] Viene disabilitata sul filesystem l'esecuzione di un
  qualunque file eseguibile eventualmente presente su di esso. L'opzione viene
  usata come misura di precauzione per rendere impossibile l'uso di programmi
  posti su filesystem che non dovrebbero contenerne.

  Anche in questo caso viene utilizzata per fornire un accesso più controllato
  a quei filesystem di cui gli utenti hanno il controllo dei contenuti. Da
  questo punto di vista l'opzione è meno importante delle analoghe
  \const{MS\_NODEV} e \const{MS\_NOSUID} in quanto l'esecuzione di un
  programma creato dall'utente pone un livello di rischio nettamente
  inferiore, ed è in genere consentita per i file contenuti nella sua home
  directory.\footnote{cosa che renderebbe superfluo l'attivazione di questa
    opzione, il cui uso ha senso solo per ambienti molto controllati in cui si
    vuole che gli utenti eseguano solo i programmi forniti
    dall'amministratore.}

\item[\constd{MS\_NOSUID}] Viene disabilitato sul filesystem l'effetto dei bit
  dei permessi \acr{suid} e \acr{sgid} (vedi sez.~\ref{sec:file_special_perm})
  eventualmente presenti sui file in esso contenuti. L'opzione viene usata
  come misura di precauzione per rendere inefficace l'effetto di questi bit
  per filesystem in cui non ci dovrebbero essere file dotati di questi
  permessi.

  Di nuovo viene utilizzata, analogamente a \const{MS\_NOEXEC} e
  \const{MS\_NODEV}, per fornire un accesso più controllato a quei filesystem
  di cui gli utenti hanno il controllo dei contenuti. In questo caso si evita
  che un utente malizioso possa inserire su uno di questi filesystem un
  eseguibile con il bit \acr{suid} attivo e di proprietà dell'amministratore o
  di un altro utente, che gli consentirebbe di eseguirlo per conto di
  quest'ultimo.

\item[\constd{MS\_RDONLY}] Esegue il montaggio del filesystem in sola lettura,
  non sarà possibile nessuna modifica ai suoi contenuti. Viene usato tutte le
  volte che si deve accedere ai contenuti di un filesystem con la certezza che
  questo non venga modificato (ad esempio per ispezionare un filesystem
  corrotto). All'avvio di default il kernel monta la radice in questa
  modalità. Si tenga presente che se non viene indicato il filesystem verrà
  montato, o rimontato nel caso lo si usi con \const{MS\_REMOUNT}, in
  lettura/scrittura; questo significa in sostanza che non esiste una opzione
  separata per indicare il montaggio in lettura/scrittura.

\item[\constd{MS\_REC}] Applica ricorsivamente a tutti i \textit{mount point}
  presenti al di sotto del \textit{mount point} indicato gli effetti della
  opzione degli \textit{shared subtree} associata. In questo caso l'argomento
  \param{target} deve fare riferimento ad un \textit{mount point} e tutti gli
  altri argomenti sono ignorati, ed il flag deve essere indicato con uno fra
  \const{MS\_PRIVATE}, \const{MS\_SHARED}, \const{MS\_SLAVE} e
  \const{MS\_UNBINDABLE}. Può anche essere usato con \const{MS\_BIND} per
  richiedere il montaggio ricorsivo anche degli eventuali ulteriori
  \textit{bind mount} presenti sotto \param{target}.

\item[\constd{MS\_RELATIME}] Indica di effettuare l'aggiornamento
  dell'\textit{access time} sul filesystem soltanto quando questo risulti
  antecedente il valore corrente del \textit{modification time} o del
  \textit{change time} (per i tempi dei file si veda
  sez.~\ref{sec:file_file_times}). L'opzione è disponibile a partire dal
  kernel 2.6.20, mentre dal 2.6.30 questo è diventato il comportamento di
  default del sistema, che può essere riportato a quello tradizionale con
  l'uso di \const{MS\_STRICTATIME}. Sempre dal 2.6.30 il comportamento è stato
  anche modificato e l'\textit{access time} viene comunque aggiornato se è più
  vecchio di un giorno.

  L'opzione consente di evitare i problemi di prestazioni relativi
  all'aggiornamento dell'\textit{access time} senza avere impatti negativi
  riguardo le funzionalità, il comportamento adottato infatti consente di
  rendere evidente che vi è stato un accesso dopo la scrittura, ed evitando al
  contempo ulteriori operazioni su disco negli accessi successivi. In questo
  modo l'informazione relativa al fatto che un file sia stato letto resta
  disponibile, ed i programmi che ne fanno uso continuano a funzionare. Con
  l'introduzione di questo comportamento l'uso delle alternative
  \const{MS\_NOATIME} e \const{MS\_NODIRATIME} è sostanzialmente inutile.

\item[\constd{MS\_SILENT}] Richiede la soppressione di alcuni messaggi di
  avvertimento nei log del kernel (vedi sez.~\ref{sec:sess_daemon}). L'opzione
  è presente a partire dal kernel 2.6.17 e sostituisce, utilizzando un nome
  non fuorviante, la precedente \const{MS\_VERBOSE}, introdotta nel kernel
  2.6.12, che aveva lo stesso effetto.

\item[\constd{MS\_STRICTATIME}] Ripristina il comportamento tradizionale per
  cui l'\textit{access time} viene aggiornato ad ogni accesso al
  file. L'opzione è disponibile solo a partire dal kernel 2.6.30 quando il
  comportamento di default del kernel è diventato quello fornito da
  \const{MS\_RELATIME}.

\item[\constd{MS\_SYNCHRONOUS}] Abilita la scrittura sincrona richiedendo che
  ogni modifica al contenuto del filesystem venga immediatamente registrata su
  disco. Lo stesso comportamento può essere ottenuto con il flag
  \const{O\_SYNC} di \func{open} (vedi sez.~\ref{sec:file_open_close}).

  Questa opzione consente di ridurre al minimo il rischio di perdita dei dati
  in caso di crollo improvviso del sistema, al costo di una pesante perdita di
  prestazioni dato che tutte le funzioni di scrittura non saranno più
  bufferizzate e si bloccheranno fino all'arrivo dei dati sul disco. Per un
  compromesso in cui questo comportamento avviene solo per le directory, ed ha
  quindi una incidenza nettamente minore, si può usare \const{MS\_DIRSYNC}.

\end{basedescript}

% NOTE: per l'opzione \texttt{lazytime} introdotta con il kernel 4.0,
% vedi http://lwn.net/Articles/621046/

% NOTE per \const{MS\_SLAVE},\const{MS\_SHARE}, \const{MS\_PRIVATE} e
% \const{MS\_UNBINDABLE} dal 2.6.15 vedi shared subtrees, in particolare
%  * http://lwn.net/Articles/159077/ e
%  * Documentation/filesystems/sharedsubtree.txt

% TODO: (bassa priorità) non documentati ma presenti in sys/mount.h:
%       * MS_POSIXACL
%       * MS_KERNMOUNT
%       * MS_I_VERSION
%       * MS_ACTIVE
%       * MS_NOUSER


Una volta che non si voglia più utilizzare un certo filesystem è possibile
``\textsl{smontarlo}'' usando la funzione di sistema \funcd{umount}, il cui
prototipo è:

\begin{funcproto}{
\fhead{sys/mount.h}
\fdecl{umount(const char *target)}
\fdesc{Smonta un filesystem.} 
}
{La funzione ritorna  $0$ in caso di successo e $-1$ per un errore,
  nel qual caso \var{errno} assumerà uno dei valori: 
  \begin{errlist}
  \item[\errcode{EBUSY}] il filesystem è occupato.
  \item[\errcode{EINVAL}] \param{target} non è un \textit{mount point}.
  \item[\errcode{EPERM}] il processo non ha i privilegi di
    amministratore.\footnotemark 
  \end{errlist}
  ed inoltre \errval{EFAULT}, \errval{ELOOP}, \errval{ENAMETOOLONG},
  \errval{ENOENT}, \errval{ENOMEM} nel loro significato generico.  }
\end{funcproto}

\footnotetext{più precisamente la capacità \const{CAP\_SYS\_ADMIN}, vedi
  sez.~\ref{sec:proc_capabilities}.}

La funzione prende il nome della directory su cui il filesystem è montato e
non il file o il dispositivo che è stato montato,\footnote{questo è vero a
  partire dal kernel 2.3.99-pre7, prima esistevano due chiamate separate e la
  funzione poteva essere usata anche specificando il file di dispositivo.} in
quanto a partire dai kernel della serie 2.4.x è possibile montare lo stesso
dispositivo in più punti. Nel caso più di un filesystem sia stato montato
sullo stesso \textit{mount point} viene smontato quello che è stato montato
per ultimo. Si tenga presente che la funzione fallisce se il filesystem è
``\textsl{occupato}'', cioè quando ci sono ancora dei file aperti sul
filesystem, se questo contiene la directory di lavoro (vedi
sez.~\ref{sec:file_work_dir}) di un qualunque processo o il \textit{mount
  point} di un altro filesystem.

Linux provvede inoltre una seconda funzione di sistema, \funcd{umount2}, che
consente un maggior controllo delle operazioni, come forzare lo smontaggio di
un filesystem anche quando questo risulti occupato; il suo prototipo è:

\begin{funcproto}{
\fhead{sys/mount.h}
\fdecl{umount2(const char *target, int flags)}
\fdesc{Smonta un filesystem.} 
}
{La funzione ritorna  $0$ in caso di successo e $-1$ per un errore,
  nel qual caso \var{errno} assumerà uno dei valori: 
  \begin{errlist}
     \item[\errcode{EAGAIN}] si è chiamata la funzione con \const{MNT\_EXPIRE}
       ed il filesystem non era occupato.
     \item[\errcode{EBUSY}] \param{target} è la directory di lavoro di qualche
       processo, o contiene dei file aperti, o un altro \textit{mount point}.
     \item[\errcode{EINVAL}] \param{target} non è un \textit{mount point} o si
       è usato \const{MNT\_EXPIRE} con \const{MNT\_FORCE} o
       \const{MNT\_DETACH} o si è specificato un flag non esistente.
  \end{errlist}
  e tutti gli altri valori visti per \func{umount} con lo stesso significato.}
\end{funcproto}

Il valore di \param{flags} è una maschera binaria dei flag che controllano le
modalità di smontaggio, che deve essere specificato con un OR aritmetico delle
costanti illustrate in tab.~\ref{tab:umount2_flags}.  Specificando
\constd{MNT\_FORCE} la funzione cercherà di liberare il filesystem anche se è
occupato per via di una delle condizioni descritte in precedenza. A seconda
del tipo di filesystem alcune (o tutte) possono essere superate, evitando
l'errore di \errcode{EBUSY}. In tutti i casi prima dello smontaggio viene
eseguita una sincronizzazione dei dati.

\begin{table}[!htb]
  \centering
  \footnotesize
  \begin{tabular}[c]{|l|p{8cm}|}
    \hline
    \textbf{Costante} & \textbf{Descrizione}\\
    \hline
    \hline
    \const{MNT\_FORCE}  & Forza lo smontaggio del filesystem anche se questo è
                           occupato (presente dai kernel della serie 2.2).\\
    \const{MNT\_DETACH} & Esegue uno smontaggio ``\textsl{pigro}'', in cui si
                           blocca l'accesso ma si aspetta che il filesystem si
                           liberi (presente dal kernel 2.4.11 e dalla
                           \acr{glibc} 2.11).\\ 
    \const{MNT\_EXPIRE} & Se non occupato marca un \textit{mount point} come
                           ``\textsl{in scadenza}'' in modo che ad una
                           successiva chiamata senza utilizzo del filesystem
                           questo venga smontato (presente dal 
                           kernel 2.6.8 e dalla \acr{glibc} 2.11).\\ 
    \const{UMOUNT\_NOFOLLOW}& Non dereferenzia \param{target} se questo è un
                               collegamento simbolico (vedi
                               sez.~\ref{sec:link_symlink_rename}) evitando
                               problemi di sicurezza (presente dal kernel
                               2.6.34).\\  
    \hline
  \end{tabular}
  \caption{Costanti che identificano i bit dell'argomento \param{flags}
    della funzione \func{umount2}.} 
  \label{tab:umount2_flags}
\end{table}

Con l'opzione \constd{MNT\_DETACH} si richiede invece uno smontaggio
``\textsl{pigro}'' (o \textit{lazy umount}) in cui il filesystem diventa
inaccessibile per i nuovi processi subito dopo la chiamata della funzione, ma
resta accessibile per quelli che lo hanno ancora in uso e non viene smontato
fintanto che resta occupato.

Con \constd{MNT\_EXPIRE}, che non può essere specificato insieme agli altri
due, si marca il \textit{mount point} di un filesystem non occupato come
``\textsl{in scadenza}'', in tal caso \func{umount2} ritorna con un errore di
\errcode{EAGAIN}, mentre in caso di filesystem occupato si sarebbe ricevuto
\errcode{EBUSY}.  Una volta marcato, se nel frattempo non viene fatto nessun
uso del filesystem, ad una successiva chiamata con \const{MNT\_EXPIRE} questo
verrà smontato. Questo flag consente di realizzare un meccanismo che smonti
automaticamente i filesystem che restano inutilizzati per un certo periodo di
tempo.

Infine il flag \constd{UMOUNT\_NOFOLLOW} non dereferenzia \param{target} se
questo è un collegamento simbolico (vedi
sez.~\ref{sec:link_symlink_rename}). Questa è una misura di sicurezza
introdotta per evitare, per quei filesystem per il quale è prevista una
gestione diretta da parte degli utenti, come quelli basati su \itindex{FUSE}
FUSE,\footnote{il \textit{Filesystem in USEr space} (FUSE) è una delle più
  interessanti applicazioni del VFS che consente, tramite un opportuno modulo,
  di implementarne le funzioni in \textit{user space}, così da rendere
  possibile l'implementazione di un qualunque filesystem (con applicazioni di
  grande interesse come il filesystem cifrato \textit{encfs} o il filesystem
  di rete \textit{sshfs}) che possa essere usato direttamente per conto degli
  utenti.}  che si possano passare ai programmi che effettuano lo smontaggio
dei filesystem, che in genere sono privilegiati ma consentono di agire solo
sui propri \textit{mount point}, dei collegamenti simbolici che puntano ad
altri \textit{mount point}, ottenendo così la possibilità di smontare
qualunque filesystem.


Altre due funzioni di sistema specifiche di Linux,\footnote{esse si trovano
  anche su BSD, ma con una struttura diversa.} utili per ottenere in maniera
diretta informazioni riguardo al filesystem su cui si trova un certo file,
sono \funcd{statfs} e \funcd{fstatfs}, i cui prototipi sono:

\begin{funcproto}{
\fhead{sys/vfs.h}
\fdecl{int statfs(const char *path, struct statfs *buf)}
\fdecl{int fstatfs(int fd, struct statfs *buf)}
\fdesc{Restituiscono informazioni relative ad un filesystem.} 
}
{Le funzioni ritornano $0$ in caso di successo e $-1$ per un errore,
  nel qual caso \var{errno} assumerà uno dei valori: 
  \begin{errlist}
  \item[\errcode{ENOSYS}] il filesystem su cui si trova il file specificato
    non supporta la funzione.
  \end{errlist} ed inoltre \errval{EFAULT} ed \errval{EIO} per entrambe,
  \errval{EBADF} per \func{fstatfs}, \errval{ENOTDIR}, \errval{ENAMETOOLONG},
  \errval{ENOENT}, \errval{EACCES}, \errval{ELOOP} per \func{statfs} nel loro
  significato generico.}
\end{funcproto}

Queste funzioni permettono di ottenere una serie di informazioni generali
riguardo al filesystem su cui si trova il file specificato con un
\textit{pathname} per \func{statfs} e con un file descriptor (vedi
sez.~\ref{sec:file_fd}) per \func{statfs}.  Le informazioni vengono restituite
all'indirizzo \param{buf} di una struttura \struct{statfs} definita come in
fig.~\ref{fig:sys_statfs}, ed i campi che sono indefiniti per il filesystem in
esame sono impostati a zero.  I valori del campo \var{f\_type} sono definiti
per i vari filesystem nei relativi file di header dei sorgenti del kernel da
costanti del tipo \var{XXX\_SUPER\_MAGIC}, dove \var{XXX} in genere è il nome
del filesystem stesso.

\begin{figure}[!htb]
  \footnotesize \centering
  \begin{minipage}[c]{0.8\textwidth}
    \includestruct{listati/statfs.h}
  \end{minipage}
  \normalsize 
  \caption{La struttura \structd{statfs}.} 
  \label{fig:sys_statfs}
\end{figure}

\conffilebeg{/etc/mtab} 

La \acr{glibc} provvede infine una serie di funzioni per la gestione dei due
file \conffiled{/etc/fstab}\footnote{più precisamente \funcm{setfsent},
  \funcm{getfsent}, \funcm{getfsfile}, \funcm{getfsspec}, \funcm{endfsent}.}
ed \conffile{/etc/mtab}\footnote{più precisamente \funcm{setmntent},
  \funcm{getmntent},\funcm{getmntent\_r}, \funcm{addmntent},\funcm{endmntent},
  \funcm{hasmntopt}.} che convenzionalmente sono usati in quasi tutti i
sistemi unix-like per mantenere rispettivamente le informazioni riguardo ai
filesystem da montare e a quelli correntemente montati. Le funzioni servono a
leggere il contenuto di questi file in opportune strutture \structd{fstab} e
\structd{mntent}, e, nel caso di \conffile{/etc/mtab}, per inserire e
rimuovere le voci presenti nel file.

In generale si dovrebbero usare queste funzioni, in particolare quelle
relative a \conffile{/etc/mtab}, quando si debba scrivere un programma che
effettua il montaggio di un filesystem. In realtà in questi casi è molto più
semplice invocare direttamente il programma \cmd{mount}. Inoltre l'uso stesso
di \conffile{/etc/mtab} è considerato una pratica obsoleta, in quanto se non
aggiornato correttamente (cosa che è impossibile se la radice è montata in
sola lettura) il suo contenuto diventa fuorviante.

Per questo motivo il suo utilizzo viene deprecato ed in molti casi viene già
oggi sostituito da un collegamento simbolico a \procfile{/proc/mounts}, che
contiene una versione degli stessi contenuti (vale a dire l'elenco dei
filesystem montati) generata direttamente dal kernel, e quindi sempre
disponibile e sempre aggiornata. Per questo motivo tralasceremo la
trattazione, di queste funzioni, rimandando al manuale della \acr{glibc}
\cite{GlibcMan} per la documentazione completa.

\conffileend{/etc/mtab}

% TODO (bassa priorità) scrivere delle funzioni (getfsent e getmntent &C)
% TODO (bassa priorità) documentare ? swapon e swapoff (man 2 ...) 

% TODO con il 5.2 è stata introdotta una serie di nuove syscall per montare un
% filesystem, vedi https://lwn.net/Articles/759499/ e
% https://git.kernel.org/linus/f1b5618e013a 


\section{La gestione di file e directory}
\label{sec:file_dir}

In questa sezione esamineremo le funzioni usate per la manipolazione dei nomi
file e directory, per la creazione di collegamenti simbolici e diretti, per la
gestione e la lettura delle directory.  In particolare ci soffermeremo sulle
conseguenze che derivano dalla architettura di un filesystem unix-like
illustrata in sez.~\ref{sec:file_filesystem} per quanto attiene il
comportamento e gli effetti delle varie funzioni. Tratteremo infine la
directory di lavoro e le funzioni per la gestione di file speciali e
temporanei.


\subsection{La gestione dei nomi dei file}
\label{sec:link_symlink_rename}

% \subsection{Le funzioni \func{link} e \func{unlink}}
% \label{sec:file_link}

Una caratteristica comune a diversi sistemi operativi è quella di poter creare
dei nomi alternativi, come gli alias del vecchio MacOS o i collegamenti di
Windows o i nomi logici del VMS, che permettono di fare riferimento allo
stesso file chiamandolo con nomi diversi o accedendovi da directory diverse.
Questo è possibile anche in ambiente Unix, dove un nome alternativo viene
usualmente chiamato ``\textsl{collegamento}'' (o \textit{link}).  Data
l'architettura del sistema riguardo la gestione dei file vedremo però che ci
sono due metodi sostanzialmente diversi per fare questa operazione.

\itindbeg{hard~link}
\index{collegamento!diretto|(}

In sez.~\ref{sec:file_filesystem} abbiamo spiegato come la capacità di
chiamare un file con nomi diversi sia connaturata con l'architettura di un
filesystem per un sistema Unix, in quanto il nome del file che si trova in una
directory è solo un'etichetta associata ad un puntatore che permette di
ottenere il riferimento ad un \textit{inode}, e che è quest'ultimo che viene
usato dal kernel per identificare univocamente gli oggetti sul filesystem.

Questo significa che fintanto che si resta sullo stesso filesystem la
realizzazione di un \textit{link} è immediata: uno stesso file può avere tanti
nomi diversi, dati da altrettante associazioni diverse allo stesso
\textit{inode} effettuate tramite ``etichette'' diverse in directory
diverse. Si noti anche come nessuno di questi nomi possa assumere una
particolare preferenza o originalità rispetto agli altri, in quanto tutti
fanno comunque riferimento allo stesso \textit{inode} e quindi tutti
otterranno lo stesso file.

Quando si vuole aggiungere ad una directory una voce che faccia riferimento ad
un file già esistente come appena descritto, per ottenere quello che viene
denominato ``\textsl{collegamento diretto}'' (o \textit{hard link}), si deve
usare la funzione di sistema \funcd{link}, il cui prototipo è:

\begin{funcproto}{
\fhead{unistd.h}
\fdecl{int link(const char *oldpath, const char *newpath)}
\fdesc{Crea un nuovo collegamento diretto (\textit{hard link}).} 
}
{La funzione ritorna $0$ in caso di successo e $-1$ per un errore,
  nel qual caso \var{errno} assumerà uno dei valori: 
  \begin{errlist}
  \item[\errcode{EEXIST}] un file (o una directory) di nome \param{newpath}
    esiste già.
  \item[\errcode{EMLINK}] ci sono troppi collegamenti al file \param{oldpath}
    (il numero massimo è specificato dalla variabile \const{LINK\_MAX}, vedi
    sez.~\ref{sec:sys_limits}).
  \item[\errcode{EPERM}] il filesystem che contiene \param{oldpath} e
    \param{newpath} non supporta i collegamenti diretti, è una directory o per
    \param{oldpath} non si rispettano i criteri per i \textit{protected
      hardlink}.\footnotemark 
  \item[\errcode{EXDEV}] i file \param{oldpath} e \param{newpath} non fanno
    riferimento ad un filesystem montato sullo stesso 
    \textit{mount point}.
  \end{errlist} ed inoltre \errval{EACCES}, \errval{EDQUOT}, \errval{EFAULT},
  \errval{EIO}, \errval{ELOOP}, \errval{ENAMETOOLONG}, \errval{ENOENT},
  \errval{ENOMEM}, \errval{ENOSPC}, \errval{ENOTDIR}, \errval{EROFS} nel loro
  significato generico.}
\end{funcproto}

\footnotetext{i \textit{protected hardlink} sono una funzionalità di
  protezione introdotta con il kernel 3.16 (si veda
  sez.~\ref{sec:procadv_security_misc} per i dettagli) che limita la capacità
  di creare un \textit{hard link} ad un file qualunque.}

La funzione crea in \param{newpath} un collegamento diretto al file indicato
da \param{oldpath}. Per quanto detto la creazione di un nuovo collegamento
diretto non copia il contenuto del file, ma si limita a creare la voce
specificata da \param{newpath} nella directory corrispondente e l'unica
proprietà del file che verrà modificata sarà il numero di riferimenti al file
(il campo \var{i\_nlink} della struttura \kstruct{inode}, vedi
fig.~\ref{fig:kstruct_inode}) che verrà aumentato di uno. In questo modo lo
stesso file potrà essere acceduto sia con \param{newpath} che
con \param{oldpath}.

Per quanto dicevamo in sez.~\ref{sec:file_filesystem} la creazione di un
collegamento diretto è possibile solo se entrambi i \textit{pathname} sono
nello stesso filesystem ed inoltre esso deve supportare gli \textit{hard link}
(il meccanismo non è disponibile ad esempio con il filesystem \acr{vfat} di
Windows). In realtà la funzione ha un ulteriore requisito, e cioè che non solo
che i due file siano sullo stesso filesystem, ma anche che si faccia
riferimento ad essi all'interno dello stesso \textit{mount point}.\footnote{si
  tenga presente infatti, come detto in sez.~\ref{sec:filesystem_mounting},
  che a partire dal kernel 2.4 uno stesso filesystem può essere montato più
  volte su directory diverse.}
La funzione inoltre opera sia sui file ordinari che sugli altri oggetti del
filesystem, con l'eccezione delle directory. In alcune versioni di Unix solo
l'amministratore è in grado di creare un collegamento diretto ad un'altra
directory: questo viene fatto perché con una tale operazione è possibile
creare dei \textit{loop} nel filesystem (vedi fig.~\ref{fig:file_link_loop})
la cui rimozione diventerebbe piuttosto complicata.\footnote{occorrerebbe
  infatti eseguire il programma \cmd{fsck} per riparare il filesystem, perché
  in caso di \textit{loop} la directory non potrebbe essere più svuotata,
  contenendo comunque se stessa, e quindi non potrebbe essere rimossa.}

Data la pericolosità di questa operazione, e visto che i collegamenti
simbolici (che tratteremo a breve) ed i \textit{bind mount} (già visti in
sez.~\ref{sec:filesystem_mounting}) possono fornire la stessa funzionalità
senza questi problemi, nel caso di Linux questa capacità è stata completamente
disabilitata, e al tentativo di creare un collegamento diretto ad una
directory la funzione \func{link} restituisce sempre l'errore \errcode{EPERM}.

Un ulteriore comportamento peculiare di Linux è quello in cui si crea un
\textit{hard link} ad un collegamento simbolico. In questo caso lo standard
POSIX.1-2001 prevederebbe che quest'ultimo venga risolto e che il collegamento
sia effettuato rispetto al file cui esso punta, e che venga riportato un
errore qualora questo non esista o non sia un file. Questo era anche il
comportamento iniziale di Linux ma a partire dai kernel della serie
2.0.x\footnote{per la precisione il comportamento era quello previsto dallo
  standard POSIX fino al kernel di sviluppo 1.3.56, ed è stato temporaneamente
  ripristinato anche durante lo sviluppo della serie 2.1.x, per poi tornare al
  comportamento attuale fino ad oggi (per riferimento si veda
  \url{http://lwn.net/Articles/293902}).} è stato adottato un comportamento
che non segue più lo standard per cui l'\textit{hard link} viene creato nei
confronti del collegamento simbolico, e non del file cui questo punta. La
revisione POSIX.1-2008 lascia invece il comportamento dipendente
dall'implementazione, cosa che rende Linux conforme a questa versione
successiva dello standard.

\itindbeg{symbolic~link}
\index{collegamento!simbolico|(}

La ragione di questa differenza rispetto al vecchio standard, presente anche
in altri sistemi unix-like, è dovuta al fatto che un collegamento simbolico
può fare riferimento anche ad un file non esistente o a una directory, per i
quali l'\textit{hard link} non può essere creato, per cui la scelta di seguire
il collegamento simbolico è stata ritenuta una scelta scorretta nella
progettazione dell'interfaccia.  Infatti se non ci fosse il comportamento
adottato da Linux sarebbe impossibile creare un \textit{hard link} ad un
collegamento simbolico, perché la funzione lo risolverebbe e l'\textit{hard
  link} verrebbe creato verso la destinazione. Invece evitando di seguire lo
standard l'operazione diventa possibile, ed anche il comportamento della
funzione risulta molto più comprensibile. Tanto più che se proprio se si vuole
creare un \textit{hard link} rispetto alla destinazione di un collegamento
simbolico è sempre possibile farlo direttamente.\footnote{ciò non toglie che
  questo comportamento possa causare problemi, come nell'esempio descritto
  nell'articolo citato nella nota precedente, a programmi che non si aspettano
  questa differenza rispetto allo standard POSIX.}

Dato che \func{link} crea semplicemente dei nomi che fanno riferimenti agli
\textit{inode}, essa può funzionare soltanto per file che risiedono sullo
stesso filesystem e solo per un filesystem di tipo Unix.  Inoltre abbiamo
visto che in Linux non è consentito eseguire un collegamento diretto ad una
directory.

Per ovviare a queste limitazioni, come accennato all'inizio, i sistemi
unix-like supportano un'altra forma di collegamento, detta
``\textsl{collegamento simbolico}'' (o anche \textit{soft link} o
\textit{symbolic link}). In questo caso si tratta, come avviene in altri
sistemi operativi, di file speciali che contengono semplicemente il
riferimento ad un altro file (o directory). In questo modo è possibile
effettuare \textit{link} anche attraverso filesystem diversi, a file posti in
filesystem che non supportano i collegamenti diretti, a delle directory, ed
anche a file che non esistono ancora.

\itindend{hard~link}
\index{collegamento!diretto|)}

Il meccanismo funziona in quanto i \textit{symbolic link} sono riconosciuti
come tali dal kernel\footnote{è uno dei diversi tipi di file visti in
  tab.~\ref{tab:file_file_types}, contrassegnato come tale nell'\textit{inode}
  e riconoscibile dal valore del campo \var{st\_mode} della struttura
  \struct{stat} (vedi sez.~\ref{sec:file_stat}).} e tutta una serie di
funzioni di sistema (come \func{open} o \func{stat}) quando ricevono come
argomento il \textit{pathname} di un collegamento simbolico vanno
automaticamente ad operare sul file da esso specificato. La funzione di
sistema che permette di creare un nuovo collegamento simbolico è
\funcd{symlink}, ed il suo prototipo è:

\begin{funcproto}{
\fhead{unistd.h}
\fdecl{int symlink(const char *oldpath, const char *newpath)}
\fdesc{Crea un nuovo collegamento simbolico (\textit{symbolic link}).} 
}
{La funzione ritorna $0$ in caso di successo e $-1$ per un errore,
  nel qual caso \var{errno} assumerà uno dei valori: 
  \begin{errlist}
  \item[\errcode{EACCES}]  o non si hanno i permessi sulla directory in cui
    creare il \textit{link}.
  \item[\errcode{EEXIST}] esiste già un file \param{newpath}.
  \item[\errcode{ENOENT}] una componente di \param{newpath} non esiste o
    \param{oldpath} è una stringa vuota.
  \item[\errcode{EPERM}] il filesystem che contiene \param{newpath} non
    supporta i collegamenti simbolici.
  \item[\errcode{EROFS}] \param{newpath} è su un filesystem montato in sola
    lettura.
  \end{errlist} ed inoltre \errval{EDQUOT}, \errval{EFAULT}, \errval{EIO},
  \errval{ELOOP}, \errval{ENAMETOOLONG}, \errval{ENOMEM}, \errval{ENOSPC} e
  \errval{ENOTDIR} nel loro significato generico.}
\end{funcproto}


La funzione crea un nuovo collegamento simbolico \param{newpath} che fa
riferimento ad \param{oldpath}.  Si tenga presente che la funzione non
effettua nessun controllo sull'esistenza di un file di nome \param{oldpath},
ma si limita ad inserire il \textit{pathname} nel collegamento
simbolico. Pertanto un collegamento simbolico può anche riferirsi ad un file
che non esiste ed in questo caso si ha quello che viene chiamato un
\itindex{dangling~link} \textit{dangling link}, letteralmente un
\index{collegamento!ciondolante} ``\textsl{collegamento ciondolante}''. Ad
esempio possiamo usare il comando \cmd{ln} per creare un collegamento
simbolico nella nostra directory con:
\begin{Console}
piccardi@hain:~/gapil$ \textbf{ln -s /tmp/tmp_file symlink}
\end{Console}
%$
e questo avrà successo anche se \file{/tmp/tmp\_file} non esiste:
\begin{Console}
piccardi@hain:~/gapil$ \textbf{ls symlink}
symlink
\end{Console}
%$
ma questo può generare confusione, perché accedendo in lettura a
\file{symlink}, ad esempio con \cmd{cat}, otterremmo:
\begin{Console}
piccardi@hain:~/gapil$ \textbf{cat symlink}
cat: symlink: No such file or directory
\end{Console}
%$
con un errore che può sembrare sbagliato, dato che \cmd{ls} ci ha mostrato in
precedenza l'esistenza di \file{symlink}. Se invece andassimo a scrivere su
\file{symlink}, l'effetto sarebbe quello di ottenere la creazione di
\file{/tmp/tmp\_file} (che a quel punto verrebbe creato) senza errori.

Come accennato i collegamenti simbolici sono risolti automaticamente dal
kernel all'invocazione delle varie \textit{system call}. In
tab.~\ref{tab:file_symb_effect} si è riportato un elenco dei comportamenti
delle varie funzioni di sistema che operano sui file nei confronti della
risoluzione dei collegamenti simbolici, specificando quali li seguono e quali
invece possono operare direttamente sui loro contenuti.
\begin{table}[htb]
  \centering
  \footnotesize
  \begin{tabular}[c]{|l|c|c|}
    \hline
    \textbf{Funzione} & \textbf{Segue il link} & \textbf{Non segue il link} \\
    \hline 
    \hline 
    \func{access}   & $\bullet$ & --        \\
    \func{chdir}    & $\bullet$ & --        \\
    \func{chmod}    & $\bullet$ & --        \\
    \func{chown}    & --        & $\bullet$ \\
    \func{creat}    & $\bullet$ & --        \\
    \func{exec}     & $\bullet$ & --        \\
    \func{lchown}   & $\bullet$ & --        \\
    \func{link}\footnotemark & --        & $\bullet$ \\
    \func{lstat}    & --        & $\bullet$ \\
    \func{mkdir}    & $\bullet$ & --        \\
    \func{mkfifo}   & $\bullet$ & --        \\
    \func{mknod}    & $\bullet$ & --        \\
    \func{open}     & $\bullet$ & --        \\
    \func{opendir}  & $\bullet$ & --        \\
    \func{pathconf} & $\bullet$ & --        \\
    \func{readlink} & --        & $\bullet$ \\
    \func{remove}   & --        & $\bullet$ \\
    \func{rename}   & --        & $\bullet$ \\
    \func{stat}     & $\bullet$ & --        \\
    \func{truncate} & $\bullet$ & --        \\
    \func{unlink}   & --        & $\bullet$ \\
    \hline 
  \end{tabular}
  \caption{Uso dei collegamenti simbolici da parte di alcune funzioni.}
  \label{tab:file_symb_effect}
\end{table}

\footnotetext{a partire dalla serie 2.0, e contrariamente a quanto indicato
  dallo standard POSIX.1-2001.}

Si noti che non si è specificato il comportamento delle funzioni che operano
con i file descriptor (che tratteremo nel prossimo capitolo), in quanto la
risoluzione del collegamento simbolico viene in genere effettuata dalla
funzione che restituisce il file descriptor (normalmente la \func{open}, vedi
sez.~\ref{sec:file_open_close}) e tutte le operazioni seguenti fanno
riferimento solo a quest'ultimo.

Si tenga anche presente che a partire dal kernel 3.16, se si abilita la
funzionalità dei \textit{protected symlinks} (attiva di default in tutte le
distribuzioni più recenti) la risoluzione dei nomi attraverso un collegamento
simbolico può fallire per una serie di restrizione di sicurezza aggiuntive
imposte dal meccanismo (si consulti sez.~\ref{sec:procadv_security_misc} per i
dettagli).

Dato che, come indicato in tab.~\ref{tab:file_symb_effect}, funzioni come la
\func{open} seguono i collegamenti simbolici, occorrono funzioni apposite per
accedere alle informazioni del collegamento invece che a quelle del file a cui
esso fa riferimento. Quando si vuole leggere il contenuto di un collegamento
simbolico si usa la funzione di sistema \funcd{readlink}, il cui prototipo è:

\begin{funcproto}{
\fhead{unistd.h}
\fdecl{int readlink(const char *pathname, char *buff, size\_t size)}
\fdesc{Legge il contenuto di un collegamento simbolico.} 
}
{La funzione ritorna il numero di caratteri letti dentro \param{buff} in caso
  di successo e $-1$ per un errore,  nel qual caso \var{errno} assumerà uno
  dei valori:
  \begin{errlist}
  \item[\errcode{EACCES}] non si hanno i permessi di attraversamento di una
    delle directory del pathname
  \item[\errcode{EINVAL}] \param{pathname} non è un collegamento simbolico
    o \param{size} non è positiva.
  \end{errlist} ed inoltre \errval{EFAULT}, \errval{EIO}, \errval{ELOOP},
  \errval{ENAMETOOLONG}, \errval{ENOENT}, \errval{ENOMEM} e \errval{ENOTDIR}
  nel loro significato generico.}
\end{funcproto}

La funzione legge il \textit{pathname} a cui fa riferimento il collegamento
simbolico indicato dall'argomento \param{pathname} scrivendolo sul buffer
\param{buff} di dimensione \param{size}. Si tenga presente che la funzione non
termina la stringa con un carattere nullo e che se questa è troppo lunga la
tronca alla dimensione specificata da \param{size} per evitare di scrivere
dati oltre le dimensioni del buffer.

\begin{figure}[htb]
  \centering
  \includegraphics[width=8cm]{img/link_loop}
  \caption{Esempio di loop nel filesystem creato con un collegamento
    simbolico.}
  \label{fig:file_link_loop}
\end{figure}

Come accennato uno dei motivi per cui non sono consentiti \textit{hard link}
alle directory è che questi possono creare dei \textit{loop} nella risoluzione
dei nomi che non possono essere eliminati facilmente. Invece è sempre
possibile, ed in genere anche molto utile, creare un collegamento simbolico ad
una directory, anche se in questo caso si potranno ottenere anche dei
\textit{loop}.

La situazione è illustrata in fig.~\ref{fig:file_link_loop}, che riporta la
struttura della directory \file{/boot}. Come si vede si è creato al suo
interno un collegamento simbolico che punta di nuovo a
\file{/boot}.\footnote{il \textit{loop} mostrato in
  fig.~\ref{fig:file_link_loop} è stato usato per poter permettere a al
  \textit{bootloader} \cmd{grub} di vedere i file contenuti nella directory
  \file{/boot} con lo stesso \textit{pathname} con cui verrebbero visti dal
  sistema operativo, anche quando si trovano, come accade spesso, su una
  partizione separata (che \cmd{grub} all'avvio vedrebbe come \file{/}).} Un
\textit{loop} di di questo tipo però può causare problemi per tutti i
programmi che effettuano la scansione di una directory, e ad esempio se
lanciassimo un comando come \code{grep -r linux *}, il \textit{loop} nella
directory porterebbe ad esaminare \file{/boot}, \file{/boot/boot},
\file{/boot/boot/boot} e così via.

Per questo motivo il kernel e le librerie prevedono che nella risoluzione di
un \textit{pathname} possano essere seguiti fino ad un certo numero massimo di
collegamenti simbolici, il cui valore limite è specificato dalla costante
\constd{MAXSYMLINKS}. Se il limite viene superato si ha un errore ed
\var{errno} viene impostata al valore \errcode{ELOOP}, che nella quasi
totalità dei casi indica appunto che si è creato un collegamento simbolico che
fa riferimento ad una directory del suo stesso \textit{pathname}.


\itindend{symbolic~link}
\index{collegamento!simbolico|)}

Un'altra funzione relativa alla gestione dei nomi dei file, anche se a prima
vista parrebbe riguardare un argomento completamente diverso, è quella per la
cancellazione di un file. In realtà una \textit{system call} che serva proprio
a cancellare un file non esiste neanche perché, come accennato in
sez.~\ref{sec:file_filesystem}, quando in un sistema unix-like si richiede la
rimozione di un file, quello che si va a cancellare è soltanto la voce che
referenzia il suo \textit{inode} all'interno di una directory.

La funzione di sistema che consente di effettuare questa operazione, il cui
nome come si può notare ha poco a che fare con il concetto di rimozione, è
\funcd{unlink}, ed il suo prototipo è:

\begin{funcproto}{
\fhead{unistd.h}
\fdecl{int unlink(const char *pathname)}
\fdesc{Cancella un file.} 
}
{La funzione ritorna $0$ in caso di successo e $-1$ per un errore,
  nel qual caso \var{errno} assumerà uno dei valori:\footnotemark  
  \begin{errlist}
  \item[\errcode{EACCES}] non si ha il permesso di scrittura sulla directory
    che contiene \param{pathname} o quello di attraversamento per una delle
    directory superiori.
  \item[\errcode{EBUSY}] \param{pathname} non può essere rimosso perché è in
    uso da parte del sistema (in particolare per i cosiddetti \textit{silly
      renames} di NFS).
  \item[\errcode{EISDIR}] \param{pathname} si riferisce ad una
    directory.
  \item[\errcode{EPERM}] il filesystem non consente l'operazione, o la
    directory che contiene \param{pathname} ha lo \textit{sticky bit} e non si
    è il proprietario del file o non si hanno privilegi amministrativi. 
  \end{errlist} ed inoltre \errval{EFAULT}, \errval{EIO}, \errval{ELOOP},
  \errval{ENOENT}, \errval{ENOMEM}, \errval{ENOTDIR}, \errval{EROFS} nel loro
  significato generico.}
\end{funcproto}

\footnotetext{questa funzione su Linux ha alcune peculiarità nei codici di
  errore, in particolare riguardo la rimozione delle directory che non è mai
  permessa e che causa l'errore \errcode{EISDIR}; questo è un valore specifico
  di Linux non conforme allo standard POSIX che prescrive invece l'uso di
  \errcode{EPERM} in caso l'operazione non sia consentita o il processo non
  abbia privilegi sufficienti, valore che invece Linux usa anche se il
  filesystem non supporta la funzione, inoltre il codice \errcode{EBUSY} nel
  caso la directory sia occupata su Linux non esiste.}

La funzione elimina il nome specificato dall'argomento \param{pathname} nella
directory che lo contiene e decrementa il numero di riferimenti nel relativo
\textit{inode}; come per \func{link} queste due operazioni sono effettuate
all'interno della \textit{system call} in maniera atomica rispetto ai
processi.

Si ricordi che, anche se se ne è rimosso il nome, un file viene realmente
cancellato soltanto quando il numero di collegamenti mantenuto
nell'\textit{inode} diventa nullo; solo allora l'\textit{inode} viene
disallocato e lo spazio che il file occupava sul disco viene liberato.

Si tenga presente comunque che a questo si aggiunge sempre un'ulteriore
condizione e cioè che non ci siano processi che stiano ancora lavorando sul il
file. Come vedremo in sez.~\ref{sec:file_unix_interface} il kernel una tabella
di tutti file aperti da ciascun processo, che a sua volta contiene i
riferimenti agli \textit{inode} ad essi relativi. Prima di procedere alla
cancellazione dello spazio occupato su disco dal contenuto di un file il
kernel controlla anche questa tabella, per verificare che anche in essa non ci
sia più nessun riferimento all'\textit{inode} in questione, assicurandosi con
questo che nessun processo stia ancora usando il file.

Nel caso di socket, \textit{fifo} o file di dispositivo la funzione rimuove il
nome, e come per i file normali i processi che hanno aperto uno di questi
oggetti possono continuare ad utilizzarli.  Nel caso di cancellazione di un
\textit{link} simbolico, che consiste solo nel rimando ad un altro file,
questo viene immediatamente eliminato e non sarà più utilizzabile. 

Per cancellare una voce in una directory è necessario avere il permesso di
scrittura su di essa, dato che si va a rimuovere una voce dal suo contenuto, e
il diritto di esecuzione/attraversamento sulla directory che la contiene
(affronteremo in dettaglio l'argomento dei permessi di file e directory in
sez.~\ref{sec:file_access_control}). Se inoltre per la directory è impostato
lo \textit{sticky bit} (vedi sez.~\ref{sec:file_special_perm}), occorrerà
anche essere proprietari del file o proprietari della directory o avere i
privilegi di amministratore.

Questa caratteristica del sistema, che consente di usare un file anche se lo
si è ``cancellato'', può essere usata per essere sicuri di non lasciare file
temporanei su disco in caso di uscita imprevista di un programma. La tecnica è
quella di aprire un nuovo file e chiamare \func{unlink} su di esso subito
dopo, in questo modo il contenuto del file sarà sempre disponibile all'interno
del processo attraverso il suo file descriptor (vedi sez.~\ref{sec:file_fd}),
ma non ne resterà traccia in nessuna directory, inoltre lo spazio occupato su
disco verrà immediatamente rilasciato alla conclusione del processo, quando
tutti i file vengono chiusi.

Al contrario di quanto avviene con altri Unix, in Linux non è possibile usare
la funzione \func{unlink} sulle directory, che in tal caso fallisce con un
errore di \errcode{EISDIR}. Per cancellare una directory si deve usare la
apposita funzione di sistema \func{rmdir} (che vedremo in
sez.~\ref{sec:file_dir_creat_rem}), oppure la funzione \func{remove}.

Quest'ultima è la funzione prevista dallo standard ANSI C per effettuare una
cancellazione generica di un file o di una directory e viene usata in generale
anche per i sistemi operativi che non supportano gli \textit{hard link}. Nei
sistemi unix-like \funcd{remove} è equivalente ad usare in maniera trasparente
\func{unlink} per i file ed \func{rmdir} per le directory; il suo prototipo è:

\begin{funcproto}{
\fhead{stdio.h}
\fdecl{int remove(const char *pathname)}
\fdesc{Cancella un file o una directory.} 
}
{La funzione ritorna $0$ in caso di successo e $-1$ per un errore, nel qual
  caso \var{errno} assumerà uno dei valori relativi alla chiamata utilizzata,
  pertanto si può fare riferimento a quanto illustrato nelle descrizioni di
  \func{unlink} e \func{rmdir}.}
\end{funcproto}

La funzione utilizza la funzione \func{unlink} per cancellare i file (e si
applica anche a link simbolici, socket, \textit{fifo} e file di dispositivo) e
la funzione \func{rmdir} (vedi sez.~\ref{sec:file_dir_creat_rem}) per
cancellare le directory.\footnote{questo vale usando la \acr{glibc}; nella
  \acr{libc4} e nella \acr{libc5} la funzione \func{remove} era un semplice
  alias alla funzione \func{unlink} e quindi non poteva essere usata per le
  directory.}  Si tenga presente che, per alcune limitazioni del protocollo
NFS, utilizzare questa funzione su file che usano questo filesystem di rete
può comportare la scomparsa di file ancora in uso.

Infine per cambiare nome ad un file o a una directory si usa la funzione di
sistema \funcd{rename},\footnote{la funzione è definita dallo standard ANSI C,
  ma si applica solo per i file, lo standard POSIX estende la funzione anche
  alle directory.} il cui prototipo è:

\begin{funcproto}{
\fhead{stdio.h}
\fdecl{int rename(const char *oldpath, const char *newpath)}
\fdesc{Rinomina un file o una directory.} 
}
{La funzione ritorna $0$ in caso di successo e $-1$ per un errore,
  nel qual caso \var{errno} assumerà uno dei valori: 
  \begin{errlist}
  \item[\errcode{EACCESS}] manca il permesso di scrittura sulle directory
    contenenti \param{oldpath} e \param{newpath} o di attraversare 
    il loro \textit{pathname} o di scrivere su \param{newpath}
    se questa è una directory.
  \item[\errcode{EBUSY}] o \param{oldpath} o \param{newpath} sono in uso da
    parte di qualche processo (come directory di lavoro o come radice) o del
    sistema (come \textit{mount point}) ed il sistema non riesce a risolvere
    la situazione.
  \item[\errcode{EEXIST}] \param{newpath} è una directory già esistente e
    non è vuota (anche \errcode{ENOTEMPTY}).
  \item[\errcode{EINVAL}] \param{newpath} contiene un prefisso di
    \param{oldpath} o più in generale si è cercato di creare una directory come
    sotto-directory di sé stessa.
  \item[\errcode{EISDIR}] \param{newpath} è una directory mentre
    \param{oldpath} non è una directory.
  \item[\errcode{ENOTDIR}] uno dei componenti dei \textit{pathname} non è una
    directory o \param{oldpath} è una directory e 
    \param{newpath} esiste e non è una directory.
  \item[\errval{EPERM}] la directory contenente \param{oldpath} o quella
    contenente un \param{newpath} esistente hanno lo \textit{sticky bit} e non
    si è i proprietari dei rispettivi file (o non si hanno privilegi
    amministrativi) oppure il filesystem non supporta l'operazione. 
  \item[\errcode{EXDEV}] \param{oldpath} e \param{newpath} non sono sullo
    stesso filesystem e sotto lo stesso \textit{mount point}. 
  \end{errlist} ed inoltre \errval{EFAULT}, \errval{ELOOP}, \errval{EMLINK},
  \errval{ENAMETOOLONG}, \errval{ENOENT}, \errval{ENOMEM}, \errval{ENOSPC} e
  \errval{EROFS} nel loro significato generico.}
\end{funcproto}

La funzione rinomina in \param{newpath} il file o la directory indicati
dall'argomento \param{oldpath}. Il nome viene eliminato nella directory
originale e ricreato nella directory di destinazione mantenendo il riferimento
allo stesso \textit{inode}. Non viene spostato nessun dato e l'\textit{inode}
del file non subisce nessuna modifica in quanto le modifiche sono eseguite
sulle directory che contengono \param{newpath} ed \param{oldpath}.

Il vantaggio nell'uso di questa funzione al posto della chiamata successiva di
\func{link} e \func{unlink} è che l'operazione è eseguita atomicamente, non
c'è modifica, per quanto temporanea, al \textit{link count} del file e non può
esistere un istante in cui un altro processo possa trovare attivi entrambi i
nomi per lo stesso file se la destinazione non esiste o in cui questa sparisca
temporaneamente se già esiste.

Dato che opera in maniera analoga la funzione è soggetta alle stesse
restrizioni di \func{link}, quindi è necessario che \param{oldpath}
e \param{newpath} siano nello stesso filesystem e facciano riferimento allo
stesso \textit{mount point}, e che il filesystem supporti questo tipo di
operazione. Qualora questo non avvenga si dovrà effettuare l'operazione in
maniera non atomica copiando il file a destinazione e poi cancellando
l'originale.

Il comportamento della funzione è diverso a seconda che si voglia rinominare
un file o una directory. Se ci riferisce ad un file allora \param{newpath}, se
esiste, non deve essere una directory, altrimenti si avrà un errore di
\errcode{EISDIR}. Se \param{newpath} indica un file già esistente questo verrà
rimpiazzato atomicamente, ma nel caso in cui \func{rename} fallisca il kernel
assicura che esso non sarà toccato. I caso di sovrascrittura però potrà
esistere una breve finestra di tempo in cui sia \param{oldpath}
che \param{newpath} potranno fare entrambi riferimento al file che viene
rinominato.

Se \param{oldpath} è una directory allora \param{newpath}, se esistente, deve
essere una directory vuota, altrimenti si avranno gli errori \errcode{ENOTDIR}
(se non è una directory) o \errcode{ENOTEMPTY} o \errcode{EEXIST} (se non è
vuota). Chiaramente \param{newpath} non potrà contenere \param{oldpath} nel
suo \textit{pathname}, non essendo possibile rendere una directory una
sottodirectory di sé stessa, se questo fosse il caso si otterrebbe un errore
di \errcode{EINVAL}.

Se \param{oldpath} si riferisce ad un collegamento simbolico questo sarà
rinominato restando tale senza nessun effetto sul file a cui fa riferimento.
Se invece \param{newpath} esiste ed è un collegamento simbolico verrà
cancellato come qualunque altro file.  Infine qualora \param{oldpath}
e \param{newpath} siano due nomi che già fanno riferimento allo stesso file lo
standard POSIX prevede che la funzione non ritorni un errore, e semplicemente
non faccia nulla, lasciando entrambi i nomi.  Linux segue questo standard,
anche se, come fatto notare dal manuale della \acr{glibc}, il comportamento
più ragionevole sarebbe quello di cancellare \param{oldpath}.

In tutti i casi si dovranno avere i permessi di scrittura nelle directory
contenenti \param{oldpath} e \param{newpath}, e nel caso \param{newpath} sia
una directory vuota già esistente anche su di essa (perché dovrà essere
aggiornata la voce ``\texttt{..}''). Se poi le directory
contenenti \param{oldpath} o \param{newpath} hanno lo \textit{sticky bit}
attivo (vedi sez.~\ref{sec:file_special_perm}) si dovrà essere i proprietari
dei file (o delle directory) che si vogliono rinominare, o avere i permessi di
amministratore.


\subsection{La creazione e la cancellazione delle directory} 
\label{sec:file_dir_creat_rem}

Benché in sostanza le directory non siano altro che dei file contenenti
elenchi di nomi con riferimenti ai rispettivi \textit{inode}, non è possibile
trattarle come file ordinari e devono essere create direttamente dal kernel
attraverso una opportuna \textit{system call}.\footnote{questo è quello che
  permette anche, attraverso l'uso del VFS, l'utilizzo di diversi formati per
  la gestione dei suddetti elenchi, dalle semplici liste a strutture complesse
  come alberi binari, hash, ecc. che consentono una ricerca veloce quando il
  numero di file è molto grande.}  La funzione di sistema usata per creare una
directory è \funcd{mkdir}, ed il suo prototipo è:

\begin{funcproto}{
\fhead{sys/stat.h}
\fhead{sys/types.h}
\fdecl{int mkdir(const char *dirname, mode\_t mode)}
\fdesc{Crea una nuova directory.} 
}
{La funzione ritorna $0$ in caso di successo e $-1$ per un errore, nel qual
  caso \var{errno} assumerà uno dei valori: 
  \begin{errlist}
  \item[\errcode{EACCES}] non c'è il permesso di scrittura per la directory in
    cui si vuole inserire la nuova directory o di attraversamento per le
    directory al di sopra di essa.
  \item[\errcode{EEXIST}] un file o una directory o un collegamento simbolico
    con quel nome esiste già.
  \item[\errcode{EMLINK}] la directory in cui si vuole creare la nuova
    directory contiene troppi file; sotto Linux questo normalmente non avviene
    perché il filesystem standard consente la creazione di un numero di file
    maggiore di quelli che possono essere contenuti nel disco, ma potendo
    avere a che fare anche con filesystem di altri sistemi questo errore può
    presentarsi.
  \item[\errcode{ENOSPC}] non c'è abbastanza spazio sul file system per creare
    la nuova directory o si è esaurita la quota disco dell'utente.
  \end{errlist}
  ed inoltre \errval{EFAULT}, \errval{ELOOP}, \errval{ENAMETOOLONG},
  \errval{ENOENT}, \errval{ENOMEM}, \errval{ENOTDIR}, \errval{EPERM},
  \errval{EROFS} nel loro significato generico.}
\end{funcproto}

La funzione crea una nuova directory vuota, che contiene cioè solo le due voci
standard presenti in ogni directory (``\file{.}'' e ``\file{..}''), con il
nome indicato dall'argomento \param{dirname}. 

I permessi di accesso (vedi sez.~\ref{sec:file_access_control}) con cui la
directory viene creata sono specificati dall'argomento \param{mode}, i cui
possibili valori sono riportati in tab.~\ref{tab:file_permission_const}; si
tenga presente che questi sono modificati dalla maschera di creazione dei file
(si veda sez.~\ref{sec:file_perm_management}).  La titolarità della nuova
directory è impostata secondo quanto illustrato in
sez.~\ref{sec:file_ownership_management}.

Come accennato in precedenza per eseguire la cancellazione di una directory è
necessaria una specifica funzione di sistema, \funcd{rmdir}, il suo prototipo
è:

\begin{funcproto}{
\fhead{sys/stat.h}
\fdecl{int rmdir(const char *dirname)}
\fdesc{Cancella una directory.} 
}
{La funzione ritorna $0$ in caso di successo e $-1$ per un errore, nel qual
  caso \var{errno} assumerà uno dei valori: 
  \begin{errlist}
  \item[\errcode{EACCES}] non c'è il permesso di scrittura per la directory
    che contiene la directory che si vuole cancellare, o non c'è il permesso
    di attraversare (esecuzione) una delle directory specificate in
    \param{dirname}.
  \item[\errcode{EBUSY}] la directory specificata è la directory di lavoro o
    la radice di qualche processo o un \textit{mount point}.
  \item[\errcode{EINVAL}] si è usato ``\texttt{.}'' come ultimo componente
    di \param{dirname}.
  \item[\errcode{EPERM}] il filesystem non supporta la cancellazione di
    directory, oppure la directory che contiene \param{dirname} ha lo
    \textit{sticky bit} impostato e non si è i proprietari della directory o
    non si hanno privilegi amministrativi.
  \end{errlist}
  ed inoltre \errval{EFAULT}, \errval{ELOOP}, \errval{ENAMETOOLONG},
  \errval{ENOENT}, \errval{ENOMEM}, \errval{ENOTDIR}, \errcode{ENOTEMPTY} e
  \errval{EROFS} nel loro significato generico.}
\end{funcproto}


La funzione cancella la directory \param{dirname}, che deve essere vuota, la
directory deve cioè contenere le due voci standard ``\file{.}'' e
``\file{..}'' e niente altro.  Il nome può essere indicato con un
\textit{pathname} assoluto o relativo, ma si deve fare riferimento al nome
nella directory genitrice, questo significa che \textit{pathname} terminanti
in ``\file{.}'' e ``\file{..}'' anche se validi in altri contesti, causeranno
il fallimento della funzione.

Inoltre per eseguire la cancellazione, oltre ad essere vuota, occorre anche
che la directory non sia utilizzata, questo vuol dire anche che non deve
essere la directory di lavoro (vedi sez.~\ref{sec:file_work_dir}) o la
directory radice (vedi sez.~\ref{sec:file_chroot}) di nessun processo, od
essere utilizzata come \textit{mount point}.

Se la directory cancellata risultasse aperta in qualche processo per una
lettura dei suoi contenuti (vedi sez.~\ref{sec:file_dir_read}), pur
scomparendo dal filesystem e non essendo più possibile accedervi o crearvi
altri file, le risorse ad essa associate verrebbero disallocate solo alla
chiusura di tutti questi ulteriori riferimenti.


\subsection{Lettura e scansione delle directory}
\label{sec:file_dir_read}

Benché le directory alla fine non siano altro che dei file che contengono
delle liste di nomi associati ai relativi \textit{inode}, per il ruolo che
rivestono nella struttura del sistema non possono essere trattate come dei
normali file di dati. Ad esempio, onde evitare inconsistenze all'interno del
filesystem, solo il kernel può scrivere il contenuto di una directory, e non
può essere un processo a inserirvi direttamente delle voci con le usuali
funzioni di scrittura.

Ma se la scrittura e l'aggiornamento dei dati delle directory è compito del
kernel, sono molte le situazioni in cui i processi necessitano di poterne
leggere il contenuto. Benché questo possa essere fatto direttamente (vedremo
in sez.~\ref{sec:file_open_close} che è possibile aprire una directory come se
fosse un file, anche se solo in sola lettura) in generale il formato con cui
esse sono scritte può dipendere dal tipo di filesystem, tanto che, come
riportato in tab.~\ref{tab:file_file_operations}, il VFS prevede una apposita
funzione per la lettura delle directory.

\itindbeg{directory~stream}

Tutto questo si riflette nello standard POSIX\footnote{le funzioni erano
  presenti in SVr4 e 4.3BSD, la loro specifica è riportata in POSIX.1-2001.}
che ha introdotto una apposita interfaccia per la lettura delle directory,
basata sui cosiddetti \textit{directory stream}, chiamati così per l'analogia
con i \textit{file stream} dell'interfaccia standard ANSI C che vedremo in
sez.~\ref{sec:files_std_interface}. La prima funzione di questa interfaccia è
\funcd{opendir}, il cui prototipo è:

\begin{funcproto}{
\fhead{sys/types.h}
\fhead{dirent.h}
\fdecl{DIR *opendir(const char *dirname)}
\fdesc{Apre un \textit{directory stream}.} 
}
{La funzione ritorna un puntatore al \textit{directory stream} in caso di
  successo e \val{NULL} per un errore, nel qual caso \var{errno} assumerà uno
  dei valori \errval{EACCES}, \errval{EMFILE}, \errval{ENFILE},
  \errval{ENOENT}, \errval{ENOMEM} e \errval{ENOTDIR} nel loro significato
  generico.}
\end{funcproto}

La funzione apre un \textit{directory stream} per la directory
\param{dirname}, ritornando il puntatore ad un oggetto di tipo \typed{DIR} (che
è il tipo opaco usato dalle librerie per gestire i \textit{directory stream})
da usare per tutte le operazioni successive, la funzione inoltre posiziona lo
\textit{stream} sulla prima voce contenuta nella directory.

Si tenga presente che comunque la funzione opera associando il
\textit{directory stream} ad un opportuno file descriptor (vedi
sez.~\ref{sec:file_fd}) sottostante, sul quale vengono compiute le
operazioni. Questo viene sempre aperto impostando il flag di
\textit{close-on-exec} (si ricordi quanto detto in sez.~\ref{sec:proc_exec}),
così da evitare che resti aperto in caso di esecuzione di un altro programma.

Nel caso in cui sia necessario conoscere il file descriptor associato ad un
\textit{directory stream} si può usare la funzione
\funcd{dirfd},\footnote{questa funzione è una estensione introdotta con BSD
  4.3-Reno ed è presente in Linux con le libc5 (a partire dalla versione
  5.1.2) e con la \acr{glibc} ma non presente in POSIX fino alla revisione
  POSIX.1-2008, per questo per poterla utilizzare fino alla versione 2.10
  della \acr{glibc} era necessario definire le macro \macro{\_BSD\_SOURCE} o
  \macro{\_SVID\_SOURCE}, dalla versione 2.10 si possono usare anche
  \texttt{\macro{\_POSIX\_C\_SOURCE} >= 200809L} o
  \texttt{\macro{\_XOPEN\_SOURCE} >= 700}.}  il cui prototipo è:

\begin{funcproto}{
\fhead{sys/types.h}
\fhead{dirent.h}
\fdecl{int dirfd(DIR *dir)}
\fdesc{Legge il file descriptor associato ad un \textit{directory stream}.} 
}
{La funzione ritorna un valore positivo corrispondente al file descriptor in
  caso di successo e $-1$ per un errore, nel qual caso \var{errno} assumerà
  uno dei valori:
  \begin{errlist}
  \item[\errcode{EINVAL}] \param{dir} non è un puntatore ad un
    \textit{directory stream}. 
  \item[\errcode{ENOTSUP}] l'implementazione non supporta l'uso di un file
    descriptor per la directory.
  \end{errlist}
}
\end{funcproto}

La funzione restituisce il file descriptor associato al \textit{directory
  stream} \param{dir}. Di solito si utilizza questa funzione in abbinamento a
funzioni che operano sui file descriptor, ad esempio si potrà usare
\func{fstat} per ottenere le proprietà della directory, o \func{fchdir} per
spostare su di essa la directory di lavoro (vedi sez.~\ref{sec:file_work_dir}).

Viceversa se si è aperto un file descriptor corrispondente ad una directory è
possibile associarvi un \textit{directory stream} con la funzione
\funcd{fdopendir},\footnote{questa funzione è però disponibile solo a partire
  dalla versione 2.4 della \acr{glibc}, ed è stata introdotta nello standard
  POSIX solo a partire dalla revisione POSIX.1-2008, prima della versione 2.10
  della \acr{glibc} per poterla utilizzare era necessario definire la macro
  \macro{\_GNU\_SOURCE}, dalla versione 2.10 si possono usare anche
  \texttt{\macro{\_POSIX\_C\_SOURCE} >= 200809L} o \texttt{\_XOPEN\_SOURCE >=
    700} .}  il cui prototipo è:

\begin{funcproto}{
\fhead{sys/types.h}
\fhead{dirent.h}
\fdecl{DIR *fdopendir(int fd)}
\fdesc{Associa un \textit{directory stream} ad un file descriptor.} 
}
{La funzione ritorna un puntatore al \textit{directory stream} in caso di
  successo e \val{NULL} per un errore, nel qual caso \var{errno} assumerà uno
  dei valori \errval{EBADF} o \errval{ENOMEM} nel loro significato generico.}
\end{funcproto}

La funzione è identica a \func{opendir}, ma ritorna un \textit{directory
  stream} facendo riferimento ad un file descriptor \param{fd} che deve essere
stato aperto in precedenza; la funzione darà un errore qualora questo non
corrisponda ad una directory. L'uso di questa funzione permette di rispondere
agli stessi requisiti delle funzioni ``\textit{at}'' che vedremo in
sez.~\ref{sec:file_openat}.

Una volta utilizzata il file descriptor verrà usato internamente dalle
funzioni che operano sul \textit{directory stream} e non dovrà essere più
utilizzato all'interno del proprio programma. In particolare dovrà essere
chiuso attraverso il \textit{directory stream} con \func{closedir} e non
direttamente. Si tenga presente inoltre che \func{fdopendir} non modifica lo
stato di un eventuale flag di \textit{close-on-exec}, che pertanto dovrà
essere impostato esplicitamente in fase di apertura del file descriptor.

Una volta che si sia aperto un \textit{directory stream} la lettura del
contenuto della directory viene effettuata attraverso la funzione
\funcd{readdir}, il cui prototipo è:

\begin{funcproto}{
\fhead{sys/types.h}
\fhead{dirent.h}
\fdecl{struct dirent *readdir(DIR *dir)}
\fdesc{Legge una voce dal \textit{directory stream}.} 
}
{La funzione ritorna il puntatore alla struttura contenente i dati in caso di
  successo e \val{NULL} per un errore o se si è raggiunta la fine dello
  \textit{stream}. Il solo codice di errore restituito in \var{errno} è
  \errval{EBADF} qualora \param{dir} non indichi un \textit{directory stream}
  valido.}
\end{funcproto}

La funzione legge la voce corrente nella directory, posizionandosi sulla voce
successiva. Pertanto se si vuole leggere l'intero contenuto di una directory
occorrerà ripetere l'esecuzione della funzione fintanto che non si siano
esaurite tutte le voci in essa presenti, che viene segnalata dalla
restituzione di \val{NULL} come valore di ritorno. Si può distinguere questa
condizione da un errore in quanto in questo caso \var{errno} non verrebbe
modificata.

I dati letti da \func{readdir} vengono memorizzati in una struttura
\struct{dirent}, la cui definizione è riportata in
fig.~\ref{fig:file_dirent_struct}.\footnote{la definizione è quella usata da
  Linux, che si trova nel file \file{/usr/include/bits/dirent.h}, essa non
  contempla la presenza del campo \var{d\_namlen} che indica la lunghezza del
  nome del file.} La funzione non è rientrante e restituisce il puntatore ad
una struttura allocata staticamente, che viene sovrascritta tutte le volte che
si ripete la lettura di una voce sullo stesso \textit{directory stream}.

Di questa funzione esiste anche una versione rientrante,
\funcd{readdir\_r},\footnote{per usarla è necessario definire una qualunque
  delle macro \texttt{\macro{\_POSIX\_C\_SOURCE} >= 1},
  \macro{\_XOPEN\_SOURCE}, \macro{\_BSD\_SOURCE}, \macro{\_SVID\_SOURCE},
  \macro{\_POSIX\_SOURCE}.} che non usa una struttura allocata staticamente, e
può essere utilizzata anche con i \textit{thread}, il suo prototipo è:

\begin{funcproto}{
\fhead{sys/types.h}
\fhead{dirent.h}
\fdecl{int readdir\_r(DIR *dir, struct dirent *entry, struct dirent **result)}
\fdesc{Legge una voce dal \textit{directory stream}.} 
}
{La funzione ritorna $0$ in caso di successo ed un numero positivo per un
  errore, nel qual caso \var{errno} assumerà gli stessi valori di
  \func{readdir}.} 
\end{funcproto}

La funzione restituisce in \param{result} come \textit{value result argument}
l'indirizzo della struttura \struct{dirent} dove sono stati salvati i dati,
che deve essere allocata dal chiamante, ed il cui indirizzo deve essere
indicato con l'argomento \param{entry}.  Se si è raggiunta la fine del
\textit{directory stream} invece in \param{result} viene restituito il valore
\val{NULL}.

\begin{figure}[!htb]
  \footnotesize \centering
  \begin{minipage}[c]{0.8\textwidth}
    \includestruct{listati/dirent.c}
  \end{minipage} 
  \normalsize 
  \caption{La struttura \structd{dirent} per la lettura delle informazioni dei 
    file.}
  \label{fig:file_dirent_struct}
\end{figure}

% Lo spazio per la \struct{dirent} dove vengono restituiti i dati della
% directory deve essere allocato a cura del chiamante, dato che la dimensione


I vari campi di \struct{dirent} contengono le informazioni relative alle voci
presenti nella directory. Sia BSD che SVr4 che POSIX.1-2001\footnote{il
  vecchio standard POSIX prevedeva invece solo la presenza del campo
  \var{d\_fileno}, identico \var{d\_ino}, che in Linux era definito come alias
  di quest'ultimo, mentre il campo \var{d\_name} era considerato dipendente
  dall'implementazione.}  prevedono che siano sempre presenti il campo
\var{d\_name}, che contiene il nome del file nella forma di una stringa
terminata da uno zero, ed il campo \var{d\_ino}, che contiene il numero di
\textit{inode} cui il file è associato e corrisponde al campo \var{st\_ino} di
\struct{stat}.  La presenza di ulteriori campi opzionali oltre i due citati è
segnalata dalla definizione di altrettante macro nella forma
\code{\_DIRENT\_HAVE\_D\_XXX} dove \code{XXX} è il nome del relativo
campo. Come si può evincere da fig.~\ref{fig:file_dirent_struct} nel caso di
Linux sono pertanto definite le macro \macrod{\_DIRENT\_HAVE\_D\_TYPE},
\macrod{\_DIRENT\_HAVE\_D\_OFF} e \macrod{\_DIRENT\_HAVE\_D\_RECLEN}, mentre non
è definita la macro \macrod{\_DIRENT\_HAVE\_D\_NAMLEN}.

Dato che possono essere presenti campi opzionali e che lo standard
POSIX.1-2001 non specifica una dimensione definita per il nome dei file (che
può variare a seconda del filesystem), ma solo un limite superiore pari a
\const{NAME\_MAX} (vedi tab.~\ref{tab:sys_file_macro}), in generale per
allocare una struttura \struct{dirent} in maniera portabile occorre eseguire
un calcolo per ottenere le dimensioni appropriate per il proprio
sistema.\footnote{in SVr4 la lunghezza del campo è definita come
  \code{NAME\_MAX+1} che di norma porta al valore di 256 byte usato anche in
  fig.~\ref{fig:file_dirent_struct}.} Lo standard però richiede che il campo
\var{d\_name} sia sempre l'ultimo della struttura, questo ci consente di
ottenere la dimensione della prima parte con la macro di utilità generica
\macro{offsetof}, che si può usare con il seguente prototipo:

{\centering
\vspace{3pt}
\begin{funcbox}{
\fhead{stddef.h}
\fdecl{size\_t \macrod{offsetof}(type, member)}
\fdesc{Restituisce la posizione del campo \param{member} nella
  struttura \param{type}.}
} 
\end{funcbox}
}

Ottenuta allora con \code{offsetof(struct dirent, d\_name)} la dimensione
della parte iniziale della struttura, basterà sommarci la dimensione massima
dei nomi dei file nel filesystem che si sta usando, che si può ottenere
attraverso la funzione \func{pathconf} (per la quale si rimanda a
sez.~\ref{sec:sys_file_limits}) più un ulteriore carattere per la terminazione
della stringa.

Per quanto riguarda il significato dei campi opzionali, il campo \var{d\_type}
indica il tipo di file (se \textit{fifo}, directory, collegamento simbolico,
ecc.), e consente di evitare una successiva chiamata a \func{lstat} (vedi
sez.~\ref{sec:file_stat}) per determinarlo. I suoi possibili valori sono
riportati in tab.~\ref{tab:file_dtype_macro}. Si tenga presente che questo
valore è disponibile solo per i filesystem che ne supportano la restituzione
(fra questi i più noti sono \acr{Btrfs}, \acr{ext2}, \acr{ext3}, e
\acr{ext4}), per gli altri si otterrà sempre il valore
\const{DT\_UNKNOWN}.\footnote{inoltre fino alla versione 2.1 della
  \acr{glibc}, pur essendo il campo \var{d\_type} presente, il suo uso non era
  implementato, e veniva restituito comunque il valore \const{DT\_UNKNOWN}.}

\begin{table}[htb]
  \centering
  \footnotesize
  \begin{tabular}[c]{|l|l|}
    \hline
    \textbf{Valore} & \textbf{Tipo di file} \\
    \hline
    \hline
    \constd{DT\_UNKNOWN} & Tipo sconosciuto.\\
    \constd{DT\_REG}     & File normale.\\
    \constd{DT\_DIR}     & Directory.\\
    \constd{DT\_LNK}     & Collegamento simbolico.\\
    \constd{DT\_FIFO}    & \textit{Fifo}.\\
    \constd{DT\_SOCK}    & Socket.\\
    \constd{DT\_CHR}     & Dispositivo a caratteri.\\
    \constd{DT\_BLK}     & Dispositivo a blocchi.\\
    \hline    
  \end{tabular}
  \caption{Costanti che indicano i vari tipi di file nel campo \var{d\_type}
    della struttura \struct{dirent}.}
  \label{tab:file_dtype_macro}
\end{table}

Per la conversione da e verso l'analogo valore mantenuto dentro il campo
\var{st\_mode} di \struct{stat} (vedi fig.~\ref{fig:file_stat_struct}) sono
definite anche due macro di conversione, \macro{IFTODT} e \macro{DTTOIF}:

{\centering
\vspace{3pt}
\begin{funcbox}{
\fhead{dirent.h}
\fdecl{int \macrod{IFTODT}(mode\_t MODE)}
\fdesc{Converte il tipo di file dal formato di \var{st\_mode} a quello di
  \var{d\_type}.}
\fdecl{mode\_t \macrod{DTTOIF}(int DTYPE)}
\fdesc{Converte il tipo di file dal formato di \var{d\_type} a quello di
  \var{st\_mode}.}  
} 
\end{funcbox}
}

Il campo \var{d\_off} contiene invece la posizione della voce successiva della
directory, mentre il campo \var{d\_reclen} la lunghezza totale della voce
letta. Con questi due campi diventa possibile, determinando la posizione delle
varie voci, spostarsi all'interno dello \textit{stream} usando la funzione
\funcd{seekdir},\footnote{sia questa funzione che \func{telldir}, sono
  estensioni prese da BSD, ed introdotte nello standard POSIX solo a partire
  dalla revisione POSIX.1-2001, per poterle utilizzare deve essere definita
  una delle macro \macro{\_XOPEN\_SOURCE}, \macro{\_BSD\_SOURCE} o
  \macro{\_SVID\_SOURCE}.} il cui prototipo è:

\begin{funcproto}{
\fhead{dirent.h}
\fdecl{void seekdir(DIR *dir, off\_t offset)}
\fdesc{Cambia la posizione all'interno di un \textit{directory stream}.} 
}
{La funzione non ritorna niente e non imposta errori.}
\end{funcproto}

La funzione non ritorna nulla e non segnala errori, è però necessario che il
valore dell'argomento \param{offset} sia valido per lo
\textit{stream} \param{dir}; esso pertanto deve essere stato ottenuto o dal
valore di \var{d\_off} di \struct{dirent} o dal valore restituito dalla
funzione \funcd{telldir}, che legge la posizione corrente; il cui prototipo
è:\footnote{prima della \acr{glibc} 2.1.1 la funzione restituiva un valore di
  tipo \type{off\_t}, sostituito a partire dalla versione 2.1.2 da \ctyp{long}
  per conformità a POSIX.1-2001.}

\begin{funcproto}{
\fhead{dirent.h}
\fdecl{long telldir(DIR *dir)}
\fdesc{Ritorna la posizione corrente in un \textit{directory stream}.} 
}
{La funzione ritorna la posizione corrente nello \textit{stream} (un numero
  positivo) in caso di successo e $-1$ per un errore, nel qual caso
  \var{errno} assume solo il valore di \errval{EBADF}, corrispondente ad un
  valore errato per \param{dir}.  }
\end{funcproto}

La sola funzione di posizionamento per un \textit{directory stream} prevista
originariamente dallo standard POSIX è \funcd{rewinddir}, che riporta la
posizione a quella iniziale; il suo prototipo è:

\begin{funcproto}{
\fhead{sys/types.h}
\fhead{dirent.h}
\fdecl{void rewinddir(DIR *dir)}
\fdesc{Si posiziona all'inizio di un \textit{directory stream}.} 
}
{La funzione non ritorna niente e non imposta errori.}
\end{funcproto}

Una volta completate le operazioni si può chiudere il \textit{directory
  stream}, ed il file descriptor ad esso associato, con la funzione
\funcd{closedir}, il cui prototipo è:

\begin{funcproto}{
\fhead{sys/types.h}
\fhead{dirent.h}
\fdecl{int closedir(DIR *dir)}
\fdesc{Chiude un \textit{directory stream}.} 
}
{La funzione ritorna $0$ in caso di successo e $-1$ per un errore, nel qual
  caso \var{errno} assume solo il valore \errval{EBADF}.}
\end{funcproto}

A parte queste funzioni di base in BSD 4.3 venne introdotta un'altra funzione
che permette di eseguire una scansione completa, con tanto di ricerca ed
ordinamento, del contenuto di una directory; la funzione è
\funcd{scandir}\footnote{in Linux questa funzione è stata introdotta fin dalle
  \acr{libc4} e richiede siano definite le macro \macro{\_BSD\_SOURCE} o
  \macro{\_SVID\_SOURCE}.} ed il suo prototipo è:

\begin{funcproto}{
\fhead{dirent.h}
\fdecl{int scandir(const char *dir, struct dirent ***namelist, \\
\phantom{int scandir(}int(*filter)(const struct dirent *), \\
\phantom{int scandir(}int(*compar)(const struct dirent **, const struct dirent **))}
\fdesc{Esegue una scansione di un \textit{directory stream}.} 
}
{La funzione ritorna il numero di voci trovate in caso di successo e $-1$ per
  un errore, nel qual caso \var{errno} può assumere solo il valore
  \errval{ENOMEM}.}
\end{funcproto}

Al solito, per la presenza fra gli argomenti di due puntatori a funzione, il
prototipo non è molto comprensibile; queste funzioni però sono quelle che
controllano rispettivamente la selezione di una voce, passata con
l'argomento \param{filter}, e l'ordinamento di tutte le voci selezionate,
specificata dell'argomento \param{compar}.

La funzione legge tutte le voci della directory indicata dall'argomento
\param{dir}, passando un puntatore a ciascuna di esse (una struttura
\struct{dirent}) come argomento della funzione di selezione specificata da
\param{filter}; se questa ritorna un valore diverso da zero il puntatore viene
inserito in un vettore che viene allocato dinamicamente con \func{malloc}.
Qualora si specifichi un valore \val{NULL} per l'argomento \param{filter} non
viene fatta nessuna selezione e si ottengono tutte le voci presenti.

Le voci selezionate possono essere riordinate tramite \funcm{qsort}, le
modalità del riordinamento possono essere personalizzate usando la funzione
\param{compar} come criterio di ordinamento di \funcm{qsort}, la funzione
prende come argomenti le due strutture \struct{dirent} da confrontare
restituendo un valore positivo, nullo o negativo per indicarne l'ordinamento;
alla fine l'indirizzo della lista ordinata dei puntatori alle strutture
\struct{dirent} viene restituito nell'argomento
\param{namelist}.\footnote{la funzione alloca automaticamente la lista, e
  restituisce, come \textit{value result argument}, l'indirizzo della stessa;
  questo significa che \param{namelist} deve essere dichiarato come
  \code{struct dirent **namelist} ed alla funzione si deve passare il suo
  indirizzo.}

\itindend{directory~stream}

Per l'ordinamento, vale a dire come valori possibili per l'argomento
\param{compar}, sono disponibili due funzioni predefinite, \funcd{alphasort} e
\funcd{versionsort}, i cui prototipi sono:

\begin{funcproto}{
\fhead{dirent.h}
\fdecl{int alphasort(const void *a, const void *b)}
\fdecl{int versionsort(const void *a, const void *b)}
\fdesc{Riordinano le voci di \textit{directory stream}.} 
}
{Le funzioni restituiscono un valore minore, uguale o maggiore di zero qualora
  il primo argomento sia rispettivamente minore, uguale o maggiore del secondo
  e non forniscono errori.}
\end{funcproto}

La funzione \func{alphasort} deriva da BSD ed è presente in Linux fin dalle
\acr{libc4}\footnote{la versione delle \acr{libc4} e \acr{libc5} usa però come
  argomenti dei puntatori a delle strutture \struct{dirent}; la glibc usa il
  prototipo originario di BSD, mostrato anche nella definizione, che prevede
  puntatori a \ctyp{void}.} e deve essere specificata come argomento
\param{compar} per ottenere un ordinamento alfabetico secondo il valore del
campo \var{d\_name} delle varie voci. La \acr{glibc} prevede come
estensione\footnote{la \acr{glibc}, a partire dalla versione 2.1, effettua
  anche l'ordinamento alfabetico tenendo conto delle varie localizzazioni,
  usando \funcm{strcoll} al posto di \funcm{strcmp}.} anche
\func{versionsort}, che ordina i nomi tenendo conto del numero di versione,
cioè qualcosa per cui \texttt{file10} viene comunque dopo \texttt{file4}.

\begin{figure}[!htb]
  \footnotesize \centering
  \begin{minipage}[c]{\codesamplewidth}
    \includecodesample{listati/my_ls.c}
  \end{minipage}
  \caption{Esempio di codice per eseguire la lista dei file contenuti in una
    directory.} 
  \label{fig:file_my_ls}
\end{figure}

Un semplice esempio dell'uso di queste funzioni è riportato in
fig.~\ref{fig:file_my_ls}, dove si è riportata la sezione principale di un
programma che, usando la funzione di scansione illustrata in
fig.~\ref{fig:file_dirscan}, stampa i nomi dei file contenuti in una directory
e la relativa dimensione, in sostanza una versione semplificata del comando
\cmd{ls}.

Il programma è estremamente semplice; in fig.~\ref{fig:file_my_ls} si è omessa
la parte di gestione delle opzioni, che prevede solo l'uso di una funzione per
la stampa della sintassi, anch'essa omessa, ma il codice completo può essere
trovato coi sorgenti allegati alla guida nel file \file{myls.c}.

In sostanza tutto quello che fa il programma, dopo aver controllato
(\texttt{\small 12-15}) di avere almeno un argomento, che indicherà la
directory da esaminare, è chiamare (\texttt{\small 16}) la funzione
\myfunc{dir\_scan} per eseguire la scansione, usando la funzione \code{do\_ls}
(\texttt{\small 22-29}) per fare tutto il lavoro.

Quest'ultima si limita (\texttt{\small 26}) a chiamare \func{stat} sul file
indicato dalla directory entry passata come argomento (il cui nome è appunto
\var{direntry->d\_name}), memorizzando in una opportuna struttura \var{data} i
dati ad esso relativi, per poi provvedere (\texttt{\small 27}) a stampare il
nome del file e la dimensione riportata in \var{data}.

Dato che la funzione verrà chiamata all'interno di \myfunc{dir\_scan} per ogni
voce presente, questo è sufficiente a stampare la lista completa dei file e
delle relative dimensioni. Si noti infine come si restituisca sempre 0 come
valore di ritorno per indicare una esecuzione senza errori.

\begin{figure}[!htb]
  \footnotesize \centering
  \begin{minipage}[c]{\codesamplewidth}
    \includecodesample{listati/dir_scan.c}
  \end{minipage}
  \caption{Codice della funzione di scansione di una directory contenuta nel
    file \file{dir\_scan.c}.} 
  \label{fig:file_dirscan}
\end{figure}

Tutto il grosso del lavoro è svolto dalla funzione \myfunc{dir\_scan},
riportata in fig.~\ref{fig:file_dirscan}. La funzione è volutamente generica e
permette di eseguire una funzione, passata come secondo argomento, su tutte le
voci di una directory.  La funzione inizia con l'aprire (\texttt{\small
  18-22}) uno \textit{stream} sulla directory passata come primo argomento,
stampando un messaggio in caso di errore.

Il passo successivo (\texttt{\small 23-24}) è cambiare directory di lavoro
(vedi sez.~\ref{sec:file_work_dir}), usando in sequenza le funzioni
\func{dirfd} e \func{fchdir} (in realtà si sarebbe potuto usare direttamente
\func{chdir} su \var{dirname}), in modo che durante il successivo ciclo
(\texttt{\small 26-30}) sulle singole voci dello \textit{stream} ci si trovi
all'interno della directory.\footnote{questo è essenziale al funzionamento
  della funzione \code{do\_ls}, e ad ogni funzione che debba usare il campo
  \var{d\_name}, in quanto i nomi dei file memorizzati all'interno di una
  struttura \struct{dirent} sono sempre relativi alla directory in questione,
  e senza questo posizionamento non si sarebbe potuto usare \func{stat} per
  ottenere le dimensioni.}

Avendo usato lo stratagemma di fare eseguire tutte le manipolazioni necessarie
alla funzione passata come secondo argomento, il ciclo di scansione della
directory è molto semplice; si legge una voce alla volta (\texttt{\small 26})
all'interno di una istruzione di \code{while} e fintanto che si riceve una
voce valida, cioè un puntatore diverso da \val{NULL}, si esegue
(\texttt{\small 27}) la funzione di elaborazione \var{compare} (che nel nostro
caso sarà \code{do\_ls}), ritornando con un codice di errore (\texttt{\small
  28}) qualora questa presenti una anomalia, identificata da un codice di
ritorno negativo. Una volta terminato il ciclo la funzione si conclude con la
chiusura (\texttt{\small 31}) dello \textit{stream}\footnote{nel nostro caso,
  uscendo subito dopo la chiamata, questo non servirebbe, in generale però
  l'operazione è necessaria, dato che la funzione può essere invocata molte
  volte all'interno dello stesso processo, per cui non chiudere i
  \textit{directory stream} comporterebbe un consumo progressivo di risorse,
  con conseguente rischio di esaurimento delle stesse.} e la restituzione
(\texttt{\small 32}) del codice di operazioni concluse con successo.



\subsection{La directory di lavoro}
\label{sec:file_work_dir}

\index{directory~di~lavoro|(} 

Come accennato in sez.~\ref{sec:proc_fork} a ciascun processo è associata una
directory nell'albero dei file,\footnote{questa viene mantenuta all'interno
  dei dati della sua \kstruct{task\_struct} (vedi
  fig.~\ref{fig:proc_task_struct}), più precisamente nel campo \texttt{pwd}
  della sotto-struttura \kstruct{fs\_struct}.} che è chiamata
\textsl{directory corrente} o \textsl{directory di lavoro} (in inglese
\textit{current working directory}).  La directory di lavoro è quella da cui
si parte quando un \textit{pathname} è espresso in forma relativa, dove il
``\textsl{relativa}'' fa riferimento appunto a questa directory.

Quando un utente effettua il login, questa directory viene impostata alla
\textit{home directory} del suo account. Il comando \cmd{cd} della shell
consente di cambiarla a piacere, spostandosi da una directory ad un'altra, il
comando \cmd{pwd} la stampa sul terminale.  Siccome la directory di lavoro
resta la stessa quando viene creato un processo figlio (vedi
sez.~\ref{sec:proc_fork}), la directory di lavoro della shell diventa anche la
directory di lavoro di qualunque comando da essa lanciato.

Dato che è il kernel che tiene traccia dell'\textit{inode} della directory di
lavoro di ciascun processo, per ottenerne il \textit{pathname} occorre usare
una apposita funzione, \funcd{getcwd},\footnote{con Linux \func{getcwd} è una
  \textit{system call} dalla versione 2.1.9, in precedenza il valore doveva
  essere ottenuto tramite il filesystem \texttt{/proc} da
  \procfile{/proc/self/cwd}.} il cui prototipo è:

\begin{funcproto}{
\fhead{unistd.h}
\fdecl{char *getcwd(char *buffer, size\_t size)}
\fdesc{Legge il \textit{pathname} della directory di lavoro corrente.} 
}
{La funzione ritorna il puntatore a \param{buffer} in caso di successo e
  \val{NULL} per un errore, nel qual caso \var{errno} assumerà uno dei valori: 
  \begin{errlist}
  \item[\errcode{EACCES}] manca il permesso di lettura o di attraversamento  su
    uno dei componenti del \textit{pathname} (cioè su una delle directory
    superiori alla corrente).
  \item[\errcode{EINVAL}] l'argomento \param{size} è zero e \param{buffer} non
    è nullo.
  \item[\errcode{ENOENT}] la directory di lavoro è stata eliminata. 
  \item[\errcode{ERANGE}] l'argomento \param{size} è più piccolo della
    lunghezza del \textit{pathname}. 
  \end{errlist}
  ed inoltre \errcode{EFAULT} ed \errcode{ENOMEM} nel loro significato
  generico.}
\end{funcproto}

La funzione restituisce il \textit{pathname} completo della directory di
lavoro corrente nella stringa puntata da \param{buffer}, che deve essere
precedentemente allocata, per una dimensione massima di \param{size}.  Il
buffer deve essere sufficientemente largo da poter contenere il
\textit{pathname} completo più lo zero di terminazione della stringa. Qualora
esso ecceda le dimensioni specificate con \param{size} la funzione restituisce
un errore.

A partire dal kernel Linux 2.6.36 il nome può avere come prefisso la stringa
\texttt{(unreachable)} se la directory di lavoro resta fuori dalla directory
radice del processo dopo un \func{chroot} (torneremo su questi argomenti in
sez.~\ref{sec:file_chroot}); pertanto è sempre opportuno controllare il primo
carattere della stringa restituita dalla funzione per evitare di interpretare
mare un \textit{pathname} irraggiungibile.

Come estensione allo standard POSIX.1, supportata da Linux e dalla
\acr{glibc}, si può anche specificare un puntatore nullo come \param{buffer}
nel qual caso la stringa sarà allocata automaticamente per una dimensione pari
a \param{size} qualora questa sia diversa da zero, o della lunghezza esatta
del \textit{pathname} altrimenti. In questo caso ci si deve ricordare di
disallocare la stringa con \func{free} una volta cessato il suo utilizzo.

Un uso comune di \func{getcwd} è quello di salvarsi la directory di lavoro
all'avvio del programma per poi potervi tornare in un tempo successivo, un
metodo alternativo più veloce, se non si è a corto di file descriptor, è
invece quello di aprire all'inizio la directory corrente (vale a dire
``\texttt{.}'') e tornarvi in seguito con \func{fchdir}.

Di questa funzione esiste una versione alternativa per compatibilità
all'indietro con BSD, \funcm{getwd}, che non prevede l'argomento \param{size}
e quindi non consente di specificare la dimensione di \param{buffer} che
dovrebbe essere allocato in precedenza ed avere una dimensione sufficiente
(per BSD maggiore \const{PATH\_MAX}, che di solito 256 byte, vedi
sez.~\ref{sec:sys_limits}). Il problema è che su Linux non esiste una
dimensione superiore per la lunghezza di un \textit{pathname}, per cui non è
detto che il buffer sia sufficiente a contenere il nome del file, e questa è
la ragione principale per cui questa funzione è deprecata, e non la tratteremo.

Una seconda funzione usata per ottenere la directory di lavoro è
\funcm{get\_current\_dir\_name} (la funzione è una estensione GNU e presente
solo nella \acr{glibc}) che non prende nessun argomento ed è sostanzialmente
equivalente ad una \code{getcwd(NULL, 0)}, con la differenza che se
disponibile essa ritorna il valore della variabile di ambiente \envvar{PWD},
che essendo costruita dalla shell può contenere un \textit{pathname}
comprendente anche dei collegamenti simbolici. Usando \func{getcwd} infatti,
essendo il \textit{pathname} ricavato risalendo all'indietro l'albero della
directory, si perderebbe traccia di ogni passaggio attraverso eventuali
collegamenti simbolici.

Per cambiare la directory di lavoro si può usare la funzione di sistema
\funcd{chdir}, equivalente del comando di shell \cmd{cd}, il cui nome sta
appunto per \textit{change directory}, il suo prototipo è:

\begin{funcproto}{
\fhead{unistd.h}
\fdecl{int chdir(const char *pathname)}
\fdesc{Cambia la directory di lavoro per \textit{pathname}.} 
}
{La funzione ritorna $0$ in caso di successo e $-1$ per un errore, nel qual
  caso \var{errno} assumerà uno dei valori: 
  \begin{errlist}
  \item[\errcode{EACCES}] manca il permesso di ricerca su uno dei componenti
    di \param{pathname}.
  \item[\errcode{ENAMETOOLONG}] il nome indicato in \param{path} è troppo lungo.
  \item[\errcode{ENOTDIR}] non si è specificata una directory.
  \end{errlist}
  ed inoltre \errval{EFAULT}, \errval{EIO}, \errval{ELOOP}, \errval{ENOENT} e
  \errval{ENOMEM} nel loro significato generico.}
\end{funcproto}

La funzione cambia la directory di lavoro in \param{pathname} ed
ovviamente \param{pathname} deve indicare una directory per la quale si hanno
i permessi di accesso.

Dato che ci si può riferire ad una directory anche tramite un file descriptor,
per cambiare directory di lavoro è disponibile anche la funzione di sistema
\funcd{fchdir}, il cui prototipo è:

\begin{funcproto}{
\fhead{unistd.h}
\fdecl{int fchdir(int fd)}
\fdesc{Cambia la directory di lavoro per file descriptor.} 
}
{La funzione ritorna $0$ in caso di successo e $-1$ per un errore, nel qual
  caso \var{errno} assumerà i valori \errval{EBADF} o \errval{EACCES} nel loro
  significato generico.}
\end{funcproto}

La funzione è identica a \func{chdir}, ma prende come argomento un file
descriptor \param{fd} invece di un \textit{pathname}. Anche in questo
caso \param{fd} deve essere un file descriptor valido che fa riferimento ad
una directory. Inoltre l'unico errore di accesso possibile (tutti gli altri
sarebbero occorsi all'apertura di \param{fd}), è quello in cui il processo non
ha il permesso di attraversamento alla directory specificata da \param{fd}.

\index{directory~di~lavoro|)} 


\subsection{La creazione dei \textsl{file speciali}}
\label{sec:file_mknod}

\index{file!di~dispositivo|(} 
\index{file!speciali|(} 

Finora abbiamo parlato esclusivamente di file, directory e collegamenti
simbolici, ma in sez.~\ref{sec:file_file_types} abbiamo visto che il sistema
prevede anche degli altri tipi di file, che in genere vanno sotto il nome
generico di \textsl{file speciali}, come i file di dispositivo, le
\textit{fifo} ed i socket.

La manipolazione delle caratteristiche di questi file speciali, il cambiamento
di nome o la loro cancellazione può essere effettuata con le stesse funzioni
che operano sugli altri file, ma quando li si devono creare sono necessarie,
come per le directory, delle funzioni apposite. La prima di queste è la
funzione di sistema \funcd{mknod}, il cui prototipo è:

\begin{funcproto}{
\fhead{sys/types.h}
\fhead{sys/stat.h}
\fhead{fcntl.h}
\fhead{unistd.h}
\fdecl{int mknod(const char *pathname, mode\_t mode, dev\_t dev)}
\fdesc{Crea un file speciale sul filesystem.} 
}
{La funzione ritorna $0$ in caso di successo e $-1$ per un errore, nel qual
  caso \var{errno} assumerà uno dei valori: 
  \begin{errlist}
  \item[\errcode{EEXIST}] \param{pathname} esiste già o è un collegamento
    simbolico. 
  \item[\errcode{EINVAL}] il valore di \param{mode} non indica un file, una
    \textit{fifo}, un socket o un dispositivo.
  \item[\errcode{EPERM}] non si hanno privilegi sufficienti a creare
    l'\texttt{inode}, o il filesystem su cui si è cercato di
    creare \param{pathname} non supporta l'operazione.
  \end{errlist}
  ed inoltre \errval{EACCES}, \errval{EFAULT}, \errval{ELOOP},
  \errval{ENAMETOOLONG}, \errval{ENOENT}, \errval{ENOMEM}, \errval{ENOSPC},
  \errval{ENOTDIR} e \errval{EROFS} nel loro significato generico.}
\end{funcproto}

La funzione permette di creare un \textit{inode} di tipo generico sul
filesystem, e viene in genere utilizzata per creare i file di dispositivo, ma
si può usare anche per creare qualunque tipo di file speciale ed anche file
regolari. L'argomento \param{mode} specifica sia il tipo di file che si vuole
creare che i relativi permessi, secondo i valori riportati in
tab.~\ref{tab:file_mode_flags}, che vanno combinati con un OR aritmetico. I
permessi sono comunque modificati nella maniera usuale dal valore di
\textit{umask} (si veda sez.~\ref{sec:file_perm_management}).

Per il tipo di file può essere specificato solo uno fra i seguenti valori:
\const{S\_IFREG} per un file regolare (che sarà creato vuoto),
\const{S\_IFBLK} per un dispositivo a blocchi, \const{S\_IFCHR} per un
dispositivo a caratteri, \const{S\_IFSOCK} per un socket e \const{S\_IFIFO}
per una \textit{fifo};\footnote{con Linux la funzione non può essere usata per
  creare directory o collegamenti simbolici, si dovranno usare le funzioni
  \func{mkdir} e \func{symlink} a questo dedicate.} un valore diverso
comporterà l'errore \errcode{EINVAL}. Inoltre \param{pathname} non deve
esistere, neanche come collegamento simbolico.

Qualora si sia specificato in \param{mode} un file di dispositivo (vale a dire
o \const{S\_IFBLK} o \const{S\_IFCHR}), il valore di \param{dev} dovrà essere
usato per indicare a quale dispositivo si fa riferimento, altrimenti il suo
valore verrà ignorato.  Solo l'amministratore può creare un file di
dispositivo usando questa funzione (il processo deve avere la capacità
\const{CAP\_MKNOD}), ma in Linux\footnote{questo è un comportamento specifico
  di Linux, la funzione non è prevista dallo standard POSIX.1 originale,
  mentre è presente in SVr4 e 4.4BSD, ma esistono differenze nei comportamenti
  e nei codici di errore, tanto che questa è stata introdotta in POSIX.1-2001
  con una nota che la definisce portabile solo quando viene usata per creare
  delle \textit{fifo}, ma comunque deprecata essendo utilizzabile a tale scopo
  la specifica \func{mkfifo}.} l'uso per la creazione di un file ordinario, di
una \textit{fifo} o di un socket è consentito anche agli utenti normali.

Gli \textit{inode} creati con \func{mknod} apparterranno al proprietario e al
gruppo del processo (usando \ids{UID} e \ids{GID} del gruppo effettivo) che li
ha creati a meno non sia presente il bit \acr{sgid} per la directory o sia
stata attivata la semantica BSD per il filesystem (si veda
sez.~\ref{sec:file_ownership_management}) in cui si va a creare
l'\textit{inode}, nel qual caso per il gruppo verrà usato il \ids{GID} del
proprietario della directory.

\itindbeg{major~number}
\itindbeg{minor~number}

Nella creazione di un file di dispositivo occorre poi specificare
correttamente il valore di \param{dev}; questo infatti è di tipo
\type{dev\_t}, che è un tipo primitivo (vedi
tab.~\ref{tab:intro_primitive_types}) riservato per indicare un
\textsl{numero} di dispositivo. Il kernel infatti identifica ciascun
dispositivo con un valore numerico, originariamente questo era un intero a 16
bit diviso in due parti di 8 bit chiamate rispettivamente \textit{major
  number} e \textit{minor number}, che sono poi i due numeri mostrati dal
comando \texttt{ls -l} al posto della dimensione quando lo si esegue su un
file di dispositivo.

Il \textit{major number} identifica una classe di dispositivi (ad esempio la
seriale, o i dischi IDE) e serve in sostanza per indicare al kernel quale è il
modulo che gestisce quella classe di dispositivi. Per identificare uno
specifico dispositivo di quella classe (ad esempio una singola porta seriale,
o uno dei dischi presenti) si usa invece il \textit{minor number}. L'elenco
aggiornato di questi numeri con le relative corrispondenze ai vari dispositivi
può essere trovato nel file \texttt{Documentation/devices.txt} allegato alla
documentazione dei sorgenti del kernel.

Data la crescita nel numero di dispositivi supportati, ben presto il limite
massimo di 256 si è rivelato troppo basso, e nel passaggio dai kernel della
serie 2.4 alla serie 2.6 è stata aumentata a 32 bit la dimensione del tipo
\type{dev\_t}, con delle dimensioni passate a 12 bit per il \textit{major
  number} e 20 bit per il \textit{minor number}. La transizione però ha
comportato il fatto che \type{dev\_t} è diventato un tipo opaco, e la
necessità di specificare il numero tramite delle opportune macro, così da non
avere problemi di compatibilità con eventuali ulteriori estensioni.

Le macro sono definite nel file \headfiled{sys/sysmacros.h},\footnote{se si
  usa la \acr{glibc} dalla versione 2.3.3 queste macro sono degli alias alle
  versioni specifiche di questa libreria, \macrod{gnu\_dev\_major},
  \macrod{gnu\_dev\_minor} e \macrod{gnu\_dev\_makedev} che si possono usare
  direttamente, al costo di una minore portabilità.} che viene automaticamente
incluso quando si include \headfile{sys/types.h}. Si possono pertanto ottenere
i valori del \textit{major number} e \textit{minor number} di un dispositivo
rispettivamente con le macro \macro{major} e \macro{minor}:

{\centering
\vspace{3pt}
\begin{funcbox}{
\fhead{sys/types.h}
\fdecl{int \macrod{major}(dev\_t dev)}
\fdesc{Restituisce il \textit{major number} del dispositivo \param{dev}.}
\fdecl{int \macrod{minor}(dev\_t dev)}
\fdesc{Restituisce il \textit{minor number} del dispositivo \param{dev}.}  
} 
\end{funcbox}
}

\noindent mentre una volta che siano noti \textit{major number} e
\textit{minor number} si potrà costruire il relativo identificativo con la
macro \macro{makedev}:

{\centering
\vspace{3pt}
\begin{funcbox}{
\fhead{sys/types.h}
\fdecl{dev\_t \macrod{makedev}(int major, int minor)}
\fdesc{Dati \textit{major number} e \textit{minor number} restituisce
  l'identificativo di un dispositivo.} 
} 
\end{funcbox}
}


\itindend{major~number}
\itindend{minor~number}
\index{file!di~dispositivo|)}

Dato che la funzione di sistema \func{mknod} presenta diverse varianti nei
vari sistemi unix-like, lo standard POSIX.1-2001 la dichiara portabile solo in
caso di creazione delle \textit{fifo}, ma anche in questo caso alcune
combinazioni degli argomenti restano non specificate, per cui nello stesso
standard è stata introdotta una funzione specifica per creare una
\textit{fifo} deprecando l'uso di \func{mknod} a tale riguardo.  La funzione è
\funcd{mkfifo} ed il suo prototipo è:

\begin{funcproto}{
\fhead{sys/types.h}
\fhead{sys/stat.h}
\fdecl{int mkfifo(const char *pathname, mode\_t mode)}
\fdesc{Crea una \textit{fifo}.} 
}
{La funzione ritorna $0$ in caso di successo e $-1$ per un errore, nel qual
  caso \var{errno} assumerà \errval{EACCES}, \errval{EEXIST},
  \errval{ENAMETOOLONG}, \errval{ENOENT}, \errval{ENOSPC}, \errval{ENOTDIR} e
  \errval{EROFS} nel loro significato generico.}
\end{funcproto}

La funzione crea la \textit{fifo} \param{pathname} con i
permessi \param{mode}. Come per \func{mknod} il file \param{pathname} non deve
esistere (neanche come collegamento simbolico); al solito i permessi
specificati da \param{mode} vengono modificati dal valore di \textit{umask}
(vedi sez.~\ref{sec:file_perm_management}).

\index{file!speciali|)} 


\subsection{I file temporanei}
\label{sec:file_temp_file}

In molte occasioni è utile poter creare dei file temporanei; benché la cosa
sembri semplice, in realtà il problema è più sottile di quanto non appaia a
prima vista. Infatti anche se sembrerebbe banale generare un nome a caso e
creare il file dopo aver controllato che questo non esista, nel momento fra il
controllo e la creazione si ha giusto lo spazio per una possibile \textit{race
  condition} (si ricordi quanto visto in sez.~\ref{sec:proc_race_cond}).

\itindbeg{symlink~attack}

Molti problemi di sicurezza derivano proprio da una creazione non accorta di
file temporanei che lascia aperta questa \textit{race condition}. Un
attaccante allora potrà sfruttarla con quello che viene chiamato
``\textit{symlink attack}'' dove nell'intervallo fra la generazione di un nome
e l'accesso allo stesso, viene creato un collegamento simbolico con quel nome
verso un file diverso, ottenendo, se il programma sotto attacco ne ha la
capacità, un accesso privilegiato.\footnote{dal kernel 3.6 sono state
  introdotte delle contromisure, illustrate in
  sez.~\ref{sec:procadv_security_misc}, che rendono impraticabili questo tipo
  di attacchi, ma questa non è una buona scusa per ignorare il problema.}

\itindend{symlink~attack}

La \acr{glibc} provvede varie funzioni per generare nomi di file temporanei,
di cui si abbia certezza di unicità al momento della generazione; storicamente
la prima di queste funzioni create a questo scopo era
\funcd{tmpnam},\footnote{la funzione è stata deprecata nella revisione
  POSIX.1-2008 dello standard POSIX.} il cui prototipo è:

\begin{funcproto}{
\fhead{stdio.h}
\fdecl{char *tmpnam(char *string)}
\fdesc{Genera un nome univoco per un file temporaneo.} 
}
{La funzione ritorna il puntatore alla stringa con il nome in caso di successo
  e \val{NULL} in caso di fallimento, non sono definiti errori.}
\end{funcproto}

La funzione restituisce il puntatore ad una stringa contente un nome di file
valido e non esistente al momento dell'invocazione. Se si è passato come
argomento \param{string} un puntatore non nullo ad un buffer di caratteri
questo deve essere di dimensione \constd{L\_tmpnam} ed il nome generato vi
verrà copiato automaticamente, altrimenti il nome sarà generato in un buffer
statico interno che verrà sovrascritto ad una chiamata successiva.  Successive
invocazioni della funzione continueranno a restituire nomi unici fino ad un
massimo di \constd{TMP\_MAX} volte, limite oltre il quale il comportamento è
indefinito. Al nome viene automaticamente aggiunto come prefisso la directory
specificata dalla costante \constd{P\_tmpdir}.\footnote{le costanti
  \const{L\_tmpnam}, \const{P\_tmpdir} e \const{TMP\_MAX} sono definite in
  \headfile{stdio.h}.}

Di questa funzione esiste una versione rientrante, \funcm{tmpnam\_r}, che non
fa nulla quando si passa \val{NULL} come argomento.  Una funzione simile,
\funcd{tempnam}, permette di specificare un prefisso per il file
esplicitamente, il suo prototipo è:

\begin{funcproto}{
\fhead{stdio.h}
\fdecl{char *tempnam(const char *dir, const char *pfx)}
\fdesc{Genera un nome univoco per un file temporaneo.} 
}
{La funzione ritorna il puntatore alla stringa con il nome in caso di successo
  e \val{NULL} per un errore, nel qual caso \var{errno} potrà assumere solo il
  valore \errval{ENOMEM} qualora fallisca l'allocazione della stringa.}
\end{funcproto}

La funzione alloca con \code{malloc} la stringa in cui restituisce il nome,
per cui è sempre rientrante, occorre però ricordarsi di disallocare con
\code{free} il puntatore che restituisce.  L'argomento \param{pfx} specifica
un prefisso di massimo 5 caratteri per il nome provvisorio. La funzione
assegna come directory per il file temporaneo, verificando che esista e sia
accessibile, la prima valida fra le seguenti:
\begin{itemize*}
\item la variabile di ambiente \envvar{TMPDIR} (non ha effetto se non è
  definita o se il programma chiamante è \acr{suid} o \acr{sgid}, vedi
  sez.~\ref{sec:file_special_perm}),
\item il valore dell'argomento \param{dir} (se diverso da \val{NULL}),
\item il valore della costante \const{P\_tmpdir},
\item la directory \file{/tmp}.
\end{itemize*}

In ogni caso, anche se con queste funzioni la generazione del nome è casuale,
ed è molto difficile ottenere un nome duplicato, nulla assicura che un altro
processo non possa avere creato, fra l'ottenimento del nome e l'apertura del
file, un altro file o un collegamento simbolico con lo stesso nome. Per questo
motivo quando si usa il nome ottenuto da una di queste funzioni occorre sempre
assicurarsi che non si stia usando un collegamento simbolico e aprire il nuovo
file in modalità di esclusione (cioè con l'opzione \const{O\_EXCL} per i file
descriptor o con il flag ``\texttt{x}'' per gli \textit{stream}) che fa
fallire l'apertura in caso il file sia già esistente. Essendo disponibili
alternative migliori l'uso di queste funzioni è deprecato.

Per evitare di dovere effettuare a mano tutti questi controlli, lo standard
POSIX definisce la funzione \funcd{tmpfile}, che permette di ottenere in
maniera sicura l'accesso ad un file temporaneo, il suo prototipo è:

\begin{funcproto}{
\fhead{stdio.h}
\fdecl{FILE *tmpfile(void)}
\fdesc{Apre un file temporaneo in lettura/scrittura.} 
}
{La funzione ritorna il puntatore allo \textit{stream} associato al file
  temporaneo in caso di successo e \val{NULL} per un errore, nel qual caso
  \var{errno} assumerà uno dei valori:
  \begin{errlist}
    \item[\errcode{EEXIST}] non è stato possibile generare un nome univoco.
    \item[\errcode{EINTR}] la funzione è stata interrotta da un segnale.
  \end{errlist}
  ed inoltre \errval{EFAULT}, \errval{EMFILE}, \errval{ENFILE},
  \errval{ENOSPC}, \errval{EROFS} e \errval{EACCES} nel loro significato
  generico.}
\end{funcproto}


La funzione restituisce direttamente uno \textit{stream} già aperto (in
modalità \code{w+b}, si veda sez.~\ref{sec:file_fopen}) e pronto per l'uso,
che viene automaticamente cancellato alla sua chiusura o all'uscita dal
programma. Lo standard non specifica in quale directory verrà aperto il file,
ma la \acr{glibc} prima tenta con \const{P\_tmpdir} e poi con
\file{/tmp}. Questa funzione è rientrante e non soffre di problemi di
\textit{race condition}.

Alcune versioni meno recenti di Unix non supportano queste funzioni; in questo
caso si possono usare le vecchie funzioni \funcd{mktemp} e \func{mkstemp} che
modificano una stringa di input che serve da modello e che deve essere
conclusa da 6 caratteri ``\texttt{X}'' che verranno sostituiti da un codice
unico. La prima delle due è analoga a \func{tmpnam} e genera soltanto un nome
casuale, il suo prototipo è:

\begin{funcproto}{
\fhead{stlib.h}
\fdecl{char *mktemp(char *template)}
\fdesc{Genera un nome univoco per un file temporaneo.} 
}
{La funzione ritorna  il puntatore a \param{template} in caso di successo e
  \val{NULL} per un errore, nel qual caso \var{errno} assumerà uno dei valori: 
  \begin{errlist}
    \item[\errcode{EINVAL}] \param{template} non termina con \code{XXXXXX}.
  \end{errlist}}
\end{funcproto}

La funzione genera un nome univoco sostituendo le \code{XXXXXX} finali di
\param{template}; dato che \param{template} deve poter essere modificata dalla
funzione non si può usare una stringa costante.  Tutte le avvertenze riguardo
alle possibili \textit{race condition} date per \func{tmpnam} continuano a
valere; inoltre in alcune vecchie implementazioni il valore usato per
sostituire le \code{XXXXXX} viene formato con il \ids{PID} del processo più
una lettera, il che mette a disposizione solo 26 possibilità diverse per il
nome del file, e rende il nome temporaneo facile da indovinare.  Per tutti
questi motivi la funzione è deprecata e non dovrebbe mai essere usata.

La seconda funzione, \funcd{mkstemp} è sostanzialmente equivalente a
\func{tmpfile}, ma restituisce un file descriptor invece di un nome; il suo
prototipo è:

\begin{funcproto}{
\fhead{stlib.h}
\fdecl{int mkstemp(char *template)}
\fdesc{Apre un file temporaneo.} 
}

{La funzione ritorna il file descriptor in caso di successo e $-1$ per un
  errore, nel qual 
  caso \var{errno} assumerà uno dei valori: 
  \begin{errlist}
    \item[\errcode{EEXIST}] non è riuscita a creare un file temporaneo, il
      contenuto di \param{template} è indefinito.
    \item[\errcode{EINVAL}] \param{template} non termina con \code{XXXXXX}.
  \end{errlist}}
\end{funcproto}

Come per \func{mktemp} anche in questo caso \param{template} non può essere
una stringa costante. La funzione apre un file in lettura/scrittura con la
funzione \func{open}, usando l'opzione \const{O\_EXCL} (si veda
sez.~\ref{sec:file_open_close}), in questo modo al ritorno della funzione si
ha la certezza di essere stati i creatori del file, i cui permessi (si veda
sez.~\ref{sec:file_perm_overview}) sono impostati al valore \code{0600}
(lettura e scrittura solo per il proprietario).\footnote{questo è vero a
  partire dalla \acr{glibc} 2.0.7, le versioni precedenti della \acr{glibc} e
  le vecchie \acr{libc5} e \acr{libc4} usavano il valore \code{0666} che
  permetteva a chiunque di leggere e scrivere i contenuti del file.}  Di
questa funzione esiste una variante \funcd{mkostemp}, introdotta
specificamente dalla \acr{glibc},\footnote{la funzione è stata introdotta
  nella versione 2.7 delle librerie e richiede che sia definita la macro
  \macro{\_GNU\_SOURCE}.} il cui prototipo è:

\begin{funcproto}{
\fhead{stlib.h}
\fdecl{int mkostemp(char *template, int flags)}
\fdesc{Apre un file temporaneo.} 
}
{La funzione ritorna un file descriptor in caso di successo e $-1$ per un
  errore, nel qual caso \var{errno} assumerà  gli stessi valori di
  \func{mkstemp}.} 
\end{funcproto}
\noindent la cui sola differenza è la presenza dell'ulteriore argomento
\var{flags} che consente di specificare alcuni ulteriori flag (come
\const{O\_APPEND}, \const{O\_CLOEXEC}, \const{O\_SYNC}, il cui significato
vedremo in sez.~\ref{sec:file_open_close}) da passare ad \func{open}
nell'apertura del file.\footnote{si tenga presente che \func{mkostemp}
  utilizza già \const{O\_CREAT}, \const{O\_EXCL} e \const{O\_RDWR}, che non è
  il caso di riindicare, dato che ciò potrebbe portare ad errori in altri
  sistemi operativi.}

Di queste due funzioni sono state poi introdotte, a partire dalla \acr{glibc}
2.11 due varianti, \funcd{mkstemps} e \funcd{mkostemps}, che consentono di
indicare anche un suffisso, i loro prototipi sono:

\begin{funcproto}{
\fhead{stlib.h}
\fdecl{int mkstemps(char *template, int suffixlen)}
\fdesc{Apre un file temporaneo.} 
\fdecl{int mkostemps(char *template, int suffixlen, int flags)}
\fdesc{Apre un file temporaneo.} 
}

{Le funzioni hanno gli stessi valori di ritorno e gli stessi errori di
  \func{mkstemp} con lo stesso significato, tranne \errval{EINVAL} che viene
  restituito se \param{template} non è di lunghezza pari ad almeno
  $6+$\param{suffixlen} ed i 6 caratteri prima del suffisso non sono
  \code{XXXXXX}.}
\end{funcproto}

Le due funzioni, un'estensione non standard fornita dalla \acr{glibc}, sono
identiche a \funcd{mkstemp} e \funcd{mkostemp}, ma consentono di avere un nome
del file nella forma \texttt{prefissoXXXXXXsuffisso} dove la lunghezza del
suffisso deve essere indicata con \param{suffixlen}.

Infine con OpenBSD è stata introdotta un'altra funzione simile alle
precedenti, \funcd{mkdtemp}, che crea invece una directory
temporanea;\footnote{la funzione è stata introdotta nella \acr{glibc} a
  partire dalla versione 2.1.91 ed inserita nello standard POSIX.1-2008.}  il
suo prototipo è:

\begin{funcproto}{
\fhead{stlib.h}
\fdecl{char *mkdtemp(char *template)}
\fdesc{Crea una directory temporanea.} 
}
{La funzione ritorna il puntatore al nome della directory in caso di successo
  e \val{NULL} per un errore, nel qual caso \var{errno} assumerà uno dei
  valori:
  \begin{errlist}
    \item[\errcode{EINVAL}] \param{template} non termina con \code{XXXXXX}.
  \end{errlist}
  più gli altri eventuali codici di errore di \func{mkdir}.}
\end{funcproto}

La funzione crea una directory temporanea il cui nome è ottenuto sostituendo
le \code{XXXXXX} finali di \param{template} con permessi \code{0700} (si veda
sez.~\ref{sec:file_perm_overview} per i dettagli). Dato che la creazione della
directory è sempre atomica i precedenti problemi di \textit{race condition}
non si pongono.


\section{La manipolazione delle caratteristiche dei file}
\label{sec:file_infos}

Come spiegato in sez.~\ref{sec:file_filesystem} tutte le informazioni generali
relative alle caratteristiche di ciascun file, a partire dalle informazioni
relative al controllo di accesso, sono mantenute nell'\textit{inode}. Vedremo
in questa sezione come sia possibile leggere tutte queste informazioni usando
la funzione \func{stat}, che permette l'accesso a tutti i dati memorizzati
nell'\textit{inode}; esamineremo poi le varie funzioni usate per manipolare
tutte queste informazioni, eccetto quelle che riguardano la gestione del
controllo di accesso, trattate in sez.~\ref{sec:file_access_control}.


\subsection{La lettura delle caratteristiche dei file}
\label{sec:file_stat}

La lettura delle informazioni relative ai file è fatta attraverso la famiglia
delle funzioni \func{stat} che sono quelle che usa il comando \cmd{ls} per
poter ottenere e mostrare tutti i dati relativi ad un file; ne fanno parte le
funzioni di sistema \funcd{stat}, \funcd{fstat} e \funcd{lstat}, i cui
prototipi sono:

\begin{funcproto}{
\fhead{sys/types.h}
\fhead{sys/stat.h}
\fhead{unistd.h}
\fdecl{int stat(const char *file\_name, struct stat *buf)}
\fdecl{int lstat(const char *file\_name, struct stat *buf)}
\fdecl{int fstat(int filedes, struct stat *buf)}
\fdesc{Leggono le informazioni di un file.} 
}
{Le funzioni ritornano $0$ in caso di successo e $-1$ per un errore, nel qual
  caso \var{errno} assumerà uno dei valori:
  \begin{errlist}
    \item[\errcode{EOVERFLOW}] il file ha una dimensione che non può essere
      rappresentata nel tipo \type{off\_t} (può avvenire solo in caso di
      programmi compilati su piattaforme a 32 bit senza le estensioni
      (\texttt{-D \_FILE\_OFFSET\_BITS=64}) per file a 64 bit).
  \end{errlist}
  ed inoltre \errval{EFAULT} ed \errval{ENOMEM}, per \func{stat} e
  \func{lstat} anche \errval{EACCES}, \errval{ELOOP}, \errval{ENAMETOOLONG},
  \errval{ENOENT}, \errval{ENOTDIR}, per \func{fstat} anche \errval{EBADF}, 
  nel loro significato generico.}
\end{funcproto}

La funzione \func{stat} legge le informazioni del file indicato
da \param{file\_name} e le inserisce nel buffer puntato
dall'argomento \param{buf}; la funzione \func{lstat} è identica a \func{stat}
eccetto che se \param{file\_name} è un collegamento simbolico vengono lette le
informazioni relative ad esso e non al file a cui fa riferimento. Infine
\func{fstat} esegue la stessa operazione su un file già aperto, specificato
tramite il suo file descriptor \param{filedes}.

La struttura \struct{stat} usata da queste funzioni è definita nell'header
\headfiled{sys/stat.h} e in generale dipende dall'implementazione; la versione
usata da Linux è mostrata in fig.~\ref{fig:file_stat_struct}, così come
riportata dalla pagina di manuale di \func{stat}. In realtà la definizione
effettivamente usata nel kernel dipende dall'architettura e ha altri campi
riservati per estensioni come tempi dei file più precisi (vedi
sez.~\ref{sec:file_file_times}).

\begin{figure}[!htb]
  \footnotesize
  \centering
  \begin{minipage}[c]{0.8\textwidth}
    \includestruct{listati/stat.h}
  \end{minipage} 
  \normalsize 
  \caption{La struttura \structd{stat} per la lettura delle informazioni dei 
    file.}
  \label{fig:file_stat_struct}
\end{figure}

Si noti come i vari membri della struttura siano specificati come tipi
primitivi del sistema, di quelli definiti in
tab.~\ref{tab:intro_primitive_types}, e dichiarati in \headfile{sys/types.h},
con l'eccezione di \typed{blksize\_t} e \typed{blkcnt\_t} che sono nuovi tipi
introdotti per rendersi indipendenti dalla piattaforma. 

Benché la descrizione dei commenti di fig.~\ref{fig:file_stat_struct} sia
abbastanza chiara, vale la pena illustrare maggiormente il significato dei
campi di \struct{stat} su cui non torneremo in maggior dettaglio nel resto di
questa sezione:
\begin{itemize*}
\item Il campo \var{st\_nlink} contiene il numero di \textit{hard link} che
  fanno riferimento al file (il cosiddetto \textit{link count}) di cui abbiamo
  già parlato in numerose occasioni.
\item Il campo \var{st\_ino} contiene il numero di \textit{inode} del file,
  quello viene usato all'interno del filesystem per identificarlo e che può
  essere usato da un programma per determinare se due \textit{pathname} fanno
  riferimento allo stesso file.
\item Il campo \var{st\_dev} contiene il numero del dispositivo su cui risiede
  il file (o meglio il suo filesystem). Si tratta dello stesso numero che si
  usa con \func{mknod} e che può essere decomposto in \textit{major number} e
  \textit{minor number} con le macro \macro{major} e \macro{minor} viste in
  sez.~\ref{sec:file_mknod}.
\item Il campo \var{st\_rdev} contiene il numero di dispositivo associato al
  file stesso ed ovviamente ha un valore significativo soltanto quando il file
  è un dispositivo a caratteri o a blocchi.
\item Il campo \var{st\_blksize} contiene la dimensione dei blocchi di dati
  usati nell'I/O su disco, che è anche la dimensione usata per la
  bufferizzazione dei dati dalle librerie del C per l'interfaccia degli
  \textit{stream}.  Leggere o scrivere blocchi di dati in dimensioni inferiori
  a questo valore è inefficiente in quanto le operazioni su disco usano
  comunque trasferimenti di questa dimensione.
\end{itemize*}

Nell'evoluzione del kernel la \textit{system call} che fornisce \func{stat} è
stata modificata più volte per tener conto dei cambiamenti fatti alla
struttura \struct{stat},\footnote{questo ha significato l'utilizzo a basso
  livello di diverse \textit{system call} e diverse versioni della struttura.}
in particolare a riguardo ai tempi dei file, di cui è stata aumentata la
precisione (torneremo su questo in sez.~\ref{sec:file_file_times}) ma anche
per gli aggiornamenti fatti ai campi \var{st\_ino}, \var{st\_uid} e
\var{st\_gid}.

Sulle piattaforme a 32 bit questi cambiamenti, che han visto un aumento delle
dimensioni dei campi della struttura per adattarli alle nuove esigenze, sono
mascherati dalla \acr{glibc} che attraverso \func{stat} invoca la versione più
recente della \textit{system call} e rimpacchetta i dati se questo è
necessario per eseguire dei vecchi programmi. Nelle piattaforme a 64 bit
invece è presente un'unica versione della \textit{system call} e la struttura
\struct{stat} ha campi di dimensione sufficiente.

Infine a partire dal kernel 2.6.16 è stata introdotta una ulteriore funzione
della famiglia, \func{fstatat} che consente di trattare con sicurezza i
\textit{pathname} relativi, la tratteremo in sez.~\ref{sec:file_openat},
insieme alla nuova \textit{system call} \func{statx}, introdotta dal kernel
4.11 per estendere l'interfaccia di \func{stat} e le informazioni che essa può
restituire.


\subsection{I tipi di file}
\label{sec:file_types}

Abbiamo sottolineato fin dall'introduzione che Linux, come ogni sistema
unix-like, supporta oltre ai file ordinari e alle directory una serie di altri
``\textsl{tipi}'' di file che possono stare su un filesystem (elencati in
tab.~\ref{tab:file_file_types}).  Il tipo di file viene ritornato dalle
funzioni della famiglia \func{stat} all'interno del campo \var{st\_mode} di
una struttura \struct{stat}. 

\begin{table}[htb]
  \centering
  \footnotesize
  \begin{tabular}[c]{|l|l|}
    \hline
    \textbf{Macro} & \textbf{Tipo del file} \\
    \hline
    \hline
    \macrod{S\_ISREG}\texttt{(m)}  & File normale.\\
    \macrod{S\_ISDIR}\texttt{(m)}  & Directory.\\
    \macrod{S\_ISCHR}\texttt{(m)}  & Dispositivo a caratteri.\\
    \macrod{S\_ISBLK}\texttt{(m)}  & Dispositivo a blocchi.\\
    \macrod{S\_ISFIFO}\texttt{(m)} & \textit{Fifo}.\\
    \macrod{S\_ISLNK}\texttt{(m)}  & Collegamento simbolico.\\
    \macrod{S\_ISSOCK}\texttt{(m)} & Socket.\\
    \hline    
  \end{tabular}
  \caption{Macro per i tipi di file (definite in \headfile{sys/stat.h}).}
  \label{tab:file_type_macro}
\end{table}

Il campo \var{st\_mode} è una maschera binaria in cui l'informazione viene
suddivisa nei vari bit che compongono, ed oltre a quelle sul tipo di file,
contiene anche le informazioni relative ai permessi su cui torneremo in
sez.~\ref{sec:file_perm_overview}. Dato che i valori numerici usati per
definire il tipo di file possono variare a seconda delle implementazioni, lo
standard POSIX definisce un insieme di macro che consentono di verificare il
tipo di file in maniera standardizzata.

Queste macro vengono usate anche da Linux che supporta pure le estensioni allo
standard per i collegamenti simbolici e i socket definite da BSD.\footnote{le
  ultime due macro di tab.~\ref{tab:file_type_macro}, che non sono presenti
  nello standard POSIX fino alla revisione POSIX.1-1996.}  L'elenco completo
delle macro con cui è possibile estrarre da \var{st\_mode} l'informazione
relativa al tipo di file è riportato in tab.~\ref{tab:file_type_macro}, tutte
le macro restituiscono un valore intero da usare come valore logico e prendono
come argomento il valore di \var{st\_mode}.

\begin{table}[htb]
  \centering
  \footnotesize
  \begin{tabular}[c]{|l|c|l|}
    \hline
    \textbf{Flag} & \textbf{Valore} & \textbf{Significato} \\
    \hline
    \hline
    \constd{S\_IFMT}   &  0170000 & Maschera per i bit del tipo di file.\\
    \constd{S\_IFSOCK} &  0140000 & Socket.\\
    \constd{S\_IFLNK}  &  0120000 & Collegamento simbolico.\\
    \constd{S\_IFREG}  &  0100000 & File regolare.\\ 
    \constd{S\_IFBLK}  &  0060000 & Dispositivo a blocchi.\\
    \constd{S\_IFDIR}  &  0040000 & Directory.\\
    \constd{S\_IFCHR}  &  0020000 & Dispositivo a caratteri.\\
    \constd{S\_IFIFO}  &  0010000 & \textit{Fifo}.\\
    \hline
    \constd{S\_ISUID}  &  0004000 & Set user ID (\acr{suid}) bit, vedi
                                   sez.~\ref{sec:file_special_perm}).\\
    \constd{S\_ISGID}  &  0002000 & Set group ID (\acr{sgid}) bit, vedi
                                   sez.~\ref{sec:file_special_perm}).\\
    \constd{S\_ISVTX}  &  0001000 & \acr{Sticky} bit, vedi
                                   sez.~\ref{sec:file_special_perm}).\\
    \hline
    \constd{S\_IRWXU}  &  00700   & Maschera per i permessi del proprietario.\\
    \constd{S\_IRUSR}  &  00400   & Il proprietario ha permesso di lettura.\\
    \constd{S\_IWUSR}  &  00200   & Il proprietario ha permesso di scrittura.\\
    \constd{S\_IXUSR}  &  00100   & Il proprietario ha permesso di esecuzione.\\
    \hline
    \constd{S\_IRWXG}  &  00070   & Maschera per i permessi del gruppo.\\
    \constd{S\_IRGRP}  &  00040   & Il gruppo ha permesso di lettura.\\
    \constd{S\_IWGRP}  &  00020   & Il gruppo ha permesso di scrittura.\\
    \constd{S\_IXGRP}  &  00010   & Il gruppo ha permesso di esecuzione.\\
    \hline
    \constd{S\_IRWXO}  &  00007   & Maschera per i permessi di tutti gli altri\\
    \constd{S\_IROTH}  &  00004   & Gli altri hanno permesso di lettura.\\
    \constd{S\_IWOTH}  &  00002   & Gli altri hanno permesso di esecuzione.\\
    \constd{S\_IXOTH}  &  00001   & Gli altri hanno permesso di esecuzione.\\
    \hline    
  \end{tabular}
  \caption{Costanti per l'identificazione dei vari bit che compongono il campo
    \var{st\_mode} (definite in \headfile{sys/stat.h}).}
  \label{tab:file_mode_flags}
\end{table}

Oltre alle macro di tab.~\ref{tab:file_type_macro}, che semplificano
l'operazione, è possibile usare direttamente il valore di \var{st\_mode} per
ricavare il tipo di file controllando direttamente i vari bit in esso
memorizzati. Per questo sempre in \headfile{sys/stat.h} sono definite le varie
costanti numeriche riportate in tab.~\ref{tab:file_mode_flags}, che
definiscono le maschere che consentono di selezionare non solo i dati relativi
al tipo di file, ma anche le informazioni relative ai permessi su cui
torneremo in sez.~\ref{sec:file_access_control}, ed identificare i rispettivi
valori.

Le costanti che servono per la identificazione del tipo di file sono riportate
nella prima sezione di tab.~\ref{tab:file_mode_flags}, mentre le sezioni
successive attengono alle costanti usate per i permessi.  Il primo valore
dell'elenco è la maschera binaria \const{S\_IFMT} che permette di estrarre da
\var{st\_mode} (con un AND aritmetico) il blocco di bit nei quali viene
memorizzato il tipo di file. I valori successivi sono le costanti
corrispondenti ai vari tipi di file, e possono essere usate per verificare la
presenza del tipo di file voluto ed anche, con opportune combinazioni,
alternative fra più tipi di file. 

Si tenga presente però che a differenza dei permessi, l'informazione relativa
al tipo di file non è una maschera binaria, per questo motivo se si volesse
impostare una condizione che permetta di controllare se un file è una
directory o un file ordinario non si possono controllare dei singoli bit, ma
si dovrebbe usare una macro di preprocessore come:
\includecodesnip{listati/is_regdir.h}
in cui si estraggono ogni volta da \var{st\_mode} i bit relativi al tipo di
file e poi si effettua il confronto con i due possibili tipi di file.


\subsection{Le dimensioni dei file}
\label{sec:file_file_size}

Abbiamo visto in fig.~\ref{fig:file_stat_struct} che campo \var{st\_size} di
una struttura \struct{stat} contiene la dimensione del file in byte. In realtà
questo è vero solo se si tratta di un file regolare contenente dei dati; nel
caso di un collegamento simbolico invece la dimensione è quella del
\textit{pathname} che il collegamento stesso contiene, e per una directory
quella dello spazio occupato per le voci della stessa (che dipende da come
queste vengono mantenute dal filesystem), infine per le \textit{fifo}, i socket
ed i file di dispositivo questo campo è sempre nullo.

Il campo \var{st\_blocks} invece definisce la lunghezza del file espressa in
numero di blocchi di 512 byte. La differenza con \var{st\_size} è che in
questo caso si fa riferimento alla quantità di spazio disco allocata per il
file, e non alla dimensione dello stesso che si otterrebbe leggendolo
sequenzialmente.

Si deve tener presente infatti che la lunghezza del file riportata in
\var{st\_size} non è detto che corrisponda all'occupazione dello spazio su
disco, e non solo perché la parte finale del file potrebbe riempire
parzialmente un blocco. In un sistema unix-like infatti è possibile
l'esistenza dei cosiddetti \textit{sparse file}, cioè file in cui sono
presenti dei ``\textsl{buchi}'' (\textit{holes} nella nomenclatura inglese)
che si formano tutte le volte che si va a scrivere su un file dopo aver
eseguito uno spostamento oltre la sua fine (tratteremo in dettaglio
l'argomento in sez.~\ref{sec:file_lseek}).

In questo caso si avranno risultati differenti a seconda del modo in cui si
calcola la lunghezza del file, ad esempio il comando \cmd{du}, (che riporta il
numero di blocchi occupati) potrà dare una dimensione inferiore, mentre se si
legge il contenuto del file (ad esempio usando il comando \cmd{wc -c}), dato
che in tal caso per i ``\textsl{buchi}'' vengono restituiti degli zeri, si
avrà lo stesso risultato di \cmd{ls}.

Se è sempre possibile allargare un file, scrivendoci sopra o usando la
funzione \func{lseek} (vedi sez.~\ref{sec:file_lseek}) per spostarsi oltre la
sua fine, esistono anche casi in cui si può avere bisogno di effettuare un
troncamento, scartando i dati presenti al di là della dimensione scelta come
nuova fine del file.

Un file può sempre essere troncato a zero aprendolo con il flag
\const{O\_TRUNC} (vedi sez.~\ref{sec:file_open_close}), ma questo è un caso
particolare; per qualunque altra dimensione si possono usare le due funzioni
di sistema \funcd{truncate} e \funcd{ftruncate}, i cui prototipi sono:

\begin{funcproto}{
\fhead{unistd.h}
\fdecl{int ftruncate(int fd, off\_t length))}
\fdecl{int truncate(const char *file\_name, off\_t length)}
\fdesc{Troncano un file.} 
}
{Le funzioni ritornano $0$ in caso di successo e $-1$ per un errore, nel qual
  caso \var{errno} assumerà uno dei valori: 
  \begin{errlist}
  \item[\errcode{EINTR}] si è stati interrotti da un segnale.
  \item[\errcode{EINVAL}] \param{length} è negativa o maggiore delle
    dimensioni massime di un file.
  \item[\errcode{EPERM}] il filesystem non supporta il troncamento di un file.
  \item[\errcode{ETXTBSY}] il file è un programma in esecuzione.
  \end{errlist} 
  per entrambe, mentre per \func{ftruncate} si avranno anche: 
  \begin{errlist}
  \item[\errcode{EBADF}] \param{fd} non è un file descriptor.
  \item[\errcode{EINVAL}] \param{fd} non è un riferimento a un file o non è
    aperto in scrittura. 
  \end{errlist}
  e per \func{truncate} si avranno anche: 
  \begin{errlist}
  \item[\errcode{EACCES}] non si ha il permesso di scrittura sul file o il
    permesso di attraversamento di una delle componenti del \textit{pathname}.
  \item[\errcode{EISDIR}] \param{file\_name} fa riferimento ad una directory.
  \end{errlist}
  ed inoltre \errval{EFAULT}, \errval{EIO}, \errval{ELOOP},
  \errval{ENAMETOOLONG}, \errval{ENOENT}, \errval{ENOTDIR} e \errval{EROFS}
  nel loro significato generico.}
\end{funcproto}

Entrambe le funzioni fan sì che la dimensione del file sia troncata ad un
valore massimo specificato da \param{length}, e si distinguono solo per il
fatto che il file viene indicato con un \textit{pathname} per \func{truncate}
e con un file descriptor per \funcd{ftruncate}. Si tenga presente che se il
file è più lungo della lunghezza specificata i dati in eccesso saranno
perduti.

Il comportamento in caso di lunghezza del file inferiore a \param{length} non
è specificato e dipende dall'implementazione: il file può essere lasciato
invariato o esteso fino alla lunghezza scelta. Nel caso di Linux viene esteso
con la creazione di un \textsl{buco} nel file e ad una lettura si otterranno
degli zeri, si tenga presente però che questo comportamento è supportato solo
per filesystem nativi, ad esempio su un filesystem non nativo come il VFAT di
Windows questo non è possibile.


\subsection{I tempi dei file}
\label{sec:file_file_times}

Il sistema mantiene per ciascun file tre tempi, che sono registrati
nell'\textit{inode} insieme agli altri attributi del file. Questi possono
essere letti tramite la funzione \func{stat}, che li restituisce attraverso
tre campi della struttura \struct{stat} di fig.~\ref{fig:file_stat_struct}. Il
significato di questi tempi e dei relativi campi della struttura \struct{stat}
è illustrato nello schema di tab.~\ref{tab:file_file_times}, dove è anche
riportato un esempio delle funzioni che effettuano cambiamenti su di essi. Il
valore del tempo è espresso nel cosiddetto \textit{calendar time}, su cui
torneremo in dettaglio in sez.~\ref{sec:sys_time}.

\begin{table}[htb]
  \centering
  \footnotesize
  \begin{tabular}[c]{|c|l|l|c|}
    \hline
    \textbf{Membro} & \textbf{Significato} & \textbf{Funzione} 
    & \textbf{Opzione di \cmd{ls}} \\
    \hline
    \hline
    \var{st\_atime}& ultimo accesso ai dati del file    &
                     \func{read}, \func{utime}          & \cmd{-u}\\
    \var{st\_mtime}& ultima modifica ai dati del file   &
                     \func{write}, \func{utime}         & default\\
    \var{st\_ctime}& ultima modifica ai dati dell'\textit{inode} &
                     \func{chmod}, \func{utime}         & \cmd{-c}\\
    \hline
  \end{tabular}
  \caption{I tre tempi associati a ciascun file.}
  \label{tab:file_file_times}
\end{table}

Il primo punto da tenere presente è la differenza fra il cosiddetto tempo di
ultima modifica (il \textit{modification time}) riportato in \var{st\_mtime},
ed il tempo di ultimo cambiamento di stato (il \textit{change status time})
riportato in \var{st\_ctime}. Il primo infatti fa riferimento ad una modifica
del contenuto di un file, mentre il secondo ad una modifica dei metadati
dell'\textit{inode}. Dato che esistono molte operazioni, come la funzione
\func{link} e altre che vedremo in seguito, che modificano solo le
informazioni contenute nell'\textit{inode} senza toccare il contenuto del
file, diventa necessario l'utilizzo di questo secondo tempo.

Il tempo di ultima modifica viene usato ad esempio da programmi come
\cmd{make} per decidere quali file necessitano di essere ricompilati perché
più recenti dei loro sorgenti oppure dai programmi di backup, talvolta insieme
anche al tempo di cambiamento di stato, per decidere quali file devono essere
aggiornati nell'archiviazione.  Il tempo di ultimo accesso viene di solito
usato per identificare i file che non vengono più utilizzati per un certo
lasso di tempo. Ad esempio un programma come \texttt{leafnode} lo usa per
cancellare gli articoli letti più vecchi, mentre \texttt{mutt} lo usa per
marcare i messaggi di posta che risultano letti.  

Il sistema non tiene mai conto dell'ultimo accesso all'\textit{inode},
pertanto funzioni come \func{access} o \func{stat} non hanno alcuna influenza
sui tre tempi. Il comando \cmd{ls} (quando usato con le opzioni \cmd{-l} o
\cmd{-t}) mostra i tempi dei file secondo lo schema riportato nell'ultima
colonna di tab.~\ref{tab:file_file_times}. Si noti anche come in
tab.~\ref{tab:file_file_times} non venga riportato il \textsl{tempo di
  creazione} di un file; in un sistema unix-like infatti questo non esiste, e
non è previsto dall'interfaccia classica, ma essendo usato da altri sistemi
operativi (in particolare Windows) in tutti i filesystem più recenti ne viene
supportata la registrazione, ed a partire dal kernel 4.11 è diventato
possibile anche ottenerne la lettura con la nuova \textit{system call}
\func{statx} (che tratteremo in sez.~\ref{sec:file_openat}).

L'aggiornamento del tempo di ultimo accesso è stato a lungo considerato un
difetto progettuale di Unix, questo infatti comporta la necessità di
effettuare un accesso in scrittura sul disco anche in tutti i casi in cui
questa informazione non interessa e sarebbe possibile avere un semplice
accesso in lettura sui dati bufferizzati. Questo comporta un ovvio costo sia
in termini di prestazioni, che di consumo di risorse come la batteria per i
portatili, o i cicli di riscrittura per i dischi su memorie riscrivibili.

Per questo motivo abbiamo visto in sez.~\ref{sec:filesystem_mounting} come nel
corso dello sviluppo del kernel siano stati introdotti degli opportuni
\textit{mount flag} che consentono di evitare di aggiornare continuamente una
informazione che nella maggior parte dei casi non ha un interesse
rilevante. Per questo motivo i valori dell'\textit{access time} possono
dipendere dalle opzioni di montaggio, ed anche, essendo stato cambiato il
comportamento di default a partire dalla versione 2.6.30, dal kernel che si
sta usando.

In generale quello che avviene con i kernel più recenti è che il tempo di
ultimo accesso viene aggiornato solo se è precedente al tempo di ultima
modifica o cambiamento, o se è cambiato ed passato più di un giorno
dall'ultimo aggiornamento. Così si può rendere evidente che vi è stato un
accesso dopo una modifica, e che il file viene comunque osservato a cadenza
regolare, conservando le informazioni veramente utili senza consumare
inutilmente risorse in continue scritture per mantenere costantemente
aggiornata una informazione che a questo punto non ha più nessuna rilevanza
pratica.\footnote{qualora ce ne fosse la necessità è comunque possibile,
  tramite l'opzione di montaggio \texttt{strictatime}, richiedere in ogni caso
  il comportamento tradizionale.}

\begin{table}[htb]
  \centering
  \footnotesize
  \begin{tabular}[c]{|l|c|c|c|c|c|c|l|}
    \hline
    \multicolumn{1}{|p{2.3cm}|}{\centering{\vspace{6pt}\textbf{Funzione}}} &
    \multicolumn{3}{|p{3.2cm}|}{\centering{
        \textbf{File o directory del riferimento}}}&
    \multicolumn{3}{|p{3.2cm}|}{\centering{
        \textbf{Directory contenente il riferimento}}} 
    &\multicolumn{1}{|p{3.cm}|}{\centering{\vspace{6pt}\textbf{Note}}} \\
    \cline{2-7}
    \cline{2-7}
    \multicolumn{1}{|p{2.3cm}|}{} 
    &\multicolumn{1}{|p{.8cm}|}{\centering{\textsl{(a)}}}
    &\multicolumn{1}{|p{.8cm}|}{\centering{\textsl{(m)}}}
    &\multicolumn{1}{|p{.8cm}|}{\centering{\textsl{(c)}}}
    &\multicolumn{1}{|p{.8cm}|}{\centering{\textsl{(a)}}}
    &\multicolumn{1}{|p{.8cm}|}{\centering{\textsl{(m)}}}
    &\multicolumn{1}{|p{.8cm}|}{\centering{\textsl{(c)}}}
    &\multicolumn{1}{|p{3cm}|}{} \\
    \hline
    \hline
    \func{chmod}, \func{fchmod} 
             & --      & --      &$\bullet$& --      & --      & --      &\\
    \func{chown}, \func{fchown} 
             & --      & --      &$\bullet$& --      & --      & --      &\\
    \func{creat}  
             &$\bullet$&$\bullet$&$\bullet$& --      &$\bullet$&$\bullet$&  
             con \const{O\_CREATE} \\
    \func{creat}  
             & --      &$\bullet$&$\bullet$& --      &$\bullet$&$\bullet$&   
             con \const{O\_TRUNC} \\
    \func{exec}  
             &$\bullet$& --      & --      & --      & --      & --      &\\
    \func{lchown}  
             & --      & --      &$\bullet$& --      & --      & --      &\\
    \func{link}
             & --      & --      &$\bullet$& --      &$\bullet$&$\bullet$&\\
    \func{mkdir}
             &$\bullet$&$\bullet$&$\bullet$& --      &$\bullet$&$\bullet$&\\
    \func{mknod}
             &$\bullet$&$\bullet$&$\bullet$& --      &$\bullet$&$\bullet$&\\
    \func{mkfifo}
             &$\bullet$&$\bullet$&$\bullet$& --      &$\bullet$&$\bullet$&\\
    \func{open}
             &$\bullet$&$\bullet$&$\bullet$& --      &$\bullet$&$\bullet$& 
             con \const{O\_CREATE} \\
    \func{open}
             & --      &$\bullet$&$\bullet$& --      & --      & --      & 
             con \const{O\_TRUNC}  \\
    \func{pipe}
             &$\bullet$&$\bullet$&$\bullet$& --      & --      & --      &\\
    \func{read}
             &$\bullet$& --      & --      & --      & --      & --      &\\
    \func{remove}
             & --      & --      &$\bullet$& --      &$\bullet$&$\bullet$& 
             se esegue \func{unlink}\\
    \func{remove}
              & --      & --      & --      & --      &$\bullet$&$\bullet$& 
              se esegue \func{rmdir}\\
    \func{rename}
              & --      & --      &$\bullet$& --      &$\bullet$&$\bullet$& 
              per ambo gli argomenti\\
    \func{rmdir}
              & --      & --      & --      & --      &$\bullet$&$\bullet$&\\ 
    \func{truncate}
              & --      &$\bullet$&$\bullet$& --      & --      & --      &\\ 
    \func{ftruncate}
              & --      &$\bullet$&$\bullet$& --      & --      & --      &\\ 
    \func{unlink}
              & --      & --      &$\bullet$& --      &$\bullet$&$\bullet$&\\ 
    \func{utime}
              &$\bullet$&$\bullet$&$\bullet$& --      & --      & --      &\\ 
    \func{utimes}
              &$\bullet$&$\bullet$&$\bullet$& --      & --      & --      &\\ 
    \func{write}
              & --      &$\bullet$&$\bullet$& --      & --      & --      &\\ 
    \hline
  \end{tabular}
  \caption{Prospetto dei cambiamenti effettuati sui tempi di ultimo 
    accesso \textsl{(a)}, ultima modifica \textsl{(m)} e ultimo cambiamento di
    stato \textsl{(c)} dalle varie funzioni operanti su file e directory.}
  \label{tab:file_times_effects}  
\end{table}


L'effetto delle varie funzioni di manipolazione dei file sui relativi tempi è
illustrato in tab.~\ref{tab:file_times_effects}, facendo riferimento al
comportamento classico per quanto riguarda \var{st\_atime}. Si sono riportati
gli effetti sia per il file a cui si fa riferimento, sia per la directory che
lo contiene. Questi ultimi possono essere capiti immediatamente se si tiene
conto di quanto già detto e ripetuto a partire da
sez.~\ref{sec:file_filesystem}, e cioè che anche le directory sono anch'esse
file che contengono una lista di nomi, che il sistema tratta in maniera del
tutto analoga a tutti gli altri.

Per questo motivo tutte le volte che compiremo un'operazione su un file che
comporta una modifica del nome contenuto nella directory, andremo anche a
scrivere sulla directory che lo contiene cambiandone il tempo di ultima
modifica. Un esempio di questo tipo di operazione può essere la cancellazione
di un file, invece leggere o scrivere o cambiare i permessi di un file ha
effetti solo sui tempi di quest'ultimo.

Si ricordi infine come \var{st\_ctime} non è il tempo di creazione del file,
che in Unix non esiste, anche se può corrispondervi per file che non sono mai
stati modificati. Per questo motivo, a differenza di quanto avviene con altri
sistemi operativi, quando si copia un file, a meno di preservare
esplicitamente i tempi (ad esempio con l'opzione \cmd{-p} di \cmd{cp}) esso
avrà sempre il tempo corrente in cui si è effettuata la copia come data di
ultima modifica.

I tempi di ultimo accesso ed ultima modifica possono essere modificati
esplicitamente usando la funzione di sistema \funcd{utime}, il cui prototipo
è:

\begin{funcproto}{
\fhead{utime.h}
\fdecl{int utime(const char *filename, struct utimbuf *times)}
\fdesc{Modifica i tempi di ultimo accesso ed ultima modifica di un file.} 
}

{La funzione ritorna $0$ in caso di successo e $-1$ per un errore, nel qual
  caso \var{errno} assumerà uno dei valori: 
  \begin{errlist}
  \item[\errcode{EACCES}] non c'è il permesso di attraversamento per uno dei
    componenti di \param{filename} o \param{times} è \val{NULL} e non si ha il
    permesso di scrittura sul file, o non si è proprietari del file o non si
    hanno i privilegi di amministratore.
  \item[\errcode{EPERM}] \param{times} non è \val{NULL}, e non si è
    proprietari del file o non si hanno i privilegi di amministratore.
  \end{errlist}
  ed inoltre \errval{ENOENT} e \errval{EROFS} nel loro significato generico.}
\end{funcproto}

La funzione cambia i tempi di ultimo accesso e di ultima modifica del file
specificato dall'argomento \param{filename}, e richiede come secondo argomento
il puntatore ad una struttura \struct{utimbuf}, la cui definizione è riportata
in fig.~\ref{fig:struct_utimebuf}, con i nuovi valori di detti tempi
(rispettivamente nei campi \var{actime} e \var{modtime}). Se si passa un
puntatore nullo verrà impostato il tempo corrente.

\begin{figure}[!htb]
  \footnotesize \centering
  \begin{minipage}[c]{0.8\textwidth}
    \includestruct{listati/utimbuf.h}
  \end{minipage} 
  \normalsize 
  \caption{La struttura \structd{utimbuf}, usata da \func{utime} per modificare
    i tempi dei file.}
  \label{fig:struct_utimebuf}
\end{figure}

L'effetto della funzione ed i privilegi necessari per eseguirla dipendono dal
valore dell'argomento \param{times}. Se è \val{NULL} la funzione imposta il
tempo corrente ed è sufficiente avere accesso in scrittura al file o essere
proprietari del file o avere i privilegi di amministratore. Se invece si è
specificato un valore diverso la funzione avrà successo solo se si è
proprietari del file o se si hanno i privilegi di amministratore.\footnote{per
  essere precisi la capacità \const{CAP\_FOWNER}, vedi
  sez.~\ref{sec:proc_capabilities}.} In entrambi i casi per verificare la
proprietà del file viene utilizzato l'\ids{UID} effettivo del processo.

Si tenga presente che non è possibile modificare manualmente il tempo di
cambiamento di stato del file, che viene aggiornato direttamente dal kernel
tutte le volte che si modifica l'\textit{inode}, e quindi anche alla chiamata
di \func{utime}.  Questo serve anche come misura di sicurezza per evitare che
si possa modificare un file nascondendo completamente le proprie tracce. In
realtà la cosa resta possibile se si è in grado di accedere al file di
dispositivo, scrivendo direttamente sul disco senza passare attraverso il
filesystem, ma ovviamente in questo modo la cosa è più complicata da
realizzare.\footnote{esistono comunque molti programmi che consentono di farlo
  con relativa semplicità per cui non si dia per scontato che il valore sia
  credibile in caso di macchina compromessa.}

A partire dal kernel 2.6 la risoluzione dei tempi dei file, che nei campi di
tab.~\ref{tab:file_file_times} è espressa in secondi, è stata portata ai
nanosecondi per la gran parte dei filesystem. L'ulteriore informazione può
essere ottenuta attraverso altri campi appositamente aggiunti alla struttura
\struct{stat}. Se si sono definite le macro \macro{\_BSD\_SOURCE} o
\macro{\_SVID\_SOURCE} questi sono \var{st\_atim.tv\_nsec},
\var{st\_mtim.tv\_nsec} e \var{st\_ctim.tv\_nsec} se queste non sono definite,
\var{st\_atimensec}, \var{st\_mtimensec} e \var{st\_mtimensec}. Qualora il
supporto per questa maggior precisione sia assente questi campi aggiuntivi
saranno nulli.

Per la gestione di questi nuovi valori è stata definita una seconda funzione
di sistema, \funcd{utimes}, che consente di specificare tempi con maggior
precisione; il suo prototipo è:

\begin{funcproto}{
\fhead{sys/time.h}
\fdecl{int utimes(const char *filename, struct timeval times[2])}
\fdesc{Modifica i tempi di ultimo accesso e ultima modifica di un file.} 
}
{La funzione ritorna $0$ in caso di successo e $-1$ per un errore, nel qual
  caso \var{errno} assumerà gli stessi valori di \func{utime}.}  
\end{funcproto}
 
La funzione è del tutto analoga alla precedente \func{utime} ma usa come
secondo argomento un vettore di due strutture \struct{timeval}, la cui
definizione è riportata in fig.~\ref{fig:sys_timeval_struct}, che consentono
di indicare i tempi con una precisione del microsecondo. Il primo elemento
di \param{times} indica il valore per il tempo di ultimo accesso, il secondo
quello per il tempo di ultima modifica. Se si indica come secondo argomento un
puntatore nullo di nuovo verrà utilizzato il tempo corrente.

\begin{figure}[!htb]
  \footnotesize \centering
  \begin{minipage}[c]{0.8\textwidth}
    \includestruct{listati/timeval.h}
  \end{minipage} 
  \normalsize 
  \caption{La struttura \structd{timeval} usata per indicare valori di tempo
    con la precisione del microsecondo.}
  \label{fig:sys_timeval_struct}
\end{figure}


% Oltre ad \func{utimes} su Linux sono presenti altre due funzioni per la
% manipolazione dei tempi dei file,\footnote{le due funzioni non sono definite
%   in nessuno standard, ma sono presenti, oltre che su Linux, anche su BSD;
%   sono accessibili definendo \macro{\_DEFAULT\_SOURCE} dalla \acr{glibc} 2.19
%   o \macro{\_GNU\_SOURCE} prima.} la prima è \funcd{futimes} e consente di
% effettuare la modifica utilizzando un file già aperto, il suo prototipo è:

% \begin{funcproto}{
% \fhead{sys/time.h}
% \fdecl{int futimes(int fd, const struct timeval tv[2])}
% \fdesc{Cambia i tempi di un file già aperto.} 
% }

% {La funzione ritornano $0$ in caso di successo e $-1$ per un errore, nel qual
%   caso \var{errno} assumerà uno gli stessi valori di \func{utimes} ed inoltre:
%   \begin{errlist}
%   \item[\errcode{EBADF}] \param{fd} non è un file descriptor.
%   \item[\errcode{ENOSYS}] il filesystem \texttt{/proc} non è accessibile.
%   \end{errlist}}  
% \end{funcproto}

% La seconda funzione, introdotta a partire dal kernel 2.6.22, è
% \funcd{lutimes}, e consente rispettivamente di modificare i tempi di un
% collegamento simbolico; il suo prototipo è:

% \begin{funcproto}{
% \fhead{sys/time.h}
% \fdecl{int lutimes(const char *filename, const struct timeval tv[2])}
% \fdesc{Cambia i tempi di un collegamento simbolico.} 
% }

% {La funzione ritornano $0$ in caso di successo e $-1$ per un errore, nel qual
%   caso \var{errno} assumerà uno gli stessi valori di \func{utimes}, con in più:
%   \begin{errlist}
%   \item[\errcode{ENOSYS}] la funzione non è supportata.
%   \end{errlist}}
% \end{funcproto}

% Le due funzioni hanno lo stesso comportamento di \texttt{utimes} e richiedono
% gli stessi privilegi per poter operare, la differenza è che con \func{futimes}
% si può indicare il file su cui operare se questo è già aperto facendo
% riferimento al suo file descriptor, mentre con \func{lutimes}, nel caso in cui
% \param{filename} sia un collegamento simbolico, saranno modificati i suoi
% tempi invece di quelli del file a cui esso punta.

Oltre ad \func{utimes} su Linux sono presenti altre due funzioni,\footnote{le
  due funzioni non sono definite in nessuno standard, ma sono presenti, oltre
  che su Linux, anche su BSD; sono accessibili definendo
  \macro{\_DEFAULT\_SOURCE} dalla \acr{glibc} 2.19 o \macro{\_GNU\_SOURCE}
  prima.} \funcd{futimes} e \funcd{lutimes}, che consentono rispettivamente
di effettuare la modifica utilizzando un file già aperto o di eseguirla
direttamente su un collegamento simbolico. I relativi prototipi sono:

\begin{funcproto}{
\fhead{sys/time.h}
\fdecl{int futimes(int fd, const struct timeval tv[2])}
\fdesc{Cambia i tempi di un file già aperto.} 
\fdecl{int lutimes(const char *filename, const struct timeval tv[2])}
\fdesc{Cambia i tempi di un collegamento simbolico.} 
}

{Le funzioni ritornano $0$ in caso di successo e $-1$ per un errore, nel qual
  caso \var{errno} assumerà uno gli stessi valori di \func{utimes}, con in più
  per \func{futimes}:
  \begin{errlist}
  \item[\errcode{EBADF}] \param{fd} non è un file descriptor.
  \item[\errcode{ENOSYS}] il filesystem \texttt{/proc} non è accessibile per
    \func{futimes} o la funzione non è supportata per \func{lutimes}.
  \end{errlist}}  
\end{funcproto}

Le due funzioni hanno lo stesso comportamento di \texttt{utimes} e richiedono
gli stessi privilegi per poter operare, la differenza è che con \func{futimes}
si può indicare il file su cui operare se questo è già aperto facendo
riferimento al suo file descriptor, mentre con \func{lutimes} nel caso in
cui \param{filename} sia un collegamento simbolico saranno modificati i suoi
tempi invece di quelli del file a cui esso punta.

Nonostante il kernel nelle versioni più recenti supporti, come accennato,
risoluzioni dei tempi dei file fino al nanosecondo, le funzioni fin qui
esaminate non consentono di impostare valori con questa precisione. Per questo
sono state introdotte due nuove funzioni di sistema, \func{utimensat} (che
vedremo in sez.~\ref{sec:file_openat} insieme alle altre
\textit{at-functions}), e \funcd{futimens}, il cui prototipo è:

\begin{funcproto}{
\fhead{sys/time.h}
\fdecl{int futimens(int fd, const struct timespec times[2])}
\fdesc{Cambia i tempi di un file già aperto.} 
}

{La funzione ritorna $0$ in caso di successo e $-1$ per un errore, nel qual
  caso \var{errno} assumerà uno dei valori:
  \begin{errlist}
  \item[\errcode{EACCES}] si è richiesta l'impostazione del tempo corrente ma
    non si ha il permesso di scrittura sul file, o non si è proprietari del
    file o non si hanno i privilegi di amministratore; oppure il file è
    immutabile (vedi sez.~\ref{sec:file_perm_overview}).
  \item[\errcode{EBADF}] \param{fd} non è un file descriptor valido.
  \item[\errcode{EFAULT}] \param{times} non è un puntatore valido.
  \item[\errcode{EINVAL}] si sono usati dei valori non corretti per i tempi di
    \param{times}.
  \item[\errcode{EPERM}] si è richiesto un cambiamento nei tempi non al tempo
    corrente, ma non si è proprietari del file o non si hanno i privilegi di
    amministratore; oppure il file è immutabile o \textit{append-only} (vedi
    sez.~\ref{sec:file_perm_overview}).
  \end{errlist}
  ed inoltre per entrambe \errval{EROFS} nel suo significato generico.}
\end{funcproto}

La funzione è sostanzialmente una estensione di \func{futimes} che consente di
specificare i tempi con precisione maggiore. Per questo per indicare i valori
dei tempi da impostare utilizza un vettore \param{times} di due strutture
\struct{timespec} che permettono di indicare il valore del tempo con una
precisione fino al nanosecondo (se ne è riportata la definizione in
fig.~\ref{fig:sys_timespec_struct}).

\begin{figure}[!htb]
  \footnotesize \centering
  \begin{minipage}[c]{0.8\textwidth}
    \includestruct{listati/timespec.h}
  \end{minipage} 
  \normalsize 
  \caption{La struttura \structd{timespec} usata per indicare valori di tempo
    con la precisione del nanosecondo.}
  \label{fig:sys_timespec_struct}
\end{figure}

Come per le precedenti funzioni il primo elemento di \param{times} indica il
tempo di ultimo accesso ed il secondo quello di ultima modifica, e se si usa
il valore \val{NULL} verrà impostato il tempo corrente sia per l'ultimo
accesso che per l'ultima modifica.

Nei singoli elementi di \param{times} si possono inoltre utilizzare due valori
speciali per il campo \var{tv\_nsec}: con \constd{UTIME\_NOW} si richiede
l'uso del tempo corrente, mentre con \constd{UTIME\_OMIT} si richiede di non
impostare il tempo. Si può così aggiornare in maniera specifica soltanto uno
fra il tempo di ultimo accesso e quello di ultima modifica. Quando si usa uno
di questi valori speciali per \var{tv\_nsec} il corrispondente valore di
\var{tv\_sec} viene ignorato.

Questa funzione e \func{utimensat} sono una estensione definita nella
revisione POSIX.1-2008 dello standard POSIX; in Linux sono state introdotte a
partire dal kernel 2.6.22, e sono supportate dalla \acr{glibc} a partire dalla
versione 2.6, si tenga presente però che per kernel precedenti il 2.6.26 le
due funzioni sono difettose nel rispetto di alcuni requisiti minori dello
standard e nel controllo della correttezza dei tempi, per i dettagli dei quali
si rimanda alla pagina di manuale.


\section{Il controllo di accesso ai file}
\label{sec:file_access_control}

Una delle caratteristiche fondamentali di tutti i sistemi unix-like è quella
del controllo di accesso ai file, che viene implementato per qualunque
filesystem standard.\footnote{per filesystem standard si intende un filesystem
  che implementi le caratteristiche previste dallo standard POSIX; in Linux
  sono utilizzabili anche filesystem di altri sistemi operativi, che non
  supportano queste caratteristiche.} In questa sezione ne esamineremo i
concetti essenziali e le funzioni usate per gestirne i vari aspetti.


\subsection{I permessi per l'accesso ai file}
\label{sec:file_perm_overview}

Ad ogni file Linux associa sempre, oltre ad un insieme di permessi, l'utente
che ne è proprietario (il cosiddetto \textit{owner}) ed un gruppo di
appartenenza, indicati dagli identificatori di utente e gruppo (\ids{UID} e
\ids{GID}) di cui abbiamo già parlato in
sez.~\ref{sec:proc_access_id}.\footnote{questo è vero solo per filesystem di
  tipo Unix, ad esempio non è vero per il filesystem VFAT di Windows, che non
  fornisce nessun supporto per l'accesso multiutente, e per il quale queste
  proprietà vengono assegnate in maniera fissa con opportune opzioni di
  montaggio.}  Anche questi sono mantenuti sull'\textit{inode} insieme alle
altre proprietà e sono accessibili da programma tramite la funzione
\func{stat} (trattata in sez.~\ref{sec:file_stat}), che restituisce l'utente
proprietario nel campo \var{st\_uid} ed il gruppo proprietario nel campo
\var{st\_gid} della omonima struttura \struct{stat}.

Il controllo di accesso ai file segue un modello abbastanza semplice che
prevede tre permessi fondamentali strutturati su tre livelli di accesso.
Esistono varie estensioni a questo modello,\footnote{come le \textit{Access
    Control List} che sono state aggiunte ai filesystem standard con opportune
  estensioni (vedi sez.~\ref{sec:file_ACL}) per arrivare a meccanismi di
  controllo ancora più sofisticati come il \textit{Mandatory Access Control}
  di \textit{SELinux} e delle altre estensioni come \textit{Smack} o
  \textit{AppArmor}.} ma nella maggior parte dei casi il meccanismo standard è
più che sufficiente a soddisfare tutte le necessità più comuni.  I tre
permessi di base associati ad ogni file sono:
\begin{itemize*}
\item il permesso di lettura (indicato con la lettera \texttt{r}, dall'inglese
  \textit{read}).
\item il permesso di scrittura (indicato con la lettera \texttt{w},
  dall'inglese \textit{write}).
\item il permesso di esecuzione (indicato con la lettera \texttt{x},
  dall'inglese \textit{execute}).
\end{itemize*}
mentre i tre livelli su cui sono divisi i privilegi sono:
\begin{itemize*}
\item i privilegi per l'utente proprietario del file.
\item i privilegi per un qualunque utente faccia parte del gruppo cui
  appartiene il file.
\item i privilegi per tutti gli altri utenti.
\end{itemize*}

L'insieme dei permessi viene espresso con un numero a 12 bit; di questi i nove
meno significativi sono usati a gruppi di tre per indicare i permessi base di
lettura, scrittura ed esecuzione e sono applicati rispettivamente
rispettivamente al proprietario, al gruppo, a tutti gli altri.

\begin{figure}[htb]
  \centering
  \includegraphics[width=8cm]{img/fileperm}
  \caption{Lo schema dei bit utilizzati per specificare i permessi di un file
    contenuti nel campo \var{st\_mode} di \struct{stat}.}
  \label{fig:file_perm_bit}
\end{figure}

I restanti tre bit (noti come \textit{suid bit}, \textit{sgid bit}, e
\textit{sticky bit}) sono usati per indicare alcune caratteristiche più
complesse del meccanismo del controllo di accesso su cui torneremo in seguito
(in sez.~\ref{sec:file_special_perm}), lo schema di allocazione dei bit è
riportato in fig.~\ref{fig:file_perm_bit}.  Come tutte le altre proprietà di
un file anche i permessi sono memorizzati nell'\textit{inode}, e come
accennato in sez.~\ref{sec:file_types} essi vengono restituiti in una parte
del campo \var{st\_mode} della struttura \struct{stat} (si veda di nuovo
fig.~\ref{fig:file_stat_struct}).

\begin{table}[htb]
  \centering
    \footnotesize
  \begin{tabular}[c]{|c|l|}
    \hline
    \textbf{\var{st\_mode}} bit & \textbf{Significato} \\
    \hline 
    \hline 
    \const{S\_IRUSR} & \textit{user-read}, l'utente può leggere.\\
    \const{S\_IWUSR} & \textit{user-write}, l'utente può scrivere.\\
    \const{S\_IXUSR} & \textit{user-execute}, l'utente può eseguire.\\ 
    \hline            
    \const{S\_IRGRP} & \textit{group-read}, il gruppo può leggere.\\
    \const{S\_IWGRP} & \textit{group-write}, il gruppo può scrivere.\\
    \const{S\_IXGRP} & \textit{group-execute}, il gruppo può eseguire.\\
    \hline            
    \const{S\_IROTH} & \textit{other-read}, tutti possono leggere.\\
    \const{S\_IWOTH} & \textit{other-write}, tutti possono scrivere.\\
    \const{S\_IXOTH} & \textit{other-execute}, tutti possono eseguire.\\
    \hline              
  \end{tabular}
  \caption{I bit dei permessi di accesso ai file, come definiti in 
    \texttt{<sys/stat.h>}}
  \label{tab:file_bit_perm}
\end{table}


In genere ci si riferisce ai tre livelli dei privilegi usando le lettere
\texttt{u} (per \textit{user}), \texttt{g} (per \textit{group}) e \texttt{o}
(per \textit{other}), inoltre se si vuole indicare tutti i raggruppamenti
insieme si usa la lettera \texttt{a} (per \textit{all}). Si tenga ben presente
questa distinzione dato che in certi casi, mutuando la terminologia in uso a
suo tempo nel VMS, si parla dei permessi base come di permessi per
\textit{owner}, \textit{group} ed \textit{all}, le cui iniziali possono dar
luogo a confusione.  Le costanti che permettono di accedere al valore numerico
di questi bit nel campo \var{st\_mode}, già viste in
tab.~\ref{tab:file_mode_flags}, sono riportate per chiarezza una seconda volta
in tab.~\ref{tab:file_bit_perm}.

I permessi vengono usati in maniera diversa dalle varie funzioni, e a seconda
che si riferiscano a dei file, dei collegamenti simbolici o delle directory;
qui ci limiteremo ad un riassunto delle regole generali, entrando nei dettagli
più avanti.

La prima regola è che per poter accedere ad un file attraverso il suo
\textit{pathname} occorre il permesso di esecuzione in ciascuna delle
directory che compongono il \textit{pathname}; lo stesso vale per aprire un
file nella directory corrente (per la quale appunto serve il diritto di
esecuzione). Per una directory infatti il permesso di esecuzione significa che
essa può essere attraversata nella risoluzione del \textit{pathname}, e per
questo viene anche chiamato permesso di attraversamento. Esso è sempre
distinto dal permesso di lettura che invece implica che si può leggere il
contenuto della directory.

Questo significa che se si ha il permesso di esecuzione senza permesso di
lettura si potrà lo stesso aprire un file all'interno di una directory (se si
hanno i permessi adeguati per il medesimo) ma non si potrà vederlo con
\cmd{ls} mancando il permesso di leggere il contenuto della directory. Per
crearlo o rinominarlo o cancellarlo invece occorrerà avere anche il permesso
di scrittura per la directory.

Avere il permesso di lettura per un file consente di aprirlo con le opzioni
(si veda quanto riportato in sez.~\ref{sec:file_open_close}) di sola lettura o
di lettura/scrittura e leggerne il contenuto. Avere il permesso di scrittura
consente di aprire un file in sola scrittura o lettura/scrittura e modificarne
il contenuto, lo stesso permesso è necessario per poter troncare il file o per
aggiornare il suo tempo di ultima modifica al tempo corrente, ma non per
modificare arbitrariamente quest'ultimo, operazione per la quale, come per
buona parte delle modifiche effettuate sui metadati del file, occorre esserne
i proprietari.

Non si può creare un file fintanto che non si disponga del permesso di
esecuzione e di quello di scrittura per la directory di destinazione. Gli
stessi permessi occorrono per cancellare un file da una directory (si ricordi
che questo non implica necessariamente la rimozione del contenuto del file dal
disco).

Per la cancellazione non è necessario nessun tipo di permesso per il file
stesso dato che, come illustrato in sez.~\ref{sec:link_symlink_rename} esso
non viene toccato, nella cancellazione infatti viene solo modificato il
contenuto della directory, rimuovendo la voce che ad esso fa riferimento. Lo
stesso vale per poter rinominare o spostare il file in altra directory, in
entrambi i casi occorrerà il permesso di scrittura sulle directory che si
vanno a modificare.

Per poter eseguire un file, che sia un programma compilato od uno script di
shell, od un altro tipo di file eseguibile riconosciuto dal kernel, occorre
oltre al permesso di lettura per accedere al contenuto avere anche il permesso
di esecuzione. Inoltre solo i file regolari possono essere eseguiti. Per i
file di dispositivo i permessi validi sono solo quelli di lettura e scrittura,
che corrispondono al poter eseguire dette operazioni sulla periferica
sottostante. 

I permessi per un collegamento simbolico sono ignorati, contano quelli del
file a cui fa riferimento; per questo in genere il comando \cmd{ls} riporta
per un collegamento simbolico tutti i permessi come concessi. Utente e gruppo
a cui esso appartiene vengono pure ignorati quando il collegamento viene
risolto, vengono controllati solo quando viene richiesta la rimozione del
collegamento e quest'ultimo è in una directory con lo \textit{sticky bit}
impostato (si veda sez.~\ref{sec:file_special_perm}).

La procedura con cui il kernel stabilisce se un processo possiede un certo
permesso (di lettura, scrittura o esecuzione) si basa sul confronto fra
l'utente e il gruppo a cui il file appartiene (i valori di \var{st\_uid} e
\var{st\_gid} accennati in precedenza) e l'\ids{UID} effettivo, il \ids{GID}
effettivo e gli eventuali \ids{GID} supplementari del processo.\footnote{in
  realtà Linux, per quanto riguarda l'accesso ai file, utilizza gli
  identificatori del gruppo \textit{filesystem} (si ricordi quanto esposto in
  sez.~\ref{sec:proc_perms}), ma essendo questi del tutto equivalenti ai primi,
  eccetto il caso in cui si voglia scrivere un server NFS, ignoreremo questa
  differenza.}

Per una spiegazione dettagliata degli identificatori associati ai processi si
veda sez.~\ref{sec:proc_perms}; normalmente, a parte quanto vedremo in
sez.~\ref{sec:file_special_perm}, l'\ids{UID} effettivo e il \ids{GID} effettivo
corrispondono ai valori dell'\ids{UID} e del \ids{GID} dell'utente che ha
lanciato il processo, mentre i \ids{GID} supplementari sono quelli dei gruppi
cui l'utente appartiene.

I passi attraverso i quali viene stabilito se il processo possiede il diritto
di accesso sono i seguenti:
\begin{enumerate*}
\item Se l'\ids{UID} effettivo del processo è zero (corrispondente
  all'amministratore) l'accesso è sempre garantito senza nessun controllo. Per
  questo motivo l'amministratore ha piena libertà di accesso a tutti i file.
\item Se l'\ids{UID} effettivo del processo è uguale all'\ids{UID} del
  proprietario del file (nel qual caso si dice che il processo è proprietario
  del file) allora:
  \begin{itemize*}
  \item se il relativo\footnote{per relativo si intende il bit di
      \textit{user-read} se il processo vuole accedere in lettura, quello di
      \textit{user-write} per l'accesso in scrittura, ecc.} bit dei permessi
    d'accesso dell'utente è impostato, l'accesso è consentito;
  \item altrimenti l'accesso è negato.
  \end{itemize*}
\item Se il \ids{GID} effettivo del processo o uno dei \ids{GID} supplementari
  del processo corrispondono al \ids{GID} del file allora:
  \begin{itemize*}
  \item se il bit dei permessi d'accesso del gruppo è impostato, l'accesso è
    consentito;
  \item altrimenti l'accesso è negato.
  \end{itemize*}
\item Se il bit dei permessi d'accesso per tutti gli altri è impostato,
  l'accesso è consentito, altrimenti l'accesso è negato.
\end{enumerate*}

Si tenga presente che questi passi vengono eseguiti esattamente in
quest'ordine. Questo vuol dire che se un processo è il proprietario di un
file, l'accesso è consentito o negato solo sulla base dei permessi per
l'utente; i permessi per il gruppo non vengono neanche controllati. Lo stesso
vale se il processo appartiene ad un gruppo appropriato, in questo caso i
permessi per tutti gli altri non vengono controllati. 

Questo significa che se si è proprietari di un file ma non si ha il permesso
di scrittura, non vi si potrà scrivere anche se questo fosse scrivibile per
tutti gli altri. Permessi di questo tipo sono ovviamente poco ortodossi, e
comunque, come vedremo in sez.~\ref{sec:file_perm_management}, il proprietario
di un file può sempre modificarne i permessi, e riassegnarsi un eventuale
permesso di scrittura mancante.

\itindbeg{file~attributes} 

A questi che sono i permessi ordinari si aggiungono, per i filesystem che
supportano questa estensione, due permessi speciali mantenuti nei cosiddetti
\textit{file attributes}, che si possono leggere ed impostare con i comandi
\cmd{lsattr} e \cmd{chattr}.\footnote{per l'utilizzo di questi comandi e per
  le spiegazioni riguardo tutti gli altri \textit{file attributes} si rimanda
  alla sezione 1.4.4 di \cite{AGL}.}

Il primo è il cosiddetto attributo di immutabilità (\textit{immutable},
identificato dalla lettera \texttt{i}) che impedisce ogni modifica al file,
\textit{inode} compreso. Questo significa non solo che non se ne può cambiare
il contenuto, ma neanche nessuna delle sue proprietà, ed in particolare non si
può modificare nei permessi o nei tempi o nel proprietario ed inoltre, visto
che non se può modificare il \textit{link count}, non si può neanche
cancellare, rinominare, o creare \textit{hard link} verso di esso.

Il secondo è il cosiddetto attributo di \textit{append-only}, (identificato
dalla lettera \texttt{a}) che consente soltanto la scrittura in coda al
file. Il file cioè può essere soltanto esteso nel contenuto, ma i suoi
metadati, a parte i tempi che però possono essere impostati al valore
corrente, non possono essere modificati in alcun modo, quindi di nuovo non si
potrà cancellare, rinominare, o modificare nei permessi o nelle altre
proprietà.

Entrambi questi attributi attivano queste restrizioni a livello di filesystem,
per cui a differenza dei permessi ordinari esse varranno per qualunque utente
compreso l'amministratore. L'amministratore è l'unico che può attivare o
disattivare questi attributi,\footnote{più precisamente un processo con la
  capacità \const{CAP\_LINUX\_IMMUTABLE}, vedi
  sez.~\ref{sec:proc_capabilities}.} e potendo rimuoverli è comunque capace di
tornare in grado di eseguire qualunque operazione su un file immutabile o
\textit{append-only}.

\itindend{file~attributes}



\subsection{I bit dei permessi speciali}
\label{sec:file_special_perm}

\itindbeg{suid~bit}
\itindbeg{sgid~bit}

Come si è accennato (in sez.~\ref{sec:file_perm_overview}) nei dodici bit del
campo \var{st\_mode} di \struct{stat} che vengono usati per il controllo di
accesso oltre ai bit dei permessi veri e propri, ci sono altri tre bit che
vengono usati per indicare alcune proprietà speciali dei file.  Due di questi
sono i bit detti \acr{suid} (da \textit{set-user-ID bit}) e \acr{sgid} (da
\textit{set-group-ID bit}) che sono identificati dalle costanti
\const{S\_ISUID} e \const{S\_ISGID}.

Come spiegato in dettaglio in sez.~\ref{sec:proc_exec}, quando si lancia un
programma il comportamento normale del kernel è quello di impostare gli
identificatori del gruppo \textit{effective} del nuovo processo al valore dei
corrispondenti del gruppo \textit{real} del processo corrente, che normalmente
corrispondono a quelli dell'utente con cui si è entrati nel sistema.

Se però il file del programma, che ovviamente deve essere
eseguibile,\footnote{anzi più precisamente un binario eseguibile: per motivi
  di sicurezza il kernel ignora i bit \acr{suid} e \acr{sgid} per gli script
  eseguibili.} ha il bit \acr{suid} impostato, il kernel assegnerà come
\ids{UID} effettivo al nuovo processo l'\ids{UID} del proprietario del file al
posto dell'\ids{UID} del processo originario.  Avere il bit \acr{sgid}
impostato ha lo stesso effetto sul \ids{GID} effettivo del processo. É
comunque possibile riconoscere questa situazione perché il cambiamento viene
effettuato solo sugli identificativi del gruppo \textit{effective}, mentre
quelli dei gruppi \textit{real} e \textit{saved} restano quelli dell'utente
che ha eseguito il programma.

I bit \acr{suid} e \acr{sgid} vengono usati per permettere agli utenti normali
di usare programmi che richiedono privilegi speciali. L'esempio classico è il
comando \cmd{passwd} che ha la necessità di modificare il file delle password,
quest'ultimo ovviamente può essere scritto solo dall'amministratore, ma non è
necessario chiamare l'amministratore per cambiare la propria password. Infatti
il comando \cmd{passwd} appartiene in genere all'utente \textit{root} ma ha il
bit \acr{suid} impostato, per cui quando viene lanciato da un utente normale
ottiene comunque  i privilegi di amministratore.

Chiaramente avere un processo che ha privilegi superiori a quelli che avrebbe
normalmente l'utente che lo ha lanciato comporta vari rischi, e questo tipo di
programmi devono essere scritti accuratamente per evitare che possano essere
usati per guadagnare privilegi non consentiti (l'argomento è affrontato in
dettaglio in sez.~\ref{sec:proc_perms}).

La presenza dei bit \acr{suid} e \acr{sgid} su un file può essere rilevata con
il comando \cmd{ls -l}, che visualizza una lettera ``\cmd{s}'' al posto della
``\cmd{x}'' in corrispondenza dei permessi di utente o gruppo. La stessa
lettera ``\cmd{s}'' può essere usata nel comando \cmd{chmod} per impostare
questi bit.  Infine questi bit possono essere controllati all'interno di
\var{st\_mode} con l'uso delle due costanti \const{S\_ISUID} e
\const{S\_IGID}, i cui valori sono riportati in
tab.~\ref{tab:file_mode_flags}.

Gli stessi bit vengono ad assumere un significato completamente diverso per le
directory, normalmente infatti Linux usa la convenzione di SVr4 per indicare
con questi bit l'uso della semantica BSD nella creazione di nuovi file (si
veda sez.~\ref{sec:file_ownership_management} per una spiegazione dettagliata
al proposito).

Infine Linux utilizza il bit \acr{sgid} per un'ulteriore estensione mutuata da
SVr4. Il caso in cui un file ha il bit \acr{sgid} impostato senza che lo sia
anche il corrispondente bit di esecuzione viene utilizzato per attivare per
quel file il \textit{mandatory locking} (affronteremo questo argomento in
dettaglio più avanti, in sez.~\ref{sec:file_mand_locking}).

\itindend{suid~bit}
\itindend{sgid~bit}

\itindbeg{sticky~bit}

L'ultimo dei bit rimanenti, identificato dalla costante \const{S\_ISVTX}, è in
parte un rimasuglio delle origini dei sistemi Unix. A quell'epoca infatti la
memoria virtuale e l'accesso ai file erano molto meno sofisticati e per
ottenere la massima velocità possibile per i programmi usati più comunemente
si poteva impostare questo bit.

L'effetto di questo bit era che il segmento di testo del programma (si veda
sez.~\ref{sec:proc_mem_layout} per i dettagli) veniva scritto nella swap la
prima volta che questo veniva lanciato, e vi permaneva fino al riavvio della
macchina (da questo il nome di \textsl{sticky bit}); essendo la swap un file
continuo o una partizione indicizzata direttamente si poteva risparmiare in
tempo di caricamento rispetto alla ricerca attraverso la struttura del
filesystem. Lo \textsl{sticky bit} è indicato usando la lettera ``\texttt{t}''
al posto della ``\texttt{x}'' nei permessi per gli altri.

Ovviamente per evitare che gli utenti potessero intasare la swap solo
l'amministratore era in grado di impostare questo bit, che venne chiamato
anche con il nome di \textit{saved text bit}, da cui deriva quello della
costante.  Le attuali implementazioni di memoria virtuale e filesystem rendono
sostanzialmente inutile questo procedimento.

Benché ormai non venga più utilizzato per i file, lo \textit{sticky bit} ha
invece assunto un uso importante per le directory;\footnote{lo \textit{sticky
    bit} per le directory è un'estensione non definita nello standard POSIX,
  Linux però la supporta, così come BSD e SVr4.} in questo caso se tale bit è
impostato un file potrà essere rimosso dalla directory soltanto se l'utente ha
il permesso di scrittura su di essa ed inoltre è vera una delle seguenti
condizioni:
\begin{itemize*}
\item l'utente è proprietario del file,
\item l'utente è proprietario della directory,
\item l'utente è l'amministratore.
\end{itemize*}

Un classico esempio di directory che ha questo bit impostato è \file{/tmp}, i
cui permessi infatti di solito sono i seguenti:
\begin{Console}
piccardi@hain:~/gapil$ \textbf{ls -ld /tmp}
drwxrwxrwt    6 root     root         1024 Aug 10 01:03 /tmp
\end{Console}
%$
quindi con lo \textit{sticky bit} bit impostato. In questo modo qualunque
utente nel sistema può creare dei file in questa directory, che come
suggerisce il nome è normalmente utilizzata per la creazione di file
temporanei, ma solo l'utente che ha creato un certo file potrà cancellarlo o
rinominarlo. In questo modo si evita che un utente possa, più o meno
consapevolmente, cancellare i file temporanei creati degli altri utenti.

\itindend{sticky~bit}



\subsection{Le funzioni per la gestione dei permessi dei file}
\label{sec:file_perm_management}

Come visto in sez.~\ref{sec:file_perm_overview} il controllo di accesso ad un
file viene fatto utilizzando l'\ids{UID} ed il \ids{GID} effettivo del
processo; ci sono casi però in cui si può voler effettuare il controllo con
l'\ids{UID} reale ed il \ids{GID} reale, vale a dire usando i valori di
\ids{UID} e \ids{GID} relativi all'utente che ha lanciato il programma, e che,
come accennato in sez.~\ref{sec:file_special_perm} e spiegato in dettaglio in
sez.~\ref{sec:proc_perms}, non è detto siano uguali a quelli effettivi.

Per far questo si può usare la funzione di sistema \funcd{access}, il cui
prototipo è:

\begin{funcproto}{
\fhead{unistd.h}
\fdecl{int access(const char *pathname, int mode)}
\fdesc{Verifica i permessi di accesso.} 
}

{La funzione ritorna $0$ in caso di successo e $-1$ per un errore, nel qual
  caso \var{errno} assumerà uno dei valori: 
  \begin{errlist} 
  \item[\errcode{EACCES}] l'accesso al file non è consentito, o non si ha il
    permesso di attraversare una delle directory di \param{pathname}.
  \item[\errcode{EINVAL}] il valore di \param{mode} non è valido.
  \item[\errcode{EROFS}] si è richiesto l'accesso in scrittura per un file su
    un filesystem montato in sola lettura.
  \item[\errcode{ETXTBSY}] si è richiesto l'accesso in scrittura per un
    eseguibile binario correntemente in esecuzione.
  \end{errlist}
  ed inoltre \errval{EFAULT}, \errval{EIO}, \errval{ELOOP},
  \errval{ENAMETOOLONG}, \errval{ENOENT}, \errval{ENOTDIR} nel loro
  significato generico.}
\end{funcproto}

La funzione verifica i permessi di accesso, indicati da \param{mode}, per il
file indicato da \param{pathname}. I valori possibili per l'argomento
\param{mode} sono esprimibili come combinazione delle costanti numeriche
riportate in tab.~\ref{tab:file_access_mode_val} (attraverso un OR binario
delle stesse). I primi tre valori implicano anche la verifica dell'esistenza
del file, se si vuole verificare solo quest'ultima si può usare \const{F\_OK},
o anche direttamente \func{stat}. Nel caso in cui \param{pathname} si
riferisca ad un collegamento simbolico, questo viene seguito ed il controllo è
fatto sul file a cui esso fa riferimento.

La funzione controlla solo i bit dei permessi di accesso per \param{pathname},
ma occorre poter risolvere quest'ultimo, e se non c'è il permesso di
esecuzione per una qualunque delle sue componenti la funzione fallirà
indipendentemente dai permessi del file.  Si tenga presente poi che il fatto
che una directory abbia permesso di scrittura non significa che vi si possa
scrivere come fosse un file, e che se un file ha il permesso di esecuzione non
è detto che sia eseguibile. La funzione ha successo solo se tutti i permessi
controllati sono disponibili.

\begin{table}[htb]
  \centering
  \footnotesize
  \begin{tabular}{|c|l|}
    \hline
    \textbf{\param{mode}} & \textbf{Significato} \\
    \hline
    \hline
    \constd{R\_OK} & Verifica il permesso di lettura. \\
    \constd{W\_OK} & Verifica il permesso di scrittura. \\
    \constd{X\_OK} & Verifica il permesso di esecuzione. \\
    \constd{F\_OK} & Verifica l'esistenza del file. \\
    \hline
  \end{tabular}
  \caption{Valori possibili per l'argomento \param{mode} della funzione 
    \func{access}.}
  \label{tab:file_access_mode_val}
\end{table}

Un esempio tipico per l'uso di questa funzione è quello di un processo che sta
eseguendo un programma coi privilegi di un altro utente (ad esempio attraverso
l'uso del \textit{suid bit}) che vuole controllare se l'utente originale ha i
permessi per accedere ad un certo file, ma eseguire questo controllo prima di
aprire il file espone al rischio di una \textit{race condition} che apre ad un
possibile \textit{symlink attack} fra il controllo e l'apertura del file. In
questo caso è sempre opportuno usare invece la funzione \func{faccessat} che
tratteremo insieme alle altre \textit{at-functions} in
sez.~\ref{sec:file_openat}.


Del tutto analoghe a \func{access} sono le due funzioni \funcm{euidaccess} e
\funcm{eaccess} che ripetono lo stesso controllo usando però gli
identificatori del gruppo effettivo, verificando quindi le effettive capacità
di accesso ad un file. Le funzioni hanno entrambe lo stesso
prototipo\footnote{in realtà \funcm{eaccess} è solo un sinonimo di
  \funcm{euidaccess} fornita per compatibilità con l'uso di questo nome in
  altri sistemi.} che è del tutto identico a quello di \func{access}. Prendono
anche gli stessi valori e restituiscono gli stessi risultati e gli stessi
codici di errore.

Le due funzioni non sono previste da nessuno standard, ed utilizzabili solo
avendo definito \macro{\_GNU\_SOURCE}; inoltre qualora le si vogliano
utilizzare per verificare che un processo abbia i permessi per accedere ad un
file prima di farlo effettivamente, ci si esporrebbe ad una \textit{race
  condition}, dato che i permessi potrebbero cambiare nell'intervallo fra il
controllo e l'accesso effettivo. Per questo motivo in questo caso è molto più
semplice e sicuro tentare direttamente l'accesso, e trattare opportunamente
l'eventuale fallimento per mancanza di permessi.

Per cambiare i permessi di un file sono invece disponibili due funzioni di
sistema \funcd{chmod} e \funcd{fchmod}, che operano rispettivamente usando il
\textit{pathname} di un file o se questo è già aperto sul relativo file
descriptor; i loro prototipi sono:

\begin{funcproto}{
\fhead{sys/types.h}
\fhead{sys/stat.h}
\fdecl{int chmod(const char *path, mode\_t mode)}
\fdesc{Cambia i permessi del file indicato da \param{path} al valore indicato
  da \param{mode}.} 
\fdecl{int fchmod(int fd, mode\_t mode)}
\fdesc{Analoga alla precedente, ma usa il file descriptor \param{fd} per
  indicare il file.}  

}

{La funzione ritorna $0$ in caso di successo e $-1$ per un errore, nel qual
  caso \var{errno} assumerà uno dei valori: 
  \begin{errlist}
  \item[\errcode{EPERM}] l'\ids{UID} effettivo non corrisponde a quello del
    proprietario del file o non si hanno i privilegi di amministratore.
  \item[\errcode{EROFS}] il file è su un filesystem in sola lettura.
  \end{errlist}
  ed inoltre per entrambe \errval{EIO}, per \func{chmod} \errval{EACCES},
  \errval{EFAULT}, \errval{ELOOP}, \errval{ENAMETOOLONG}, \errval{ENOENT},
  \errval{ENOMEM}, \errval{ENOTDIR}, per \func{fchmod} \errval{EBADF} nel loro
  significato generico.}
\end{funcproto}


Entrambe le funzioni utilizzano come secondo argomento \param{mode}, una
variabile dell'apposito tipo primitivo \type{mode\_t} (vedi
tab.~\ref{tab:intro_primitive_types}), che è una maschera binaria da
utilizzare per specificare i permessi sui file.

\begin{table}[!htb]
  \centering
  \footnotesize
  \begin{tabular}[c]{|c|c|l|}
    \hline
    \textbf{\param{mode}} & \textbf{Valore} & \textbf{Significato} \\
    \hline
    \hline
    \const{S\_ISUID} & 04000 & Set user ID  bit.\\
    \const{S\_ISGID} & 02000 & Set group ID bit.\\
    \const{S\_ISVTX} & 01000 & Sticky bit.\\
    \hline
    \const{S\_IRWXU} & 00700 & L'utente ha tutti i permessi.\\
    \const{S\_IRUSR} & 00400 & L'utente ha il permesso di lettura.\\
    \const{S\_IWUSR} & 00200 & L'utente ha il permesso di scrittura.\\
    \const{S\_IXUSR} & 00100 & L'utente ha il permesso di esecuzione.\\
    \hline
    \const{S\_IRWXG} & 00070 & Il gruppo ha tutti i permessi.\\
    \const{S\_IRGRP} & 00040 & Il gruppo ha il permesso di lettura.\\
    \const{S\_IWGRP} & 00020 & Il gruppo ha il permesso di scrittura.\\
    \const{S\_IXGRP} & 00010 & Il gruppo ha il permesso di esecuzione.\\
    \hline
    \const{S\_IRWXO} & 00007 & Gli altri hanno tutti i permessi.\\
    \const{S\_IROTH} & 00004 & Gli altri hanno il permesso di lettura.\\
    \const{S\_IWOTH} & 00002 & Gli altri hanno il permesso di scrittura.\\
    \const{S\_IXOTH} & 00001 & Gli altri hanno il permesso di esecuzione.\\
    \hline
  \end{tabular}
  \caption{Valori delle costanti usate per indicare i vari bit di
    \param{mode} utilizzato per impostare i permessi dei file.}
  \label{tab:file_permission_const}
\end{table}

Si sono riportate in tab.~\ref{tab:file_permission_const} le costanti con cui
indicare i singoli bit di \param{mode}; si noti come corrispondano agli stessi
valori usati per \var{st\_mode} già visti in tab.~\ref{tab:file_mode_flags}.
Il valore di \param{mode} può essere ottenuto combinando fra loro con un OR
binario le costanti simboliche relative ai vari bit, o specificato
direttamente, come per l'omonimo comando di shell \texttt{chmod}, con un
valore numerico (con la shell in genere lo si scrive in ottale, dato che i bit
dei permessi sono divisibili in gruppi di tre), che si può calcolare
direttamente usando lo schema di utilizzo dei bit illustrato in
fig.~\ref{fig:file_perm_bit}.

Ad esempio i permessi standard assegnati ai nuovi file (lettura e scrittura
per il proprietario, sola lettura per il gruppo e gli altri) sono
corrispondenti al valore ottale \texttt{0644}, un programma invece avrebbe
anche il bit di esecuzione attivo, con un valore ottale di \texttt{0755}, se
infine si volesse attivare anche il bit \acr{suid} il valore ottale da fornire
sarebbe \texttt{4755}.

Il cambiamento dei permessi di un file eseguito attraverso queste funzioni ha
comunque alcune limitazioni, previste per motivi di sicurezza. L'uso delle
funzioni infatti è possibile solo se l'\ids{UID} effettivo del processo
chiamante corrisponde a quello del proprietario del file o se il processo ha i
privilegi di amministratore,\footnote{per la precisione la capacità
  \const{CAP\_FOWNER}, vedi sez.~\ref{sec:proc_capabilities}.} in caso
contrario esse falliranno con un errore di \errcode{EPERM}.

Ma oltre a questa regola generale, di immediata comprensione, esistono delle
limitazioni ulteriori. Per questo motivo, anche se si è proprietari del file,
non tutti i valori possibili di \param{mode} sono permessi o hanno effetto;
in particolare accade che:
\begin{enumerate*}
\item siccome solo l'amministratore può impostare lo \textit{sticky bit}, se
  l'\ids{UID} effettivo del processo non è zero esso viene automaticamente
  cancellato, senza notifica di errore, qualora sia stato indicato
  in \param{mode}.
\item per quanto detto in sez.~\ref{sec:file_ownership_management} riguardo la
  creazione dei nuovi file, si può avere il caso in cui il file creato da un
  processo è assegnato ad un gruppo per il quale il processo non ha privilegi.
  Per evitare che si possa assegnare il bit \acr{sgid} ad un file appartenente
  ad un gruppo per cui non si hanno diritti, questo viene automaticamente
  cancellato da \param{mode}, senza notifica di errore, qualora il gruppo del
  file non corrisponda a quelli associati al processo; la cosa non avviene
  quando l'\ids{UID} effettivo del processo è zero.
\end{enumerate*}

Per alcuni filesystem\footnote{i filesystem più comuni (\acr{ext2},
  \acr{ext3}, \acr{ext4}, \acr{XFS}, \acr{Btrfs}) supportano
  questa caratteristica, che è mutuata da BSD.} è inoltre prevista
un'ulteriore misura di sicurezza, volta a scongiurare l'abuso dei bit
\acr{suid} e \acr{sgid}; essa consiste nel cancellare automaticamente questi
bit dai permessi di un file qualora un processo che non appartenga
all'amministratore\footnote{per la precisione un processo che non dispone
  della capacità \const{CAP\_FSETID}, vedi sez.~\ref{sec:proc_capabilities}.}
effettui una scrittura. In questo modo anche se un utente malizioso scopre un
file \acr{suid} su cui può scrivere, un'eventuale modifica comporterà la
perdita di questo privilegio.

Le funzioni \func{chmod} e \func{fchmod} ci permettono di modificare i
permessi di un file, resta però il problema di quali sono i permessi assegnati
quando il file viene creato. Le funzioni dell'interfaccia nativa di Unix, come
vedremo in sez.~\ref{sec:file_open_close}, permettono di indicare
esplicitamente i permessi di creazione di un file, ma questo non è possibile
per le funzioni dell'interfaccia standard ANSI C che non prevede l'esistenza
di utenti e gruppi, ed inoltre il problema si pone anche per l'interfaccia
nativa quando i permessi non vengono indicati esplicitamente.

\itindbeg{umask} 

Per le funzioni dell'interfaccia standard ANSI C l'unico riferimento possibile
è quello della modalità di apertura del nuovo file (lettura/scrittura o sola
lettura), che però può fornire un valore che è lo stesso per tutti e tre i
permessi di sez.~\ref{sec:file_perm_overview} (cioè $666$ nel primo caso e
$222$ nel secondo). Per questo motivo il sistema associa ad ogni
processo\footnote{è infatti contenuta nel campo \var{umask} della struttura
  \kstruct{fs\_struct}, vedi fig.~\ref{fig:proc_task_struct}.}  una maschera
di bit, la cosiddetta \textit{umask}, che viene utilizzata per impedire che
alcuni permessi possano essere assegnati ai nuovi file in sede di creazione. I
bit indicati nella maschera vengono infatti cancellati dai permessi quando un
nuovo file viene creato.\footnote{l'operazione viene fatta sempre: anche
  qualora si indichi esplicitamente un valore dei permessi nelle funzioni di
  creazione che lo consentono, i permessi contenuti nella \textit{umask}
  verranno tolti.}

La funzione di sistema che permette di impostare il valore di questa maschera
di controllo è \funcd{umask}, ed il suo prototipo è:

\begin{funcproto}{
\fhead{stat.h}
\fdecl{mode\_t umask(mode\_t mask)}
\fdesc{Imposta la maschera dei permessi.} 
}

{La funzione ritorna il precedente valore della maschera, non sono previste
  condizioni di errore.}
\end{funcproto}

La funzione imposta la maschera dei permessi dei bit al valore specificato
da \param{mask}, di cui vengono presi solo i 9 bit meno significativi.  In
genere si usa questa maschera per impostare un valore predefinito che escluda
preventivamente alcuni permessi, il caso più comune è eliminare il permesso di
scrittura per il gruppo e gli altri, corrispondente ad un valore
per \param{mask} pari a $022$.  In questo modo è possibile cancellare
automaticamente i permessi non voluti.  Di norma questo valore viene impostato
una volta per tutte al login (a $022$ se non indicato altrimenti), e gli
utenti non hanno motivi per modificarlo.

\itindend{umask} 


\subsection{La gestione della titolarità dei file}
\label{sec:file_ownership_management}

Vedremo in sez.~\ref{sec:file_open_close} con quali funzioni si possono creare
nuovi file, in tale occasione vedremo che è possibile specificare in sede di
creazione quali permessi applicare ad un file, però non si può indicare a
quale utente e gruppo esso deve appartenere.  Lo stesso problema si presenta
per la creazione di nuove directory (procedimento descritto in
sez.~\ref{sec:file_dir_creat_rem}).

Lo standard POSIX prescrive che l'\ids{UID} del nuovo file corrisponda
all'\ids{UID} effettivo del processo che lo crea; per il \ids{GID} invece
prevede due diverse possibilità:
\begin{itemize*}
\item che il \ids{GID} del file corrisponda al \ids{GID} effettivo del
  processo.
\item che il \ids{GID} del file corrisponda al \ids{GID} della directory in
  cui esso è creato.
\end{itemize*}

In genere BSD usa sempre la seconda possibilità, che viene per questo chiamata
semantica BSD. Linux invece segue normalmente quella che viene chiamata
semantica SVr4: di norma un nuovo file viene creato, seguendo la prima
opzione, con il \ids{GID} del processo, se però la directory in cui viene
creato ha il bit \acr{sgid} dei permessi impostato allora viene usata la
seconda opzione. L'adozione di questa semantica però può essere controllata,
all'interno di alcuni filesystem,\footnote{con il kernel 2.6.25 questi erano
  \acr{ext2}, \acr{ext3}, \acr{ext4}, e \acr{XFS}.}  con l'uso dell'opzione di
montaggio \texttt{grpid}, che se attivata fa passare all'uso della semantica
BSD.

Usare la semantica BSD ha il vantaggio che il \ids{GID} viene sempre
automaticamente propagato, restando coerente a quello della directory di
partenza, in tutte le sotto-directory. La semantica SVr4 offre la possibilità
di scegliere, ma per ottenere lo stesso risultato di coerenza che si ha con
quella di BSD necessita che quando si creano nuove directory venga propagato
il bit \acr{sgid}. Questo è il comportamento predefinito del comando
\cmd{mkdir}, ed è in questo modo ad esempio che le varie distribuzioni
assicurano che le sotto-directory create nella home di un utente restino
sempre con il \ids{GID} del gruppo primario dello stesso.

La presenza del bit \acr{sgid} è inoltre molto comoda quando si hanno
directory contenenti file condivisi da un gruppo di utenti in cui possono
scrivere tutti i membri dello stesso, dato che assicura che i file che gli
utenti vi creano appartengano allo gruppo stesso. Questo non risolve però
completamente i problemi di accesso da parte di altri utenti dello stesso
gruppo, in quanto di default i permessi assegnati al gruppo non sono
sufficienti per un accesso in scrittura; in questo caso si deve aver cura di
usare prima della creazione dei file un valore per \textit{umask} lasci il
permesso di scrittura per il gruppo.\footnote{in tal caso si può assegnare
  agli utenti del gruppo una \textit{umask} di $002$, anche se la soluzione
  migliore in questo caso è usare una ACL di default (vedi
  sez.~\ref{sec:file_ACL}).}

Come avviene nel caso dei permessi esistono anche delle appropriate funzioni
di sistema, \funcd{chown}, \funcd{fchown} e \funcd{lchown}, che permettono di
cambiare sia l'utente che il gruppo a cui un file appartiene; i rispettivi
prototipi sono:

\begin{funcproto}{
\fhead{sys/types.h}
\fhead{sys/stat.h}
\fdecl{int chown(const char *path, uid\_t owner, gid\_t group)}
\fdecl{int fchown(int fd, uid\_t owner, gid\_t group)}
\fdecl{int lchown(const char *path, uid\_t owner, gid\_t group)}
\fdesc{Cambiano proprietario e gruppo proprietario di un file.} 
}

{Le funzioni ritornano $0$ in caso di successo e $-1$ per un errore, nel qual
  caso \var{errno} assumerà uno dei valori: 
  \begin{errlist}
  \item[\errcode{EPERM}] l'\ids{UID} effettivo non corrisponde a quello del
    proprietario del file o non è zero, o utente e gruppo non sono validi.
  \end{errlist}
  ed inoltre per tutte \errval{EROFS} e \errval{EIO}, per \func{chown}
  \errval{EACCES}, \errval{EFAULT}, \errval{ELOOP}, \errval{ENAMETOOLONG},
  \errval{ENOENT}, \errval{ENOMEM}, \errval{ENOTDIR}, per \func{fchown}
  \errval{EBADF} nel loro significato generico.}
\end{funcproto}

Le funzioni cambiano utente e gruppo di appartenenza di un file ai valori
specificati dalle variabili \param{owner} e \param{group}. Con Linux solo
l'amministratore\footnote{o in generale un processo con la capacità
  \const{CAP\_CHOWN}, vedi sez.~\ref{sec:proc_capabilities}.} può cambiare il
proprietario di un file; in questo viene seguita la semantica usata da BSD che
non consente agli utenti di assegnare i loro file ad altri utenti evitando
eventuali aggiramenti delle quote.  L'amministratore può cambiare sempre il
gruppo di un file, il proprietario può cambiare il gruppo solo dei file che
gli appartengono e solo se il nuovo gruppo è il suo gruppo primario o uno dei
gruppi di cui fa parte.

La funzione \func{chown} segue i collegamenti simbolici, per operare
direttamente su un collegamento simbolico si deve usare la funzione
\func{lchown}.\footnote{fino alla versione 2.1.81 in Linux \func{chown} non
  seguiva i collegamenti simbolici, da allora questo comportamento è stato
  assegnato alla funzione \func{lchown}, introdotta per l'occasione, ed è
  stata creata una nuova \textit{system call} per \func{chown} che seguisse i
  collegamenti simbolici.} La funzione \func{fchown} opera su un file aperto,
essa è mutuata da BSD, ma non è nello standard POSIX.  Un'altra estensione
rispetto allo standard POSIX è che specificando $-1$ come valore per
\param{owner} e \param{group} i valori restano immutati.

Quando queste funzioni sono chiamate con successo da un processo senza i
privilegi di amministratore entrambi i bit \acr{suid} e \acr{sgid} vengono
cancellati. Questo non avviene per il bit \acr{sgid} nel caso in cui esso sia
usato (in assenza del corrispondente permesso di esecuzione) per indicare che
per il file è attivo il \textit{mandatory locking} (vedi
sez.~\ref{sec:file_mand_locking}).


\subsection{Un quadro d'insieme sui permessi}
\label{sec:file_riepilogo}

Avendo affrontato in maniera separata il comportamento delle varie funzioni
che operano sui permessi dei file ed avendo trattato in sezioni diverse il
significato dei singoli bit dei permessi, vale la pena di fare un riepilogo in
cui si riassumano le caratteristiche di ciascuno di essi, in modo da poter
fornire un quadro d'insieme.

\begin{table}[!htb]
  \centering
  \footnotesize
  \begin{tabular}[c]{|c|c|c|c|c|c|c|c|c|c|c|c|l|}
    \hline
    \multicolumn{3}{|c|}{special}&
    \multicolumn{3}{|c|}{user}&
    \multicolumn{3}{|c|}{group}&
    \multicolumn{3}{|c|}{other}&
    \multirow{2}{*}{\textbf{Significato per i file}} \\
    \cline{1-12}
    \texttt{s}&\texttt{s}&\texttt{t}&
    \texttt{r}&\texttt{w}&\texttt{x}&
    \texttt{r}&\texttt{w}&\texttt{x}&
    \texttt{r}&\texttt{w}&\texttt{x}& \\
    \hline
    \hline
   1&-&-&-&-&1&-&-&-&-&-&-&Eseguito conferisce l'\ids{UID} effettivo dell'utente.\\
   -&1&-&-&-&1&-&-&-&-&-&-&Eseguito conferisce il \ids{GID} effettivo del gruppo.\\
   -&1&-&-&-&0&-&-&-&-&-&-&Il \textit{mandatory locking} è abilitato.\\
   -&-&1&-&-&-&-&-&-&-&-&-&Non utilizzato.\\
   -&-&-&1&-&-&-&-&-&-&-&-&Permesso di lettura per l'utente.\\
   -&-&-&-&1&-&-&-&-&-&-&-&Permesso di scrittura per l'utente.\\
   -&-&-&-&-&1&-&-&-&-&-&-&Permesso di esecuzione per l'utente.\\
   -&-&-&-&-&-&1&-&-&-&-&-&Permesso di lettura per il gruppo.\\
   -&-&-&-&-&-&-&1&-&-&-&-&Permesso di scrittura per il gruppo.\\
   -&-&-&-&-&-&-&-&1&-&-&-&Permesso di esecuzione per il gruppo.\\
   -&-&-&-&-&-&-&-&-&1&-&-&Permesso di lettura per tutti gli altri.\\
   -&-&-&-&-&-&-&-&-&-&1&-&Permesso di scrittura per tutti gli altri.\\
   -&-&-&-&-&-&-&-&-&-&-&1&Permesso di esecuzione per tutti gli altri.\\
    \hline
    \hline
    \multicolumn{3}{|c|}{special}&
    \multicolumn{3}{|c|}{user}&
    \multicolumn{3}{|c|}{group}&
    \multicolumn{3}{|c|}{other}&
    \multirow{2}{*}{\textbf{Significato per le directory}} \\
    \cline{1-12}
    \texttt{s}&\texttt{s}&\texttt{t}&
    \texttt{r}&\texttt{w}&\texttt{x}&
    \texttt{r}&\texttt{w}&\texttt{x}&
    \texttt{r}&\texttt{w}&\texttt{x}& \\
    \hline
    \hline
    1&-&-&-&-&-&-&-&-&-&-&-&Non utilizzato.\\
    -&1&-&-&-&-&-&-&-&-&-&-&Propaga il gruppo ai nuovi file creati.\\
    -&-&1&-&-&-&-&-&-&-&-&-&Solo il proprietario di un file può rimuoverlo.\\
    -&-&-&1&-&-&-&-&-&-&-&-&Permesso di visualizzazione per l'utente.\\
    -&-&-&-&1&-&-&-&-&-&-&-&Permesso di aggiornamento per l'utente.\\
    -&-&-&-&-&1&-&-&-&-&-&-&Permesso di attraversamento per l'utente.\\
    -&-&-&-&-&-&1&-&-&-&-&-&Permesso di visualizzazione per il gruppo.\\
    -&-&-&-&-&-&-&1&-&-&-&-&Permesso di aggiornamento per il gruppo.\\
    -&-&-&-&-&-&-&-&1&-&-&-&Permesso di attraversamento per il gruppo.\\
    -&-&-&-&-&-&-&-&-&1&-&-&Permesso di visualizzazione per tutti gli altri.\\
    -&-&-&-&-&-&-&-&-&-&1&-&Permesso di aggiornamento per tutti gli altri.\\
    -&-&-&-&-&-&-&-&-&-&-&1&Permesso di attraversamento per tutti gli altri.\\
    \hline
  \end{tabular}
  \caption{Tabella riassuntiva del significato dei bit dei permessi per un
    file e directory.} 
  \label{tab:file_fileperm_bits}
\end{table}

Nella parte superiore di tab.~\ref{tab:file_fileperm_bits} si è riassunto il
significato dei vari bit dei permessi per un file ordinario; per quanto
riguarda l'applicazione dei permessi per proprietario, gruppo ed altri si
ricordi quanto illustrato in sez.~\ref{sec:file_perm_overview}.  Per
compattezza, nella tabella si sono specificati i bit di \textit{suid},
\textit{sgid} e \textit{sticky} con la notazione illustrata anche in
fig.~\ref{fig:file_perm_bit}.  Nella parte inferiore si sono invece riassunti
i significati dei vari bit dei permessi per una directory; anche in questo
caso si è riapplicato ai bit di \textit{suid}, \textit{sgid} e \textit{sticky}
la notazione illustrata in fig.~\ref{fig:file_perm_bit}.

Si ricordi infine che i permessi non hanno alcun significato per i
collegamenti simbolici, mentre per i file di dispositivo hanno senso soltanto
i permessi di lettura e scrittura, che si riflettono sulla possibilità di
compiere dette operazioni sul dispositivo stesso.

Nella tabella si è indicato con il carattere ``-'' il fatto che il valore del
bit in questione non è influente rispetto a quanto indicato nella riga della
tabella; la descrizione del significato fa riferimento soltanto alla
combinazione di bit per i quali è stato riportato esplicitamente un valore.
Si rammenti infine che il valore dei bit dei permessi non ha alcun effetto
qualora il processo possieda i privilegi di amministratore.


\section{Caratteristiche e funzionalità avanzate}
\label{sec:file_dir_advances}

Tratteremo qui alcune caratteristiche e funzionalità avanzate della gestione
di file e directory, affrontando anche una serie di estensioni
dell'interfaccia classica dei sistemi unix-like, principalmente utilizzate a
scopi di sicurezza, che sono state introdotte nelle versioni più recenti di
Linux.

\subsection{Gli attributi estesi}
\label{sec:file_xattr}

\itindbeg{Extended~Attributes}

Nelle sezioni precedenti abbiamo trattato in dettaglio le varie informazioni
che il sistema mantiene negli \textit{inode}, e le varie funzioni che
permettono di modificarle.  Si sarà notato come in realtà queste informazioni
siano estremamente ridotte.  Questo è dovuto al fatto che Unix origina negli
anni '70, quando le risorse di calcolo e di spazio disco erano minime. Con il
venir meno di queste restrizioni è incominciata ad emergere l'esigenza di
poter associare ai file delle ulteriori informazioni astratte (quelli che
abbiamo chiamato genericamente \textsl{metadati}) che però non potevano
trovare spazio nei dati classici mantenuti negli \textit{inode}.

Per risolvere questo problema alcuni sistemi unix-like (e fra questi anche
Linux) hanno introdotto un meccanismo generico, detto \textit{Extended
  Attributes} che consente di associare delle informazioni ulteriori ai
singoli file; ad esempio vengono usati per la gestione delle ACL, che
tratteremo in sez.~\ref{sec:file_ACL} e per le \textit{file capabilities}, che
vedremo in sez.~\ref{sec:proc_capabilities}. Gli \textsl{attributi estesi}
(abbreviati in \textsl{xattr}) non sono altro che delle coppie nome/valore che
sono associate permanentemente ad un oggetto sul filesystem, analoghi di
quello che sono le variabili di ambiente (vedi sez.~\ref{sec:proc_environ})
per un processo.

Altri sistemi (come Solaris, MacOS e Windows) hanno adottato un meccanismo
diverso in cui ad un file sono associati diversi flussi di dati, su cui
possono essere mantenute ulteriori informazioni, che possono essere accedute
con le normali operazioni di lettura e scrittura. Questi non vanno confusi con
gli \textit{Extended Attributes} (anche se su Solaris hanno lo stesso nome),
che sono un meccanismo molto più semplice, che pur essendo limitato (potendo
contenere solo una quantità limitata di informazione) hanno il grande
vantaggio di essere molto più semplici da realizzare, più
efficienti,\footnote{cosa molto importante, specie per le applicazioni che
  richiedono una gran numero di accessi, come le ACL.} e di garantire
l'atomicità di tutte le operazioni.

In Linux gli attributi estesi sono sempre associati al singolo \textit{inode}
e l'accesso viene sempre eseguito in forma atomica, in lettura il valore
corrente viene scritto su un buffer in memoria, mentre la scrittura prevede
che ogni valore precedente sia sovrascritto.

Si tenga presente che non tutti i filesystem supportano gli \textit{Extended
  Attributes}; al momento della scrittura di queste dispense essi sono
presenti solo sui vari \acr{extN}, \acr{ReiserFS}, \acr{JFS},
\acr{XFS}, \acr{Btrfs}, \acr{Lustre} e \acr{OCFS2}.  Inoltre a
seconda della implementazione ci possono essere dei limiti sulla quantità di
attributi che si possono utilizzare.\footnote{ad esempio nel caso di
  \acr{ext2} ed \acr{ext3} è richiesto che essi siano contenuti
  all'interno di un singolo blocco, pertanto con dimensioni massime pari a
  1024, 2048 o 4096 byte a seconda delle dimensioni di quest'ultimo impostate
  in fase di creazione del filesystem, mentre con \textsl{XFS} e
  \textsl{Btrfs} non ci sono limiti ed i dati vengono memorizzati in maniera
  diversa (nell'\textit{inode} stesso, in un blocco a parte, o in una
  struttura ad albero dedicata) per mantenerne la scalabilità; lasciamo i
  dettagli dei vari filesystem alla documentazione accessibile con \texttt{man
    xattr}.} Infine lo spazio utilizzato per mantenere gli attributi estesi
viene tenuto in conto per il calcolo delle quote di utente e gruppo
proprietari del file.

Come meccanismo per mantenere informazioni aggiuntive associate al singolo
file, gli \textit{Extended Attributes} possono avere usi anche molto diversi
fra loro.  Per poterli distinguere allora sono stati suddivisi in
\textsl{classi}, a cui poter applicare requisiti diversi per l'accesso e la
gestione. Per questo motivo il nome di un attributo deve essere sempre
specificato nella forma \texttt{namespace.attribute}, dove \texttt{namespace}
fa riferimento alla classe a cui l'attributo appartiene, mentre
\texttt{attribute} è il nome ad esso assegnato. In tale forma il nome di un
attributo esteso deve essere univoco. Al momento sono state definite le
quattro classi di attributi riportate in
tab.~\ref{tab:extended_attribute_class}.

\begin{table}[htb]
  \centering
  \footnotesize
  \begin{tabular}{|l|p{8cm}|}
    \hline
    \textbf{Nome} & \textbf{Descrizione} \\
    \hline
    \hline
    \texttt{security}&Gli \textit{extended security attributes}: vengono
                      utilizzati dalle estensioni di sicurezza del kernel (i
                      \textit{Linux Security Modules}), per le realizzazione
                      di meccanismi evoluti di controllo di accesso come
                      \textit{SELinux} o le \textit{capabilities} dei
                      file di sez.~\ref{sec:proc_capabilities}.\\ 
    \texttt{system} & Gli \textit{extended system attributes}: sono usati
                      dal kernel per memorizzare dati di sistema associati ai
                      file come le ACL (vedi sez.~\ref{sec:file_ACL}) o le
                      \textit{capabilities} (vedi
                      sez.~\ref{sec:proc_capabilities}).\\
    \texttt{trusted}& I \textit{trusted extended attributes}: vengono
                      utilizzati per poter realizzare in \textit{user space} 
                      meccanismi che consentano di mantenere delle
                      informazioni sui file che non devono essere accessibili
                      ai processi ordinari.\\
    \texttt{user}   & Gli \textit{extended user attributes}: utilizzati per
                      mantenere informazioni aggiuntive sui file (come il
                      \textit{mime-type}, la codifica dei caratteri o del
                      file) accessibili dagli utenti.\\
    \hline
  \end{tabular}
  \caption{I nomi utilizzati valore di \texttt{namespace} per distinguere le
    varie classi di \textit{Extended Attributes}.}
  \label{tab:extended_attribute_class}
\end{table}


Dato che uno degli usi degli \textit{Extended Attributes} è di impiegarli per
realizzare delle estensioni al tradizionale meccanismo dei controlli di
accesso di Unix (come le ACL, \textit{SELinux}, ecc.), l'accesso ai loro
valori viene regolato in maniera diversa a seconda sia della loro classe che
di quali, fra le estensioni che li utilizzano, sono poste in uso. In
particolare, per ciascuna delle classi riportate in
tab.~\ref{tab:extended_attribute_class}, si hanno i seguenti casi:
\begin{basedescript}{\desclabelwidth{1.7cm}\desclabelstyle{\nextlinelabel}}
\item[\texttt{security}] L'accesso agli \textit{extended security attributes}
  dipende dalle politiche di sicurezza stabilite da loro stessi tramite
  l'utilizzo di un sistema di controllo basato sui \textit{Linux Security
    Modules} (ad esempio \textit{SELinux}). Pertanto l'accesso in lettura o
  scrittura dipende dalle politiche di sicurezza implementate all'interno dal
  modulo di sicurezza che si sta utilizzando al momento (ciascuno avrà le
  sue). Se non è stato caricato nessun modulo di sicurezza l'accesso in
  lettura sarà consentito a tutti i processi, mentre quello in scrittura solo
  ai processi con privilegi amministrativi (per la precisione dotati della
  capacità \const{CAP\_SYS\_ADMIN}, vedi sez.~\ref{sec:proc_capabilities}).

\item[\texttt{system}] Anche l'accesso agli \textit{extended system
    attributes} dipende dalle politiche di accesso che il kernel realizza
  anche utilizzando gli stessi valori in essi contenuti. Ad esempio nel caso
  delle ACL (vedi sez.~\ref{sec:file_ACL}) l'accesso è consentito in lettura
  ai processi che hanno la capacità di eseguire una ricerca sul file (cioè
  hanno il permesso di lettura sulla directory che contiene il file) ed in
  scrittura al proprietario del file o ai processi dotati della capacità
  \const{CAP\_FOWNER}, vale a dire una politica di accesso analoga a quella
  impiegata per gli ordinari permessi dei file.

\item[\texttt{trusted}] L'accesso ai \textit{trusted extended attributes}, sia
  per la lettura che per la scrittura, è consentito soltanto ai processi con
  privilegi amministrativi dotati della capacità \const{CAP\_SYS\_ADMIN}. In
  questo modo si possono utilizzare questi attributi per realizzare in
  \textit{user space} dei meccanismi di controllo che accedono ad informazioni
  non disponibili ai processi ordinari.

\item[\texttt{user}] L'accesso agli \textit{extended user attributes} è
  regolato dai normali permessi dei file: occorre avere il permesso di lettura
  per leggerli e quello di scrittura per scriverli o modificarli. Dato l'uso
  di questi attributi si è scelto di applicare al loro accesso gli stessi
  criteri che si usano per l'accesso al contenuto dei file (o delle directory)
  cui essi fanno riferimento. Questa scelta vale però soltanto per i file e le
  directory ordinarie, se valesse in generale infatti si avrebbe un serio
  problema di sicurezza dato che esistono diversi oggetti sul filesystem per i
  quali è normale avere il permesso di scrittura consentito a tutti gli
  utenti, come i collegamenti simbolici, o alcuni file di dispositivo come
  \texttt{/dev/null}. Se fosse possibile usare su di essi gli \textit{extended
    user attributes} un utente qualunque potrebbe inserirvi dati a
  piacere.\footnote{la cosa è stata notata su XFS, dove questo comportamento
    permetteva, non essendovi limiti sullo spazio occupabile dagli
    \textit{Extended Attributes}, di bloccare il sistema riempiendo il disco.}

  La semantica del controllo di accesso indicata inoltre non avrebbe alcun
  senso al di fuori di file e directory: i permessi di lettura e scrittura per
  un file di dispositivo attengono alle capacità di accesso al dispositivo
  sottostante,\footnote{motivo per cui si può formattare un disco anche se
    \texttt{/dev} è su un filesystem in sola lettura.} mentre per i
  collegamenti simbolici questi vengono semplicemente ignorati: in nessuno dei
  due casi hanno a che fare con il contenuto del file, e nella discussione
  relativa all'uso degli \textit{extended user attributes} nessuno è mai stato
  capace di indicare una qualche forma sensata di utilizzo degli stessi per
  collegamenti simbolici o file di dispositivo, e neanche per le \textit{fifo}
  o i socket.

  Per questo motivo essi sono stati completamente disabilitati per tutto ciò
  che non sia un file regolare o una directory.\footnote{si può verificare la
    semantica adottata consultando il file \texttt{fs/xattr.c} dei sorgenti
    del kernel.} Inoltre per le directory è stata introdotta una ulteriore
  restrizione, dovuta di nuovo alla presenza ordinaria di permessi di
  scrittura completi su directory come \texttt{/tmp}. Per questo motivo, per
  evitare eventuali abusi, se una directory ha lo \textit{sticky bit} attivo
  sarà consentito scrivere i suoi \textit{extended user attributes} soltanto
  se si è proprietari della stessa, o si hanno i privilegi amministrativi
  della capacità \const{CAP\_FOWNER}.
\end{basedescript}

Le funzioni per la gestione degli attributi estesi, come altre funzioni di
gestione avanzate specifiche di Linux, non fanno parte della \acr{glibc}, e
sono fornite da una apposita libreria, \texttt{libattr}, che deve essere
installata a parte;\footnote{la versione corrente della libreria è
  \texttt{libattr1}.}  pertanto se un programma le utilizza si dovrà indicare
esplicitamente l'uso della suddetta libreria invocando il compilatore con
l'opzione \texttt{-lattr}.  

Per leggere gli attributi estesi sono disponibili tre diverse funzioni di
sistema, \funcd{getxattr}, \funcd{lgetxattr} e \funcd{fgetxattr}, che
consentono rispettivamente di richiedere gli attributi relativi a un file, a
un collegamento simbolico e ad un file descriptor; i rispettivi prototipi
sono:

\begin{funcproto}{
\fhead{sys/types.h}
\fhead{attr/xattr.h}
\fdecl{ssize\_t getxattr(const char *path, const char *name, void *value,
  size\_t size)}
\fdecl{ssize\_t lgetxattr(const char *path, const char *name, void *value,
  size\_t size)}
\fdecl{ssize\_t fgetxattr(int filedes, const char *name, void *value,
  size\_t size)}
\fdesc{Leggono il valore di un attributo esteso.} 
}

{Le funzioni ritornano un intero positivo che indica la dimensione
  dell'attributo richiesto in caso di successo e $-1$ per un errore, nel qual
  caso \var{errno} assumerà uno dei valori:
  \begin{errlist}
  \item[\errcode{ENOATTR}] l'attributo richiesto non esiste.
  \item[\errcode{ENOTSUP}] gli attributi estesi non sono supportati dal
    filesystem o sono disabilitati.
  \item[\errcode{ERANGE}] la dimensione \param{size} del buffer \param{value}
    non è sufficiente per contenere il risultato.
  \end{errlist}
  ed inoltre tutti gli errori delle analoghe della famiglia \func{stat} con lo
  stesso significato, ed in particolare \errcode{EPERM} se non si hanno i
  permessi di accesso all'attributo.}
\end{funcproto}

Le funzioni \func{getxattr} e \func{lgetxattr} prendono come primo argomento
un \textit{pathname} che indica il file di cui si vuole richiedere un
attributo, la sola differenza è che la seconda, se il \textit{pathname} indica
un collegamento simbolico, restituisce gli attributi di quest'ultimo e non
quelli del file a cui esso fa riferimento. La funzione \func{fgetxattr} prende
invece come primo argomento un numero di file descriptor, e richiede gli
attributi del file ad esso associato.

Tutte e tre le funzioni richiedono di specificare nell'argomento \param{name}
il nome dell'attributo di cui si vuole ottenere il valore. Il nome deve essere
indicato comprensivo di prefisso del \textit{namespace} cui appartiene (uno
dei valori di tab.~\ref{tab:extended_attribute_class}) nella forma
\texttt{namespace.attributename}, come stringa terminata da un carattere NUL.
Il valore dell'attributo richiesto verrà restituito nel buffer puntato
dall'argomento \param{value} per una dimensione massima di \param{size} byte
(gli attributi estesi possono essere costituiti arbitrariamente da dati
testuali o binari); se quest'ultima non è sufficiente si avrà un errore di
\errcode{ERANGE}.

Per evitare di dover indovinare la dimensione di un attributo per tentativi si
può eseguire una interrogazione utilizzando un valore nullo per \param{size};
in questo caso non verrà letto nessun dato, ma verrà restituito come valore di
ritorno della funzione chiamata la dimensione totale dell'attributo esteso
richiesto, che si potrà usare come stima per allocare un buffer di dimensioni
sufficienti.

Si tenga conto che questa è comunque una stima perché anche se le funzioni
restituiscono la dimensione esatta dell'attributo al momento in cui sono
eseguite, questa potrebbe essere modificata in qualunque momento da un
successivo accesso eseguito da un altro processo, pertanto si verifichi sempre
il valore di ritorno ed il codice di errore della funzione usata, senza dare
per scontato che essa abbia sempre successo.

Un secondo gruppo di funzioni sono quelle che consentono di impostare il
valore di un attributo esteso, le funzioni sono \funcd{setxattr},
\funcd{lsetxattr} e \funcd{fsetxattr} e consentono di operare rispettivamente
su un file, su un collegamento simbolico o utilizzando un file descriptor; i
rispettivi prototipi sono:

\begin{funcproto}{
\fhead{sys/types.h}
\fhead{attr/xattr.h}
\fdecl{int setxattr(const char *path, const char *name, const void *value,\\
  \phantom{int setxattr(}size\_t size, int flags)} 
\fdecl{int lsetxattr(const char *path, const char *name, const void *value,\\
  \phantom{int lsetxattr(}size\_t size, int flags)} 
\fdecl{int fsetxattr(int filedes, const char *name, const void *value,\\
  \phantom{int fsetxattr(}size\_t size, int flags)} 
\fdesc{Impostano il valore di un attributo esteso.} 
}

{Le funzioni ritornano un $0$ in caso di successo e $-1$ per un errore, nel qual
  caso \var{errno} assumerà uno dei valori:
  \begin{errlist}
  \item[\errcode{EEXIST}] si è usato il flag \const{XATTR\_CREATE} ma
    l'attributo esiste già.
  \item[\errcode{ENOATTR}] si è usato il flag \const{XATTR\_REPLACE} e
    l'attributo richiesto non esiste.
  \item[\errcode{ENOTSUP}] gli attributi estesi non sono supportati dal
    filesystem o sono disabilitati.
  \end{errlist}
  ed inoltre tutti gli errori delle analoghe della famiglia \func{stat} con lo
  stesso significato ed in particolare \errcode{EPERM} se non si hanno i
  permessi di accesso all'attributo.}
\end{funcproto}

Le tre funzioni prendono come primo argomento un valore adeguato al loro scopo
(un \textit{pathname} le prime due, un file descriptor la terza), usato in
maniera del tutto identica a quanto visto in precedenza per le analoghe che
leggono gli attributi estesi.

Il secondo argomento, \param{name}, deve indicare, anche in questo caso con
gli stessi criteri appena visti per le analoghe \func{getxattr},
\func{lgetxattr} e \func{fgetxattr}, il nome (completo di suffisso)
dell'attributo su cui si vuole operare.  Il valore che verrà assegnato
all'attributo dovrà essere preparato nel buffer puntato da \param{value}, e la
sua dimensione totale (in byte) dovrà essere indicata dall'argomento
\param{size}.

Infine l'argomento \param{flag} consente di controllare le modalità di
sovrascrittura dell'attributo esteso, esso può prendere due valori: con
\constd{XATTR\_REPLACE} si richiede che l'attributo esista, nel qual caso
verrà sovrascritto ed altrimenti si avrà errore; con \constd{XATTR\_CREATE} si
richiede che l'attributo non esista, nel qual caso verrà creato, altrimenti si
avrà errore ed il valore attuale non sarà modificato.  Utilizzando per
\param{flag} un valore nullo l'attributo verrà modificato se è già presente, o
creato se non c'è.

Le funzioni finora illustrate permettono di leggere o scrivere gli attributi
estesi presenti su un file, ma sarebbe altrettanto utile poter sapere quali
sono questi attributi; per questo sono disponibili le ulteriori tre funzioni
di sistema \funcd{listxattr}, \funcd{llistxattr} e \funcd{flistxattr} i cui
prototipi sono:

\begin{funcproto}{
\fhead{sys/types.h}
\fhead{attr/xattr.h}
\fdecl{ssize\_t listxattr(const char *path, char *list, size\_t size)} 
\fdecl{ssize\_t llistxattr(const char *path, char *list, size\_t size)} 
\fdecl{ssize\_t flistxattr(int filedes, char *list, size\_t size)} 
\fdesc{Leggono la lista degli attributi estesi di un file.} 
}

{Le funzioni ritornano un intero positivo che indica la dimensione della lista
  in caso di successo e $-1$ per un errore, nel qual caso \var{errno} assumerà
  uno dei valori:
  \begin{errlist}
  \item[\errcode{ENOTSUP}] gli attributi estesi non sono supportati dal
    filesystem o sono disabilitati.
  \item[\errcode{ERANGE}] la dimensione \param{size} del buffer \param{value}
    non è sufficiente per contenere il risultato.
  \end{errlist}
  ed inoltre tutti gli errori delle analoghe della famiglia \func{stat} con lo
  stesso significato ed in particolare \errcode{EPERM} se non si hanno i
  permessi di accesso all'attributo.}
\end{funcproto}

Come per le precedenti le tre funzioni leggono gli attributi rispettivamente
di un file, un collegamento simbolico o specificando un file descriptor, da
specificare con il loro primo argomento. Gli altri due argomenti, identici per
tutte e tre, indicano rispettivamente il puntatore \param{list} al buffer dove
deve essere letta la lista e la dimensione \param{size} di quest'ultimo.

La lista viene fornita come sequenza non ordinata dei nomi dei singoli
attributi estesi (sempre comprensivi del prefisso della loro classe) ciascuno
dei quali è terminato da un carattere nullo. I nomi sono inseriti nel buffer
uno di seguito all'altro. Il valore di ritorno della funzione indica la
dimensione totale della lista in byte.

Come per le funzioni di lettura dei singoli attributi se le dimensioni del
buffer non sono sufficienti si avrà un errore, ma è possibile ottenere dal
valore di ritorno della funzione una stima della dimensione totale della lista
usando per \param{size} un valore nullo. 

Infine per rimuovere semplicemente un attributo esteso, si ha a disposizione
un ultimo gruppo di funzioni di sistema: \funcd{removexattr},
\funcd{lremovexattr} e \funcd{fremovexattr}; i rispettivi prototipi sono:

\begin{funcproto}{
\fhead{sys/types.h}
\fhead{attr/xattr.h}
\fdecl{int removexattr(const char *path, const char *name)} 
\fdecl{int lremovexattr(const char *path, const char *name)} 
\fdecl{int fremovexattr(int filedes, const char *name)} 
\fdesc{Rimuovono un attributo esteso di un file.} 
}

{Le funzioni ritornano $0$ in caso di successo e $-1$ per un errore, nel qual
  caso \var{errno} assumerà uno dei valori:
  \begin{errlist}
  \item[\errcode{ENOATTR}] l'attributo richiesto non esiste.
  \item[\errcode{ENOTSUP}] gli attributi estesi non sono supportati dal
    filesystem o sono disabilitati.
  \end{errlist}
  ed inoltre tutti gli errori delle analoghe della famiglia \func{stat} con lo
  stesso significato ed in particolare \errcode{EPERM} se non si hanno i
  permessi di accesso all'attributo.}
\end{funcproto}

Le tre funzioni rimuovono un attributo esteso operando rispettivamente su di
un file, su un collegamento simbolico o un file descriptor, che vengono
specificati dal valore passato con il loro primo argomento.  L'attributo da
rimuovere deve essere anche in questo caso indicato con
l'argomento \param{name} secondo le stesse modalità già illustrate in
precedenza per le altre funzioni relative alla gestione degli attributi
estesi.

\itindend{Extended~Attributes}


\subsection{Le \textit{Access  Control List}}
\label{sec:file_ACL}

% la documentazione di sistema è nei pacchetti libacl1-dev e acl 
% vedi anche http://www.suse.de/~agruen/acl/linux-acls/online/

\itindbeg{Access~Control~List~(ACL)}

Il modello classico dei permessi di Unix, per quanto funzionale ed efficiente,
è comunque piuttosto limitato e per quanto possa aver coperto per lunghi anni
le esigenze più comuni con un meccanismo semplice e potente, non è in grado di
rispondere in maniera adeguata a situazioni che richiedono una gestione più
complessa dei permessi di accesso.\footnote{già un requisito come quello di
  dare accesso in scrittura ad alcune persone ed in sola lettura ad altre non
  si può soddisfare in maniera semplice.}

Per questo motivo erano state progressivamente introdotte nelle varie versioni
di Unix dei meccanismi di gestione dei permessi dei file più flessibili, nella
forma delle cosiddette \textit{Access Control List} (indicate usualmente con
la sigla ACL).  Nello sforzo di standardizzare queste funzionalità era stato
creato un gruppo di lavoro il cui scopo era estendere lo standard POSIX 1003
attraverso due nuovi insiemi di specifiche, la POSIX 1003.1e per l'interfaccia
di programmazione e la POSIX 1003.2c per i comandi di shell.

Gli obiettivi del gruppo di lavoro erano però forse troppo ambizioni, e nel
gennaio del 1998 i finanziamenti vennero ritirati senza che si fosse arrivati
alla definizione dello standard richiesto. Dato però che una parte della
documentazione prodotta era di alta qualità venne deciso di rilasciare al
pubblico la diciassettesima bozza del documento, quella che va sotto il nome
di \textit{POSIX 1003.1e Draft 17}, che è divenuta la base sulla quale si
definiscono le cosiddette \textit{Posix ACL}.

A differenza di altri sistemi, come ad esempio FreeBSD, nel caso di Linux si è
scelto di realizzare le ACL attraverso l'uso degli \textit{Extended
  Attributes} (appena trattati in sez.~\ref{sec:file_xattr}), e fornire tutte
le relative funzioni di gestione tramite una libreria, \texttt{libacl} che
nasconde i dettagli implementativi delle ACL e presenta ai programmi una
interfaccia che fa riferimento allo standard POSIX 1003.1e.

Anche in questo caso le funzioni di questa libreria non fanno parte della
\acr{glibc} e devono essere installate a parte;\footnote{la versione corrente
  della libreria è \texttt{libacl1}, e nel caso si usi Debian la si può
  installare con il pacchetto omonimo e con il collegato \texttt{libacl1-dev}
  per i file di sviluppo.}  pertanto se un programma le utilizza si dovrà
indicare esplicitamente l'uso della libreria \texttt{libacl} invocando il
compilatore con l'opzione \texttt{-lacl}. Si tenga presente inoltre che le ACL
devono essere attivate esplicitamente montando il filesystem\footnote{che deve
  supportarle, ma questo è ormai vero per praticamente tutti i filesystem più
  comuni, con l'eccezione di NFS per il quale esiste però il supporto per le
  versioni NFSv2 e NFSv3 del protocollo, con NFSv4 esistono invece delle ACL
  native che hanno una semantica diversa, su di esse possono mappare le ACL
  POSIX, ma l'inverso è possibile solo in forma incompleta.} su cui le si
vogliono utilizzare con l'opzione \texttt{acl} attiva. Dato che si tratta di
una estensione è infatti opportuno utilizzarle soltanto laddove siano
necessarie.

Una ACL è composta da un insieme di voci, e ciascuna voce è a sua volta
costituita da un \textsl{tipo}, da un eventuale \textsl{qualificatore} (deve
essere presente soltanto per le voci di tipo \const{ACL\_USER} e
\const{ACL\_GROUP}) e da un insieme di permessi.  Ad ogni oggetto sul
filesystem si può associare una ACL che ne governa i permessi di accesso,
detta \textit{access ACL}.  Inoltre per le directory si può impostare una ACL
aggiuntiva, detta ``\textit{Default ACL}'', che serve ad indicare quale dovrà
essere la ACL assegnata di default nella creazione di un file all'interno
della directory stessa. Come avviene per i permessi le ACL possono essere
impostate solo del proprietario del file, o da un processo con la capacità
\const{CAP\_FOWNER}.

\begin{table}[htb]
  \centering
  \footnotesize
  \begin{tabular}{|l|p{8cm}|}
    \hline
    \textbf{Tipo} & \textbf{Descrizione} \\
    \hline
    \hline
    \constd{ACL\_USER\_OBJ}& Voce che contiene i diritti di accesso del
                             proprietario del file.\\
    \constd{ACL\_USER}     & Voce che contiene i diritti di accesso per
                             l'utente indicato dal rispettivo
                             qualificatore.\\  
    \constd{ACL\_GROUP\_OBJ}&Voce che contiene i diritti di accesso del
                             gruppo proprietario del file.\\
    \constd{ACL\_GROUP}    & Voce che contiene i diritti di accesso per
                             il gruppo indicato dal rispettivo
                             qualificatore.\\
    \constd{ACL\_MASK}     & Voce che contiene la maschera dei massimi
                             permessi di accesso che possono essere garantiti
                             da voci del tipo \const{ACL\_USER},
                             \const{ACL\_GROUP} e \const{ACL\_GROUP\_OBJ}.\\
    \constd{ACL\_OTHER}    & Voce che contiene i diritti di accesso di chi
                             non corrisponde a nessuna altra voce dell'ACL.\\
    \hline
  \end{tabular}
  \caption{Le costanti che identificano i tipi delle voci di una ACL.}
  \label{tab:acl_tag_types}
\end{table}

L'elenco dei vari tipi di voci presenti in una ACL, con una breve descrizione
del relativo significato, è riportato in tab.~\ref{tab:acl_tag_types}. Tre di
questi tipi, \const{ACL\_USER\_OBJ}, \const{ACL\_GROUP\_OBJ} e
\const{ACL\_OTHER}, corrispondono direttamente ai tre permessi ordinari dei
file (proprietario, gruppo proprietario e tutti gli altri) e per questo una
ACL valida deve sempre contenere una ed una sola voce per ciascuno di questi
tipi.

Una ACL può poi contenere un numero arbitrario di voci di tipo
\const{ACL\_USER} e \const{ACL\_GROUP}, ciascuna delle quali indicherà i
permessi assegnati all'utente e al gruppo indicato dal relativo qualificatore.
Ovviamente ciascuna di queste voci dovrà fare riferimento a un utente o a un
gruppo diverso, e non corrispondenti a quelli proprietari del file. Inoltre se
in una ACL esiste una voce di uno di questi due tipi, è obbligatoria anche la
presenza di una ed una sola voce di tipo \const{ACL\_MASK}, che negli altri
casi è opzionale.

Una voce di tipo \const{ACL\_MASK} serve a mantenere la maschera dei permessi
che possono essere assegnati tramite voci di tipo \const{ACL\_USER},
\const{ACL\_GROUP} e \const{ACL\_GROUP\_OBJ}. Se in una di queste voci si
fosse specificato un permesso non presente in \const{ACL\_MASK} questo
verrebbe ignorato. L'uso di una ACL di tipo \const{ACL\_MASK} è di particolare
utilità quando essa associata ad una \textit{Default ACL} su una directory, in
quanto i permessi così specificati verranno ereditati da tutti i file creati
nella stessa directory. Si ottiene così una sorta di \textit{umask} associata
ad un oggetto sul filesystem piuttosto che a un processo.

Dato che le ACL vengono a costituire una estensione dei permessi ordinari, uno
dei problemi che si erano posti nella loro standardizzazione era appunto
quello della corrispondenza fra questi e le ACL. Come accennato i permessi
ordinari vengono mappati nelle tre voci di tipo \const{ACL\_USER\_OBJ},
\const{ACL\_GROUP\_OBJ} e \const{ACL\_OTHER} che devono essere presenti in
qualunque ACL; un cambiamento ad una di queste voci viene automaticamente
riflesso sui permessi ordinari dei file e viceversa.\footnote{per permessi
  ordinari si intende quelli mantenuti nell'\textit{inode}, che devono restare
  dato che un filesystem può essere montato senza abilitare le ACL.}

In realtà la mappatura è diretta solo per le voci \const{ACL\_USER\_OBJ} e
\const{ACL\_OTHER}, nel caso di \const{ACL\_GROUP\_OBJ} questo vale soltanto
se non è presente una voce di tipo \const{ACL\_MASK}, che è quanto avviene
normalmente se non sono presenti ACL aggiuntive rispetto ai permessi
ordinari. Se invece questa è presente verranno tolti dai permessi di
\const{ACL\_GROUP\_OBJ} (cioè dai permessi per il gruppo proprietario del
file) tutti quelli non presenti in \const{ACL\_MASK}.\footnote{questo diverso
  comportamento a seconda delle condizioni è stato introdotto dalla
  standardizzazione \textit{POSIX 1003.1e Draft 17} per mantenere il
  comportamento invariato sui sistemi dotati di ACL per tutte quelle
  applicazioni che sono conformi soltanto all'ordinario standard \textit{POSIX
    1003.1}.}

Un secondo aspetto dell'incidenza delle ACL sul comportamento del sistema è
quello relativo alla creazione di nuovi file,\footnote{o oggetti sul
  filesystem, il comportamento discusso vale per le funzioni \func{open} e
  \func{creat} (vedi sez.~\ref{sec:file_open_close}), \func{mkdir} (vedi
  sez.~\ref{sec:file_dir_creat_rem}), \func{mknod} e \func{mkfifo} (vedi
  sez.~\ref{sec:file_mknod}).} che come accennato può essere modificato dalla
presenza di una \textit{Default ACL} sulla directory che andrà a contenerli.
Se questa non c'è valgono le regole usuali illustrate in
sez.~\ref{sec:file_perm_management}, per cui essi sono determinati dalla
\textit{umask} del processo, e la sola differenza è che i permessi ordinari da
esse risultanti vengono automaticamente rimappati anche su una ACL di accesso
assegnata automaticamente al nuovo file, che contiene soltanto le tre
corrispondenti voci di tipo \const{ACL\_USER\_OBJ}, \const{ACL\_GROUP\_OBJ} e
\const{ACL\_OTHER}.

Se invece è presente una ACL di default sulla directory che contiene il nuovo
file, essa diventerà automaticamente anche la ACL di accesso di quest'ultimo,
a meno di non aver indicato, nelle funzioni di creazione che lo consentono,
uno specifico valore per i permessi ordinari.\footnote{tutte le funzioni
  citate in precedenza supportano un argomento \var{mode} che indichi un
  insieme di permessi iniziale.} In tal caso saranno eliminati dalle voci
corrispondenti che deriverebbero dalla ACL di default, tutti i permessi non
presenti in tale indicazione.

Dato che questa è la ragione che ha portato alla loro creazione, la principale
modifica introdotta nel sistema con la presenza della ACL è quella alle regole
del controllo di accesso ai file che si sono illustrate in
sez.~\ref{sec:file_perm_overview}.  Come nel caso ordinario per il controllo
vengono sempre utilizzati gli identificatori del gruppo \textit{effective} del
processo, ma in caso di presenza di una ACL sul file, i passi attraverso i
quali viene stabilito se il processo ha il diritto di accesso sono i seguenti:
\begin{enumerate*}
\item Se l'\ids{UID} del processo è nullo (se cioè si è l'amministratore)
  l'accesso è sempre garantito senza nessun controllo.\footnote{più
    precisamente se si devono avere le capacità \const{CAP\_DAC\_OVERRIDE} per
    i file e \const{CAP\_DAC\_READ\_SEARCH} per le directory, vedi
    sez.~\ref{sec:proc_capabilities}.}
\item Se l'\ids{UID} del processo corrisponde al proprietario del file allora:
  \begin{itemize*}
  \item se la voce \const{ACL\_USER\_OBJ} contiene il permesso richiesto,
    l'accesso è consentito;
  \item altrimenti l'accesso è negato.
  \end{itemize*}
\item Se l'\ids{UID} del processo corrisponde ad un qualunque qualificatore
  presente in una voce \const{ACL\_USER} allora:
  \begin{itemize*}
  \item se la voce \const{ACL\_USER} corrispondente e la voce
    \const{ACL\_MASK} contengono entrambe il permesso richiesto, l'accesso è
    consentito;
  \item altrimenti l'accesso è negato.
  \end{itemize*}
\item Se è il \ids{GID} del processo o uno dei \ids{GID} supplementari
  corrisponde al gruppo proprietario del file allora: 
  \begin{itemize*}
  \item se la voce \const{ACL\_GROUP\_OBJ} e una eventuale voce
    \const{ACL\_MASK} (se non vi sono voci di tipo \const{ACL\_GROUP} questa
    può non essere presente) contengono entrambe il permesso richiesto,
    l'accesso è consentito;
  \item altrimenti l'accesso è negato.
  \end{itemize*}
\item Se è il \ids{GID} del processo o uno dei \ids{GID} supplementari
  corrisponde ad un qualunque qualificatore presente in una voce
  \const{ACL\_GROUP} allora:
  \begin{itemize*}
  \item se la voce \const{ACL\_GROUP} corrispondente e la voce
    \const{ACL\_MASK} contengono entrambe il permesso richiesto, l'accesso è
    consentito;
  \item altrimenti l'accesso è negato.
  \end{itemize*}
\item Se la voce \const{ACL\_USER\_OBJ} contiene il permesso richiesto,
  l'accesso è consentito, altrimenti l'accesso è negato.
\end{enumerate*}

I passi di controllo vengono eseguiti esattamente in questa sequenza, e la
decisione viene presa non appena viene trovata una corrispondenza con gli
identificatori del processo. Questo significa che i permessi presenti in una
voce di tipo \const{ACL\_USER} hanno la precedenza sui permessi ordinari
associati al gruppo proprietario del file (vale a dire su
\const{ACL\_GROUP\_OBJ}).

Per la gestione delle ACL lo standard \textit{POSIX 1003.1e Draft 17} ha
previsto delle apposite funzioni ed tutta una serie di tipi di dati dedicati,
arrivando fino a definire un tipo di dato e delle costanti apposite per
identificare i permessi standard di lettura, scrittura ed esecuzione.  Tutte
le operazioni devono essere effettuate attraverso tramite questi tipi di dati,
che incapsulano tutte le informazioni contenute nelle ACL.

La prima di queste funzioni che prendiamo in esame (si ricordi che come per
tutte le altre per poterla usare occorre invocare il compilatore con l'opzione
\texttt{-l acl}) è \funcd{acl\_init}, il cui prototipo è:

\begin{funcproto}{
\fhead{sys/types.h}
\fhead{sys/acl.h}
\fdecl{acl\_t acl\_init(int count)}
\fdesc{Inizializza un'area di lavoro per una ACL.} 
}

{La funzione ritorna un oggetto di tipo \type{acl\_t} in caso di successo e
  \val{NULL} per un errore, nel qual caso \var{errno} assumerà uno dei valori:
  \begin{errlist}
  \item[\errcode{EINVAL}] il valore di \param{count} è negativo.
  \item[\errcode{ENOMEM}] non c'è sufficiente memoria disponibile.
  \end{errlist}
}
\end{funcproto}

La funzione alloca ed inizializza un'area di memoria che verrà usata per
mantenere i dati di una ACL contenente fino ad un massimo di \param{count}
voci. La funzione ritorna un valore di tipo \typed{acl\_t} da usare in tutte le
altre funzioni che operano sulla ACL. La funzione si limita alla allocazione
iniziale e non inserisce nessun valore nella ACL che resta vuota. 

Si tenga presente che pur essendo \typed{acl\_t} un tipo opaco che identifica
``\textsl{l'oggetto}'' ACL, il valore restituito dalla funzione non è altro
che un puntatore all'area di memoria allocata per i dati richiesti. Pertanto
in caso di fallimento verrà restituito un puntatore nullo di tipo
``\code{(acl\_t) NULL}'' e si dovrà, in questa come in tutte le funzioni
seguenti che restituiscono un oggetto di tipo \type{acl\_t}, confrontare il
valore di ritorno della funzione con \val{NULL} (anche se, a voler essere
estremamente pignoli, si dovrebbe usare ``\code{(acl\_t) NULL}'', ma è
sufficiente fare un confronto direttamente con \val{NULL} essendo cura del
compilatore fare le conversioni necessarie).

Una volta che si siano completate le operazioni sui dati di una ACL la memoria
allocata per un oggetto \type{acl\_t} dovrà essere liberata esplicitamente
attraverso una chiamata alla funzione \funcd{acl\_free}, il cui prototipo è:

\begin{funcproto}{
\fhead{sys/types.h}
\fhead{sys/acl.h}
\fdecl{int acl\_free(void *obj\_p)}
\fdesc{Disalloca la memoria riservata per una ACL.} 
}

{La funzione ritorna $0$ in caso di successo e $-1$ per un errore, nel qual
  caso \var{errno} può assumere solo il valore:
  \begin{errlist}
  \item[\errcode{EINVAL}] \param{obj\_p} non è valido.
  \end{errlist}
}
\end{funcproto}

Si noti come la funzione usi come argomento un puntatore di tipo ``\ctyp{void
  *}'', essa infatti può essere usata non solo per liberare la memoria
allocata per i dati di una ACL, ma anche per quella usata per creare le
stringhe di descrizione testuale delle ACL o per ottenere i valori dei
qualificatori della una voce di una ACL. L'uso del tipo generico ``\ctyp{void
  *}'' consente di evitare di eseguire un \textit{cast} al tipo di dato di cui
si vuole effettuare la disallocazione.

Si tenga presente poi che oltre a \func{acl\_init} ci sono molte altre
funzioni che possono allocare memoria per i dati delle ACL, è pertanto
opportuno tenere traccia di tutte le chiamate a queste funzioni perché alla
fine delle operazioni tutta la memoria allocata dovrà essere liberata con
\func{acl\_free}.

Una volta che si abbiano a disposizione i dati di una ACL tramite il
riferimento ad oggetto di tipo \type{acl\_t} questi potranno essere copiati
con la funzione \funcd{acl\_dup}, il cui prototipo è:

\begin{funcproto}{
\fhead{sys/types.h}
\fhead{sys/acl.h}
\fdecl{acl\_t acl\_dup(acl\_t acl)}
\fdesc{Crea una copia di una ACL.} 
}

{La funzione ritorna un oggetto di tipo \type{acl\_t} in caso di successo in
  caso di successo e \val{NULL} per un errore, nel qual caso \var{errno}
  assumerà uno dei valori:
  \begin{errlist}
  \item[\errcode{EINVAL}] l'argomento \param{acl} non è un puntatore valido
    per una ACL.
  \item[\errcode{ENOMEM}] non c'è sufficiente memoria disponibile per eseguire
    la copia.
  \end{errlist}
}
\end{funcproto}

La funzione crea una copia dei dati della ACL indicata tramite l'argomento
\param{acl}, allocando autonomamente tutto spazio necessario alla copia e
restituendo un secondo oggetto di tipo \type{acl\_t} come riferimento a
quest'ultima.  Valgono per questo le stesse considerazioni fatte per il valore
di ritorno di \func{acl\_init}, ed in particolare il fatto che occorrerà
prevedere una ulteriore chiamata esplicita a \func{acl\_free} per liberare la
memoria occupata dalla copia.

Se si deve creare una ACL manualmente l'uso di \func{acl\_init} è scomodo,
dato che la funzione restituisce una ACL vuota, una alternativa allora è usare
\funcd{acl\_from\_mode} che consente di creare una ACL a partire da un valore
di permessi ordinari, il prototipo della funzione è:

\begin{funcproto}{
\fhead{sys/types.h}
\fhead{sys/acl.h}
\fdecl{acl\_t acl\_from\_mode(mode\_t mode)}
\fdesc{Crea una ACL inizializzata con i permessi ordinari.} 
}

{La funzione ritorna un oggetto di tipo \type{acl\_t} in caso di successo e
  \val{NULL} per un errore, nel qual caso \var{errno} può assumere solo
  il valore \errval{ENOMEM}.}
\end{funcproto}


La funzione restituisce una ACL inizializzata con le tre voci obbligatorie
\const{ACL\_USER\_OBJ}, \const{ACL\_GROUP\_OBJ} e \const{ACL\_OTHER} già
impostate secondo la corrispondenza ai valori dei permessi ordinari indicati
dalla maschera passata nell'argomento \param{mode}. Questa funzione è una
estensione usata dalle ACL di Linux e non è portabile, ma consente di
semplificare l'inizializzazione in maniera molto comoda. 

Altre due funzioni che consentono di creare una ACL già inizializzata sono
\funcd{acl\_get\_fd} e \funcd{acl\_get\_file}, che consentono di leggere la
ACL di un file; i rispettivi prototipi sono:

\begin{funcproto}{
\fhead{sys/types.h}
\fhead{sys/acl.h}
\fdecl{acl\_t acl\_get\_file(const char *path\_p, acl\_type\_t type)}
\fdecl{acl\_t acl\_get\_fd(int fd)}
\fdesc{Leggono i dati delle ACL di un file.} 
}

{Le funzioni ritornano un oggetto di tipo \type{acl\_t} in caso di successo e
  \val{NULL} per un errore, nel qual caso \var{errno} assumerà uno dei valori:
  \begin{errlist}
  \item[\errcode{EACCESS}] non c'è accesso per una componente di
    \param{path\_p} o si è richiesta una ACL di default per un file (solo per
    \func{acl\_get\_file}).
  \item[\errcode{EINVAL}] \param{type} non ha un valore valido (solo per
    \func{acl\_get\_file}).
  \item[\errcode{ENOTSUP}] il filesystem cui fa riferimento il file non
    supporta le ACL.
  \end{errlist}
  ed inoltre \errval{ENOMEM} per entrambe, \errval{EBADF} per
  \func{acl\_get\_fd}, e \errval{ENAMETOOLONG}, \errval{ENOENT},
  \errval{ENOTDIR}, per \func{acl\_get\_file} nel loro significato generico. }
\end{funcproto}

Le due funzioni ritornano, con un oggetto di tipo \type{acl\_t}, il valore
della ACL correntemente associata ad un file, che può essere identificato
tramite un file descriptor usando \func{acl\_get\_fd} o con un
\textit{pathname} usando \func{acl\_get\_file}.

Nel caso di quest'ultima funzione, che può richiedere anche la ACL relativa ad
una directory, il secondo argomento \param{type} consente di specificare se si
vuole ottenere la ACL di default o quella di accesso. Questo argomento deve
essere di tipo \typed{acl\_type\_t} e può assumere solo i due valori riportati
in tab.~\ref{tab:acl_type}.

\begin{table}[htb]
  \centering
  \footnotesize
  \begin{tabular}{|l|l|}
    \hline
    \textbf{Tipo} & \textbf{Descrizione} \\
    \hline
    \hline
    \constd{ACL\_TYPE\_ACCESS} & Indica una ACL di accesso.\\
    \constd{ACL\_TYPE\_DEFAULT}& Indica una ACL di default.\\  
    \hline
  \end{tabular}
  \caption{Le costanti che identificano il tipo di ACL.}
  \label{tab:acl_type}
\end{table}

Si tenga presente che nel caso di \func{acl\_get\_file} occorrerà che il
processo chiamante abbia privilegi di accesso sufficienti a poter leggere gli
attributi estesi dei file (come illustrati in sez.~\ref{sec:file_xattr});
inoltre una ACL di tipo \const{ACL\_TYPE\_DEFAULT} potrà essere richiesta
soltanto per una directory, e verrà restituita solo se presente, altrimenti
verrà restituita una ACL vuota.

Infine si potrà creare una ACL direttamente dalla sua rappresentazione
testuale con la funzione \funcd{acl\_from\_text}, il cui prototipo è:

\begin{funcproto}{
\fhead{sys/types.h}
\fhead{sys/acl.h}
\fdecl{acl\_t acl\_from\_text(const char *buf\_p)}
\fdesc{Crea una ACL a partire dalla sua rappresentazione testuale.} 
}

{La funzione ritorna un oggetto di tipo \type{acl\_t} in caso di successo e
  \val{NULL} per un errore, nel qual caso \var{errno} assumerà uno
  dei valori:
  \begin{errlist}
  \item[\errcode{EINVAL}] la rappresentazione testuale all'indirizzo
    \param{buf\_p} non è valida.
  \item[\errcode{ENOMEM}] non c'è memoria sufficiente per allocare i dati.
   \end{errlist}
}
\end{funcproto}

La funzione prende come argomento il puntatore ad un buffer dove si è inserita
la rappresentazione testuale della ACL che si vuole creare, la memoria
necessaria viene automaticamente allocata ed in caso di successo viene
restituito come valore di ritorno un oggetto di tipo \type{acl\_t} con il
contenuto della stessa, che, come per le precedenti funzioni, dovrà essere
disallocato esplicitamente al termine del suo utilizzo.

La rappresentazione testuale di una ACL è quella usata anche dai comandi
ordinari per la gestione delle ACL (\texttt{getfacl} e \texttt{setfacl}), che
prevede due diverse forme, estesa e breve, entrambe supportate da
\func{acl\_from\_text}. La forma estesa prevede che sia specificata una voce
per riga, nella forma:
\begin{Example}
tipo:qualificatore:permessi
\end{Example}
dove il tipo può essere uno fra \texttt{user}, \texttt{group}, \texttt{other}
e \texttt{mask}. Il qualificatore è presente solo per \texttt{user} e
\texttt{group} ed indica l'utente o il gruppo a cui la voce si riferisce,
mentre i permessi sono espressi con una tripletta di lettere analoga a quella
usata per i permessi dei file, vale a dire ``\texttt{r}'' per il permesso di
lettura, ``\texttt{w}'' per il permesso di scrittura, ``\texttt{x}'' per il
permesso di esecuzione (scritti in quest'ordine) e ``\texttt{-}'' per
l'assenza del permesso.

Un possibile esempio di rappresentazione della ACL di un file ordinario a cui,
oltre ai permessi ordinari, si è aggiunto un altro utente con un accesso in
lettura, è il seguente:
\begin{Example}
user::rw-
group::r--
other::r--
user:piccardi:r--
group:gapil:r--
\end{Example}

Va precisato che i due tipi \texttt{user} e \texttt{group} sono usati
rispettivamente per indicare delle voci relative ad utenti e gruppi (cioè per
voci di tipo \const{ACL\_USER\_OBJ} e \const{ACL\_USER} per \texttt{user} e
\const{ACL\_GROUP\_OBJ} e \const{ACL\_GROUP} per \texttt{group}) applicate sia
a quelli proprietari del file che a quelli generici. Quelle dei proprietari si
riconoscono per l'assenza di un qualificatore, ed in genere si scrivono per
prima delle altre.  Il significato delle voci di tipo \texttt{mask} e
\texttt{mark} è evidente. Usando questa forma estesa si possono anche inserire
dei commenti nel testo precedendoli con il carattere ``\texttt{\#}''.

La forma breve prevede invece la scrittura delle singole voci su una riga,
separate da virgole; come specificatori del tipo di voce si possono usare le
iniziali dei valori usati nella forma estesa (cioè ``\texttt{u}'',
``\texttt{g}'', ``\texttt{o}'' e ``\texttt{m}''), mentre le altri parte della
voce sono le stesse. In questo caso non sono consentiti permessi.

Per la conversione inversa, che consente di ottenere la rappresentazione
testuale di una ACL, sono invece disponibili due funzioni. La prima delle due,
di uso più immediato, è \funcd{acl\_to\_text}, ed il suo prototipo è:

\begin{funcproto}{
\fhead{sys/types.h}
\fhead{sys/acl.h}
\fdecl{char *acl\_to\_text(acl\_t acl, ssize\_t *len\_p)}
\fdesc{Produce la rappresentazione testuale di una ACL.} 
}

{La funzione ritorna il puntatore ad una stringa con la rappresentazione
  testuale della ACL in caso di successo e \var{NULL} per un errore, nel qual
  caso \var{errno} assumerà uno dei valori:
  \begin{errlist}
  \item[\errcode{EINVAL}] o \param{acl} non è un puntatore ad una ACL o la ACL
    che esso indica non è valida o non può esser tradotta in forma testuale.
  \item[\errcode{ENOMEM}] non c'è memoria sufficiente per allocare i dati.
  \end{errlist}
}  
\end{funcproto}

La funzione restituisce il puntatore ad una stringa, terminata da un NUL,
contenente la rappresentazione testuale in forma estesa della ACL passata come
argomento, ed alloca automaticamente la memoria necessaria. Questa dovrà poi
essere liberata, quando non più necessaria, con \func{acl\_free}.  Se
nell'argomento \param{len\_p} si passa come valore il puntatore ad una
variabile intera, in questa verrà restituita (come \textit{value result
  argument}) la dimensione della stringa con la rappresentazione testuale, non
comprendente il carattere nullo finale.

La seconda funzione, che permette di controllare con una gran dovizia di
particolari la generazione della stringa contenente la rappresentazione
testuale della ACL, è \funcd{acl\_to\_any\_text}, ed il suo prototipo è:

\begin{funcproto}{
\fhead{sys/types.h}
\fhead{sys/acl.h}
\fdecl{char *acl\_to\_any\_text(acl\_t acl, const char *prefix, char
    separator, int options)}
\fdesc{Produce la rappresentazione testuale di una ACL.} 
}

{La funzione ritorna il puntatore ad una stringa con la rappresentazione
  testuale della ACL in caso di successo e \val{NULL} per un errore, nel qual
  caso \var{errno} assumerà uno dei valori:
  \begin{errlist}
  \item[\errcode{EINVAL}] la ACL indicata da \param{acl} non è valida.
  \item[\errcode{ENOMEM}] non c'è memoria sufficiente per allocare i dati.
  \end{errlist}
}  
\end{funcproto}

La funzione converte in formato testo la ACL indicata dall'argomento
\param{acl}, usando il carattere \param{separator} come separatore delle
singole voci; se l'argomento \param{prefix} non è nullo la stringa da esso
indicata viene utilizzata come prefisso per le singole voci. 

L'ultimo argomento, \param{options}, consente di controllare la modalità con
cui viene generata la rappresentazione testuale. Un valore nullo fa sì che
vengano usati gli identificatori standard \texttt{user}, \texttt{group},
\texttt{other} e \texttt{mask} con i nomi di utenti e gruppi risolti rispetto
ai loro valori numerici. Altrimenti si può specificare un valore in forma di
maschera binaria, da ottenere con un OR aritmetico dei valori riportati in
tab.~\ref{tab:acl_to_text_options}.

\begin{table}[htb]
  \centering
  \footnotesize
  \begin{tabular}{|l|p{8cm}|}
    \hline
    \textbf{Tipo} & \textbf{Descrizione} \\
    \hline
    \hline
    \constd{TEXT\_ABBREVIATE}    & Stampa le voci in forma abbreviata.\\
    \constd{TEXT\_NUMERIC\_IDS}  & Non effettua la risoluzione di
                                   \ids{UID} e \ids{GID} lasciando i valori
                                   numerici.\\
    \constd{TEXT\_SOME\_EFFECTIVE}&Per ciascuna voce che contiene permessi che
                                   vengono eliminati dalla \const{ACL\_MASK}
                                   viene generato un commento con i permessi 
                                   effettivamente risultanti; il commento è
                                   separato con un tabulatore.\\
    \constd{TEXT\_ALL\_EFFECTIVE}& Viene generato un commento con i permessi
                                   effettivi per ciascuna voce che contiene
                                   permessi citati nella \const{ACL\_MASK},
                                   anche quando questi non vengono modificati
                                   da essa; il commento è separato con un
                                   tabulatore.\\
    \constd{TEXT\_SMART\_INDENT} & Da usare in combinazione con le precedenti
                                   opzioni \const{TEXT\_SOME\_EFFECTIVE} e
                                   \const{TEXT\_ALL\_EFFECTIVE}, aumenta
                                   automaticamente il numero di spaziatori
                                   prima degli eventuali commenti in modo da
                                   mantenerli allineati.\\
    \hline
  \end{tabular}
  \caption{Possibili valori per l'argomento \param{options} di
    \func{acl\_to\_any\_text}.} 
  \label{tab:acl_to_text_options}
\end{table}

Come per \func{acl\_to\_text} anche in questo caso il buffer contenente la
rappresentazione testuale dell'ACL, di cui la funzione restituisce
l'indirizzo, viene allocato automaticamente, e dovrà essere esplicitamente
disallocato con una chiamata ad \func{acl\_free}. Si tenga presente infine che
questa funzione è una estensione specifica di Linux, e non è presente nella
bozza dello standard POSIX.1e.

Per quanto utile per la visualizzazione o l'impostazione da riga di comando
delle ACL, la forma testuale non è la più efficiente per poter memorizzare i
dati relativi ad una ACL, ad esempio quando si vuole eseguirne una copia a
scopo di archiviazione. Per questo è stata prevista la possibilità di
utilizzare una rappresentazione delle ACL in una apposita forma binaria
contigua e persistente. È così possibile copiare il valore di una ACL in un
buffer e da questa rappresentazione tornare indietro e generare una ACL.

Lo standard POSIX.1e prevede a tale scopo tre funzioni, la prima e più
semplice è \funcd{acl\_size}, che consente di ottenere la dimensione che avrà
la citata rappresentazione binaria, in modo da poter allocare per essa un
buffer di dimensione sufficiente, il suo prototipo è:

\begin{funcproto}{
\fhead{sys/types.h}
\fhead{sys/acl.h}
\fdecl{ssize\_t acl\_size(acl\_t acl)}
\fdesc{Determina la dimensione della rappresentazione binaria di una ACL.} 
}

{La funzione ritorna la dimensione in byte della rappresentazione binaria
  della ACL in caso di successo e $-1$ per un errore, nel qual caso
  \var{errno} può assumere solo il valore:
  \begin{errlist}
  \item[\errcode{EINVAL}] la ACL indicata da \param{acl} non è valida.
  \end{errlist}
}  
\end{funcproto}

Ottenuta con \func{acl\_size} la dimensione del buffer necessaria per potervi
memorizzare una ACL questo dovrà potrà essere allocato direttamente con
\func{malloc}, ed a questo punto vi si potrà salvare la rappresentazione
binaria della precedente ACL utilizzando la funzione \funcd{acl\_copy\_ext},
il cui prototipo è:

\begin{funcproto}{
\fhead{sys/types.h}
\fhead{sys/acl.h}
\fdecl{ssize\_t acl\_copy\_ext(void *buf\_p, acl\_t acl, ssize\_t size)}
\fdesc{Ottiene la rappresentazione binaria di una ACL.} 
}

{La funzione ritorna la dimensione in byte della rappresentazione binaria
  della ACL in caso di successo e $-1$ per un errore, nel qual caso
  \var{errno} assumerà uno dei valori:
  \begin{errlist}
  \item[\errcode{EINVAL}] la ACL indicata da \param{acl} non è valida, o
    \param{acl} non è un puntatore ad una ACL o \param{size} è negativo o
    nullo.
  \item[\errcode{ERANGE}] il valore di \param{size} è più piccolo della
    dimensione della rappresentazione binaria della ACL.
  \end{errlist}
}  
\end{funcproto}

La funzione scriverà la rappresentazione binaria della ACL indicata da
\param{acl} sul buffer di dimensione \param{size} all'indirizzo
\param{buf\_p}, restituendo la dimensione della stessa come valore di
ritorno. Qualora la dimensione della rappresentazione ecceda il valore di
\param{size} la funzione fallirà con un errore di \errcode{ERANGE}. La
funzione non ha nessun effetto sulla ACL indicata da \param{acl}.

Viceversa se si vuole ripristinare una ACL a partire da una rappresentazione
binaria si potrà usare la funzione \funcd{acl\_copy\_int}, il cui prototipo è:

\begin{funcproto}{
\fhead{sys/types.h} 
\fhead{sys/acl.h}
\fdecl{acl\_t acl\_copy\_int(const void *buf\_p)}
\fdesc{Ripristina la rappresentazione binaria di una ACL.} 
}

{La funzione ritorna un oggetto di tipo \type{acl\_t} in caso di successo e
  \val{NULL} per un errore, nel qual caso \var{errno} assumerà uno dei valori:
  \begin{errlist}
  \item[\errcode{EINVAL}] il buffer all'indirizzo \param{buf\_p} non contiene
    una rappresentazione corretta di una ACL.
  \item[\errcode{ENOMEM}] non c'è memoria sufficiente per allocare un oggetto
    \type{acl\_t} per la ACL richiesta.
  \end{errlist}
}
\end{funcproto}

La funzione alloca autonomamente un oggetto di tipo \type{acl\_t}, restituito
come valore di ritorno, con il contenuto della ACL rappresentata dai dati del
buffer puntato da \param{buf\_p}. Al solito l'oggetto \type{acl\_t} dovrà
essere disallocato esplicitamente al termine del suo utilizzo.

Una volta che si disponga della ACL desiderata, questa potrà essere impostata
su un file o una directory. Per impostare una ACL sono disponibili due
funzioni: \funcd{acl\_set\_file}, che opera sia su file che su directory
usando un \textit{pathname}, e \funcd{acl\_set\_file} che opera solo su file
usando un file descriptor; i rispettivi prototipi sono:

\begin{funcproto}{
\fhead{sys/types.h}
\fhead{sys/acl.h}
\fdecl{int acl\_set\_file(const char *path, acl\_type\_t type, acl\_t acl)}
\fdesc{Imposta una ACL su un file o una directory.} 
\fdecl{int acl\_set\_fd(int fd, acl\_t acl)}
\fdesc{Imposta una ACL su un file descriptor.} 
}

{Le funzioni ritornano $0$ in caso di successo e $-1$ per un errore, nel qual
  caso \var{errno} assumerà uno dei valori: 
  \begin{errlist}
  \item[\errcode{EACCES}] o un generico errore di accesso a \param{path} o il
    valore di \param{type} specifica una ACL il cui tipo non può essere
    assegnato a \param{path}.
  \item[\errcode{EINVAL}] o \param{acl} non è una ACL valida, o \param{type}
    ha un valore non corretto per \func{acl\_set\_file} o o ha più voci di
    quante se ne possono assegnare al file per \func{acl\_set\_fd}.
  \item[\errcode{ENOSPC}] non c'è spazio disco sufficiente per contenere i
    dati aggiuntivi della ACL.
  \item[\errcode{ENOTSUP}] si è cercato di impostare una ACL su un file
    contenuto in un filesystem che non supporta le ACL.
  \end{errlist}
  ed inoltre nel loro significato generico \errval{EPERM}, \errval{EROFS} per
  entrambe, \errval{EBADF} per \func{acl\_set\_fd}, \errval{ENAMETOOLONG},
  \errval{ENOENT}, \errval{ENOTDIR} per \func{acl\_set\_file}.}
\end{funcproto}

Con \func{acl\_set\_file} si assegna la ACL contenuta in \param{acl} al file o
alla directory indicate da \param{path}, con \param{type} che indica il tipo
di ACL con le costanti di tab.~\ref{tab:acl_type}; si tenga presente però che
le ACL di default possono essere solo impostate qualora \param{path} indichi
una directory. Inoltre perché la funzione abbia successo la ACL dovrà essere
valida, e contenere tutti le voci necessarie, con l'eccezione si specifica una
ACL vuota per cancellare la ACL di default associata a \param{path}, valida
solo in caso di directory.\footnote{questo però è una estensione della
  implementazione delle ACL di Linux, la bozza di standard POSIX.1e prevedeva
  l'uso della apposita funzione \funcd{acl\_delete\_def\_file}, che prende
  come unico argomento il \textit{pathname} della directory di cui si vuole
  cancellare l'ACL di default, per i dettagli si ricorra alla pagina di
  manuale.}

La seconda funzione, \func{acl\_set\_fd}, è del tutto è analoga alla prima, ma
non dovendo avere a che fare con directory (e la conseguente possibilità di
avere una ACL di default) non necessita che si specifichi il tipo di ACL, che
sarà sempre di accesso, e prende come unico argomento, a parte il file
descriptor, la ACL da impostare.

Le funzioni viste finora operano a livello di una intera ACL, eseguendo in una
sola volta tutte le operazioni relative a tutte le voci in essa contenuta. In
generale è possibile modificare un singolo valore all'interno di una singola
voce direttamente con le funzioni previste dallo standard POSIX.1e.  Queste
funzioni però sono alquanto macchinose da utilizzare per cui è molto più
semplice operare direttamente sulla rappresentazione testuale. Questo è il
motivo per non tratteremo nei dettagli dette funzioni, fornendone solo una
descrizione sommaria; chi fosse interessato può ricorrere alle pagine di
manuale.

Se si vuole operare direttamente sui contenuti di un oggetto di tipo
\type{acl\_t} infatti occorre fare riferimento alle singole voci tramite gli
opportuni puntatori di tipo \typed{acl\_entry\_t}, che possono essere ottenuti
dalla funzione \funcm{acl\_get\_entry} (per una voce esistente) o dalla
funzione \funcm{acl\_create\_entry} per una voce da aggiungere. Nel caso della
prima funzione si potrà poi ripetere la lettura per ottenere i puntatori alle
singole voci successive alla prima.

\begin{figure}[!htb]
  \footnotesize \centering
  \begin{minipage}[c]{\codesamplewidth}
    \includecodesample{listati/mygetfacl.c}
  \end{minipage} 
  \normalsize
  \caption{Corpo principale del programma \texttt{mygetfacl.c}.}
  \label{fig:proc_mygetfacl}
\end{figure}

Una volta ottenuti detti puntatori si potrà operare sui contenuti delle
singole voci: con le funzioni \funcm{acl\_get\_tag\_type},
\funcm{acl\_get\_qualifier}, \funcm{acl\_get\_permset} si potranno leggere
rispettivamente tipo, qualificatore e permessi, mentre con le corrispondenti
\funcm{acl\_set\_tag\_type}, \funcm{acl\_set\_qualifier},
\funcm{acl\_set\_permset} si potranno impostare i valori; in entrambi i casi
vengono utilizzati tipi di dato ad hoc, descritti nelle pagine di manuale. Si
possono poi copiare i valori di una voce da una ACL ad un altra con
\funcm{acl\_copy\_entry} o eliminare una voce da una ACL con
\funcm{acl\_delete\_entry} e verificarne la validità prima di usarla con
\funcm{acl\_valid} o \funcm{acl\_check}.

\itindend{Access~Control~List~(ACL)}

Come esempio di utilizzo di queste funzioni nei sorgenti allegati alla guida
si è distribuito il programma \texttt{mygetfacl.c}, che consente di leggere le
ACL di un file, passato come argomento.

La sezione principale del programma, da cui si è rimossa la sezione sulla
gestione delle opzioni, è riportata in fig.~\ref{fig:proc_mygetfacl}. Il
programma richiede un unico argomento (\texttt{\small 16-20}) che indica il
file di cui si vuole leggere la ACL. Se questo è presente si usa
(\texttt{\small 22}) la funzione \func{get\_acl\_file} per leggerne la ACL, e
si controlla (\texttt{\small 23-26}) se l'operazione ha successo, uscendo con
un messaggio di errore in caso contrario. 

Ottenuta la ACL la si converte in formato testuale (\texttt{\small 27}) con la
funzione \func{acl\_to\_text}, controllando di nuovo se l'operazione ha
successo (\texttt{\small 28-31}) ed uscendo in caso contrario.  Si provvede
infine a stampare la rappresentazione testuale (\texttt{\small 32}) e dopo
aver liberato (\texttt{\small 33-34}) le risorse allocate automaticamente, si
conclude l'esecuzione.


\subsection{La gestione delle quote disco}
\label{sec:disk_quota}

Quella delle quote disco è una funzionalità introdotta inizialmente da BSD e
presente in Linux fino dai kernel dalla serie 2.0, che consente di porre dei
tetti massimi al consumo delle risorse di un filesystem (spazio disco e
\textit{inode}) da parte di utenti e gruppi.

Dato che la funzionalità ha senso solo per i filesystem su cui si mantengono i
dati degli utenti essa deve essere attivata esplicitamente.\footnote{in genere
  la si attiva sul filesystem che contiene le \textit{home} degli utenti, dato
  che non avrebbe senso per i file di sistema che in genere appartengono
  all'amministratore.} Questo si fa, per tutti i filesystem che le supportano,
tramite due distinte opzioni di montaggio, \texttt{usrquota} e
\texttt{grpquota}, che abilitano le quote rispettivamente per gli utenti e per
i gruppi. Così è possibile usare le limitazioni delle quote o sugli utenti o
sui gruppi o su entrambi.

Dal kernel 4.1, ed inizialmente solo per il filesystem XFS, sono diventate
disponibili un terzo tipo di quote, dette \textit{project quota}, che
consentono di applicare delle quote ad un ``\textsl{progetto}'', identificato
come ramo di albero sotto una directory, per il quale possono essere imposti
dei limiti (di nuovo in termini di spazio disco o \textit{inode}) per i file
che ne fanno parte. Si può così porre dei limiti sul contenuto di un ramo di
albero.

Il meccanismo prevede che per ciascun filesystem che supporta le quote disco
(i vari \acr{extN}, \acr{Btrfs}, \acr{XFS}, \acr{JFS}, \acr{ReiserFS}) il
kernel provveda sia a mantenere aggiornati i dati relativi al consumo delle
risorse da parte degli utenti e dei gruppi (e del progetto), che a far
rispettare i limiti imposti dal sistema, con la generazione di un errore di
\errcode{EDQUOT} per tutte le operazioni sui file che porterebbero ad un
superamento degli stessi. Si tenga presente che questi due compiti sono
separati, il primo si attiva al montaggio del filesystem con il supporto per
le quote, il secondo deve essere abilitato esplicitamente.

Per il mantenimento dei dati di consumo delle risorse vengono usati due file
riservati nella directory radice del filesystem su cui si sono attivate le
quote, uno per le quote utente e l'altro per le quote gruppo.\footnote{la cosa
  vale per tutti i filesystem tranne \textit{XFS} che mantiene i dati
  internamente, compresi quelli per le \textit{project quota}, che pertanto,
  essendo questo l'unico filesyste che le supporta, non hanno un file ad esse
  riservato.} Con la versione 2 del supporto delle quote, che da anni è
l'unica rimasta in uso, questi file sono \texttt{aquota.user} e
\texttt{aquota.group}, in precedenza erano \texttt{quota.user} e
\texttt{quota.group}.

Dato che questi file vengono aggiornati soltanto se il filesystem è stato
montato attivando il supporto delle quote, se si abilita il supporto in un
secondo tempo e nel frattempo sono state eseguite delle operazioni sul
filesystem quando il supporto era disabilitato, i dati contenuti possono non
corrispondere esattamente allo stato corrente del consumo delle risorse. Per
questo motivo prima di montare in scrittura un filesystem su cui sono
abilitate le quote viene richiesto di utilizzare il comando \cmd{quotacheck}
per verificare e aggiornare i dati.

Le restrizioni sul consumo delle risorse previste dal sistema delle quote
prevedono sempre la presenza di due diversi limiti, il primo viene detto
\textit{soft limit} e può essere superato per brevi periodi di tempo senza che
causare errori per lo sforamento delle quote, il secondo viene detto
\textit{hard limit} e non può mai essere superato.

Il periodo di tempo per cui è possibile eccedere rispetto alle restrizioni
indicate dal \textit{soft limit} è detto ``\textsl{periodo di grazia}''
(\textit{grace period}), che si attiva non appena si supera la quota da esso
indicata. Se si continua a restare al di sopra del \textit{soft limit} una
volta scaduto il \textit{grace period} questo verrà trattato allo stesso modo
dell'\textit{hard limit} e si avrà l'emissione immediata di un errore.

Si tenga presente infine che entrambi i tipi di limiti (\textit{soft limit} e
\textit{hard limit}) possono essere disposti separatamente su entrambe le
risorse di un filesystem, essi cioè possono essere presenti in maniera
indipendente sia sullo spazio disco, con un massimo per il numero di blocchi,
che sui file, con un massimo per il numero di \textit{inode}.

La funzione di sistema che consente di controllare tutti i vari aspetti della
gestione delle quote è \funcd{quotactl}, ed il suo prototipo è:

\begin{funcproto}{
\fhead{sys/types.h}
\fhead{sys/quota.h}
\fdecl{int quotactl(int cmd, const char *dev, int id, caddr\_t addr)}
\fdesc{Esegue una operazione di controllo sulle quote disco.} 
}

{La funzione ritorna $0$ in caso di successo e $-1$ per un errore, nel qual
  caso \var{errno} assumerà uno dei valori: 
  \begin{errlist}
  \item[\errcode{EACCES}] si è richiesto \const{Q\_QUOTAON}, ma il file delle
    quote indicato da \param{addr} esiste ma non è un file ordinario o no sta
    su \param{dev}.
  \item[\errcode{EBUSY}] si è richiesto \const{Q\_QUOTAON}, ma le quote sono
    già attive.
  \item[\errcode{EFAULT}] \param{addr} o \param{dev} non sono un puntatori
    validi. 
  \item[\errcode{EINVAL}] o \param{cmd} non è un comando valido,
    o il dispositivo \param{dev} non esiste.
  \item[\errcode{EIO}] errore di lettura/scrittura sul file delle quote.
  \item[\errcode{EMFILE}] non si può aprire il file delle quote avendo
    superato il limite sul numero di file aperti nel sistema. 
  \item[\errcode{ENODEV}] \param{dev} non corrisponde ad un \textit{mount
      point} attivo.
  \item[\errcode{ENOPKG}] il kernel è stato compilato senza supporto per le
    quote. 
  \item[\errcode{ENOTBLK}] \param{dev} non è un dispositivo a blocchi.
  \item[\errcode{EPERM}] non si hanno i permessi per l'operazione richiesta.
  \item[\errcode{ESRCH}] è stato richiesto uno fra \const{Q\_GETQUOTA},
    \const{Q\_SETQUOTA}, \const{Q\_SETUSE}, \const{Q\_SETQLIM} per un
    filesystem senza quote attivate.
  \end{errlist}
}
\end{funcproto}

% TODO rivedere gli errori

La funzione richiede che il filesystem sul quale si vuole operare, che deve
essere specificato con il nome del relativo file di dispositivo nell'argomento
\param{dev}, sia montato con il supporto delle quote abilitato. Per le
operazioni che lo richiedono inoltre si dovrà indicare con l'argomento
\param{id} l'utente o il gruppo o il progetto (specificati rispettivamente per
\ids{UID}, \ids{GID} o identificativo) su cui si vuole operare,\footnote{nel
  caso di \textit{project quota} gli identificativi vengono associati alla
  directory base del progetto nel file \file{/etc/projects}, ed impostati con
  \cmd{xfs\_quota}, l'argomento è di natura sistemistica e va al di là dello
  scopo di questo testo.} o altri dati relativi all'operazione. Alcune
operazioni più complesse usano infine l'argomento \param{addr} per indicare un
indirizzo ad un area di memoria il cui utilizzo dipende dall'operazione
stessa.

La funzione prevede la possibilità di eseguire una serie operazioni sulle
quote molto diverse fra loro, la scelta viene effettuata tramite il primo
argomento, \param{cmd}, che però oltre all'operazione indica anche a quale
tipo di quota (utente o gruppo) l'operazione deve applicarsi. Per questo il
valore di questo argomento viene costruito con l'ausilio della di una apposita
macro \macro{QCMD}:

{\centering
\vspace{3pt}
\begin{funcbox}{
\fhead{sys/quota.h}
\fdecl{int \macrod{QCMD}(subcmd,type)}
\fdesc{Imposta il comando \param{subcmd} per il tipo di quote (utente o
  gruppo) \param{type}.}
} 
\end{funcbox}
}

La macro consente di specificare, oltre al tipo di operazione, da indicare con
l'argomento \param{subcmd}, se questa deve applicarsi alle quote utente o alle
quote gruppo o alle quote progetto. Questo viene indicato dall'argomento
\param{type} che deve essere sempre definito ed assegnato ad uno fra i valori
\const{USRQUOTA}, \const{GRPQUOTA} e \const{PRJQUOTA}.

\begin{table}[htb]
  \centering
  \footnotesize
  \begin{tabular}{|l|p{10cm}|}
    \hline
    \textbf{Comando} & \textbf{Descrizione} \\
    \hline
    \hline
    \constd{Q\_QUOTAON} & Attiva l'applicazione delle quote disco per il
                          filesystem indicato da \param{dev}, si deve passare
                          in \param{addr} il \textit{pathname} al file che
                          mantiene le quote, che deve esistere, e \param{id}
                          deve indicare la versione del formato con uno dei
                          valori di tab.~\ref{tab:quotactl_id_format};
                          l'operazione richiede i privilegi di
                          amministratore.\\
    \constd{Q\_QUOTAOFF}& Disattiva l'applicazione delle quote disco per il
                          filesystem indicato da \param{dev}, \param{id}
                          e \param{addr} vengono ignorati; l'operazione
                          richiede i privilegi di amministratore.\\  
    \constd{Q\_GETQUOTA}& Legge i limiti ed i valori correnti delle quote nel
                          filesystem indicato da \param{dev} per l'utente o
                          il gruppo specificato da \param{id}; si devono avere
                          i privilegi di amministratore per leggere i dati
                          relativi ad altri utenti o a gruppi di cui non si fa
                          parte, il risultato viene restituito in una struttura
                          \struct{dqblk} all'indirizzo indicato
                          da \param{addr}.\\
    \constd{Q\_SETQUOTA}& Imposta i limiti per le quote nel filesystem
                          indicato da \param{dev} per l'utente o il gruppo
                          specificato da \param{id} secondo i valori ottenuti
                          dalla struttura \struct{dqblk} puntata
                          da \param{addr}; l'operazione richiede i privilegi
                          di amministratore.\\ 
    \constd{Q\_GETINFO} & Legge le informazioni (in sostanza i \textit{grace
                            time}) delle quote del filesystem indicato
                          da \param{dev} sulla struttura \struct{dqinfo} 
                          puntata da \param{addr}, \param{id} viene ignorato.\\
    \constd{Q\_SETINFO} & Imposta le informazioni delle quote del filesystem
                          indicato da \param{dev} come ottenuti dalla
                          struttura \struct{dqinfo} puntata
                          da \param{addr}, \param{id} viene ignorato;  
                          l'operazione richiede i privilegi di amministratore.\\
    \constd{Q\_GETFMT}  & Richiede il valore identificativo (quello di
                          tab.~\ref{tab:quotactl_id_format}) per il formato
                          delle quote attualmente in uso sul filesystem
                          indicato da \param{dev}, che sarà memorizzato
                          sul buffer di 4 byte puntato da \param{addr}.\\
    \constd{Q\_SYNC}    & Aggiorna la copia su disco dei dati delle quote del
                          filesystem indicato da \param{dev}; in questo
                          caso \param{dev} può anche essere \val{NULL} nel
                          qual caso verranno aggiornati i dati per tutti i
                          filesystem con quote attive, \param{id}
                          e \param{addr} vengono comunque ignorati.\\ 
    \constd{Q\_GETSTATS}& Ottiene statistiche ed altre informazioni generali 
                          relative al sistema delle quote per il filesystem
                          indicato da \param{dev}, richiede che si
                          passi come argomento \param{addr} l'indirizzo di una
                          struttura \struct{dqstats}, mentre i valori
                          di \param{id} e \param{dev} vengono ignorati;
                          l'operazione è obsoleta e non supportata nei kernel
                          più recenti, che espongono la stessa informazione
                          nei file sotto \procfile{/proc/self/fs/quota/}.\\
%    \const{} & .\\
    \hline
  \end{tabular}
  \caption{Possibili valori per l'argomento \param{subcmd} di
    \macro{QCMD}.} 
  \label{tab:quotactl_commands}
\end{table}

I possibili valori per l'argomento \param{subcmd} di \macro{QCMD} sono
riportati in tab.~\ref{tab:quotactl_commands}, che illustra brevemente il
significato delle operazioni associate a ciascuno di essi. In generale le
operazioni di attivazione, disattivazione e di modifica dei limiti delle quote
sono riservate e richiedono i privilegi di amministratore.\footnote{per essere
  precisi tutte le operazioni indicate come privilegiate in
  tab.~\ref{tab:quotactl_commands} richiedono la capacità
  \const{CAP\_SYS\_ADMIN}.} Inoltre gli utenti possono soltanto richiedere i
dati relativi alle proprie quote, solo l'amministratore può ottenere i dati di
tutti.


Alcune delle operazioni di tab.~\ref{tab:quotactl_commands} sono alquanto
complesse e richiedono un approfondimento maggiore. Le due più rilevanti sono
probabilmente \const{Q\_GETQUOTA} e \const{Q\_SETQUOTA}, che consentono la
gestione dei limiti delle quote. Entrambe fanno riferimento ad una specifica
struttura \struct{dqblk}, la cui definizione è riportata in
fig.~\ref{fig:dqblk_struct},\footnote{la definizione mostrata è quella usata
  fino dal kernel 2.4.22, non prenderemo in considerazione le versioni
  obsolete.} nella quale vengono inseriti i dati relativi alle quote di un
singolo utente o gruppo.

\begin{figure}[!htb]
  \footnotesize \centering
  \begin{minipage}[c]{0.95\textwidth}
    \includestruct{listati/dqblk.h}
  \end{minipage} 
  \normalsize 
  \caption{La struttura \structd{dqblk} per i dati delle quote disco.}
  \label{fig:dqblk_struct}
\end{figure}

La struttura \struct{dqblk} viene usata sia con \const{Q\_GETQUOTA} per
ottenere i valori correnti dei limiti e dell'occupazione delle risorse, che
con \const{Q\_SETQUOTA} per effettuare modifiche ai limiti. Come si può notare
ci sono alcuni campi (in sostanza \val{dqb\_curspace}, \val{dqb\_curinodes},
\val{dqb\_btime}, \val{dqb\_itime}) che hanno senso solo in lettura, in quanto
riportano uno stato non modificabile da \func{quotactl} come l'uso corrente di
spazio disco ed \textit{inode}, o il tempo che resta nel caso si sia superato
un \textit{soft limit}.

Inoltre in caso di modifica di un limite si può voler operare solo su una
delle risorse (blocchi o \textit{inode}),\footnote{non è possibile modificare
  soltanto uno dei limiti (\textit{hard} o \textit{soft}) occorre sempre
  rispecificarli entrambi.} per questo la struttura prevede un campo apposito,
\val{dqb\_valid}, il cui scopo è quello di indicare quali sono gli altri campi
che devono essere considerati validi. Questo campo è una maschera binaria che
deve essere espressa nei termini di OR aritmetico delle apposite costanti di
tab.~\ref{tab:quotactl_qif_const}, dove si è riportato il significato di
ciascuna di esse ed i campi a cui fanno riferimento.

\begin{table}[!htb]
  \centering
  \footnotesize
  \begin{tabular}{|l|p{10cm}|}
    \hline
    \textbf{Costante} & \textbf{Descrizione} \\
    \hline
    \hline
    \constd{QIF\_BLIMITS}& Limiti sui blocchi di spazio disco
                           (\val{dqb\_bhardlimit} e \val{dqb\_bsoftlimit}).\\
    \constd{QIF\_SPACE}  & Uso corrente dello spazio disco
                           (\val{dqb\_curspace}).\\
    \constd{QIF\_ILIMITS}& Limiti sugli \textit{inode}
                           (\val{dqb\_ihardlimit} e \val{dqb\_isoftlimit}).\\
    \constd{QIF\_INODES} & Uso corrente degli \textit{inode}
                           (\val{dqb\_curinodes}).\\
    \constd{QIF\_BTIME}  & Tempo di sforamento del \textit{soft limit} sul
                           numero di blocchi (\val{dqb\_btime}).\\
    \constd{QIF\_ITIME}  & Tempo di sforamento del \textit{soft limit} sul
                           numero di \textit{inode} (\val{dqb\_itime}).\\ 
    \constd{QIF\_LIMITS} & L'insieme di \const{QIF\_BLIMITS} e
                           \const{QIF\_ILIMITS}.\\
    \constd{QIF\_USAGE}  & L'insieme di \const{QIF\_SPACE} e
                           \const{QIF\_INODES}.\\
    \constd{QIF\_TIMES}  & L'insieme di \const{QIF\_BTIME} e
                           \const{QIF\_ITIME}.\\ 
    \constd{QIF\_ALL}    & Tutti i precedenti.\\
    \hline
  \end{tabular}
  \caption{Costanti per il campo \val{dqb\_valid} di \struct{dqblk}.} 
  \label{tab:quotactl_qif_const}
\end{table}

In lettura con \const{Q\_SETQUOTA} eventuali valori presenti in \struct{dqblk}
vengono comunque ignorati, al momento la funzione sovrascrive tutti i campi
che restituisce e li marca come validi in \val{dqb\_valid}. Si possono invece
usare \const{QIF\_BLIMITS} o \const{QIF\_ILIMITS} per richiedere di impostare
solo la rispettiva tipologia di limiti con \const{Q\_SETQUOTA}. Si tenga
presente che il sistema delle quote richiede che l'occupazione di spazio disco
sia indicata in termini di blocchi e non di byte, dato che la dimensione dei
blocchi dipende da come si è creato il filesystem potrà essere necessario
effettuare qualche conversione per avere un valore in byte.\footnote{in genere
  viene usato un default di 1024 byte per blocco, ma quando si hanno file di
  dimensioni medie maggiori può convenire usare valori più alti per ottenere
  prestazioni migliori in conseguenza di un minore frazionamento dei dati e di
  indici più corti.}

Come accennato la realizzazione delle quote disco ha visto diverse revisioni,
con modifiche sia del formato delle stesse che dei nomi dei file
utilizzati. Per questo alcune operazioni di gestione (in particolare
\const{Q\_QUOTAON} e \const{Q\_GETFMT}) e possono fare riferimento a queste
versioni, che vengono identificate tramite le costanti di
tab.~\ref{tab:quotactl_id_format}.

\begin{table}[htb]
  \centering
  \footnotesize
  \begin{tabular}{|l|p{10cm}|}
    \hline
    \textbf{Identificatore} & \textbf{Descrizione} \\
    \hline
    \hline
    \constd{QFMT\_VFS\_OLD}& Il vecchio (ed obsoleto) formato delle quote.\\
    \constd{QFMT\_VFS\_V0} & La versione 0 usata dal VFS di Linux, supporta
                             \ids{UID} e \ids{GID} a 32 bit e limiti fino a
                             $2^{42}$ byte e $2^{32}$ file.\\
    \constd{QFMT\_VFS\_V1} & La versione 1 usata dal VFS di Linux, supporta
                             \ids{UID} e \ids{GID} a 32 bit e limiti fino a
                             $2^{64}$ byte e $2^{64}$ file.\\
    \hline
  \end{tabular}
  \caption{Valori di identificazione del formato delle quote.} 
  \label{tab:quotactl_id_format}
\end{table}

Altre due operazioni che necessitano di ulteriori spiegazioni sono
\const{Q\_GETINFO} e \const{Q\_SETINFO}, che consentono di ottenere i dati
relativi alle impostazioni delle altre proprietà delle quote, che al momento
sono solo la durata del \textit{grace time} per i due tipi di limiti. Queste
sono due proprietà generali identiche per tutti gli utenti (e i gruppi), per
cui viene usata una operazione distinta dalle precedenti. Anche in questo caso
le due operazioni richiedono l'uso di una apposita struttura \struct{dqinfo},
la cui definizione è riportata in fig.~\ref{fig:dqinfo_struct}.

\begin{figure}[!htb]
  \footnotesize \centering
  \begin{minipage}[c]{0.8\textwidth}
    \includestruct{listati/dqinfo.h}
  \end{minipage} 
  \normalsize 
  \caption{La struttura \structd{dqinfo} per i dati delle quote disco.}
  \label{fig:dqinfo_struct}
\end{figure}

Come per \struct{dqblk} anche in questo caso viene usato un campo della
struttura, \val{dqi\_valid} come maschera binaria per dichiarare quale degli
altri campi sono validi; le costanti usate per comporre questo valore sono
riportate in tab.~\ref{tab:quotactl_iif_const} dove si è riportato il
significato di ciascuna di esse ed i campi a cui fanno riferimento.

\begin{table}[htb]
  \centering
  \footnotesize
  \begin{tabular}{|l|l|}
    \hline
    \textbf{Costante} & \textbf{Descrizione} \\
    \hline
    \hline
    \constd{IIF\_BGRACE}& Il \textit{grace period} per i blocchi
                         (\val{dqi\_bgrace}).\\
    \constd{IIF\_IGRACE}& Il \textit{grace period} per gli \textit{inode} 
                         (\val{dqi\_igrace}).\\ 
    \constd{IIF\_FLAGS} & I flag delle quote (\val{dqi\_flags}) (inusato ?).\\
    \constd{IIF\_ALL}   & Tutti i precedenti.\\
    \hline
  \end{tabular}
  \caption{Costanti per il campo \val{dqi\_valid} di \struct{dqinfo}.} 
  \label{tab:quotactl_iif_const}
\end{table}

Come in precedenza con \const{Q\_GETINFO} tutti i valori vengono letti
sovrascrivendo il contenuto di \struct{dqinfo} e marcati come validi in
\val{dqi\_valid}. In scrittura con \const{Q\_SETINFO} si può scegliere quali
impostare, si tenga presente che i tempi dei campi \val{dqi\_bgrace} e
\val{dqi\_igrace} devono essere specificati in secondi.

Come esempi dell'uso di \func{quotactl} utilizzeremo estratti del codice di un
modulo Python usato per fornire una interfaccia diretta a \func{quotactl}
senza dover passare dalla scansione dei risultati di un comando. Il modulo si
trova fra i pacchetti Debian messi a disposizione da Truelite Srl,
all'indirizzo \url{http://labs.truelite.it/projects/packages}.\footnote{in
  particolare il codice C del modulo è nel file \texttt{quotamodule.c}
  visionabile a partire dall'indirizzo indicato nella sezione
  \textit{Repository}.}

\begin{figure}[!htbp]
  \footnotesize \centering
  \begin{minipage}[c]{\codesamplewidth}
    \includecodesample{listati/get_quota.c}
  \end{minipage}
  \caption{Esempio di codice per ottenere i dati delle quote.} 
  \label{fig:get_quota}
\end{figure}

Il primo esempio, riportato in fig.~\ref{fig:get_quota}, riporta il codice
della funzione che consente di leggere le quote. La funzione fa uso
dell'interfaccia dal C verso Python, che definisce i vari simboli \texttt{Py*}
(tipi di dato e funzioni). Non staremo ad approfondire i dettagli di questa
interfaccia, per la quale esistono numerose trattazioni dettagliate, ci
interessa solo esaminare l'uso di \func{quotactl}. 

In questo caso la funzione prende come argomenti (\texttt{\small 1}) l'intero
\texttt{who} che indica se si vuole operare sulle quote utente o gruppo,
l'identificatore \texttt{id} dell'utente o del gruppo scelto, ed il nome del
file di dispositivo del filesystem su cui si sono attivate le
quote.\footnote{questi vengono passati come argomenti dalle funzioni mappate
  come interfaccia pubblica del modulo (una per gruppi ed una per gli utenti)
  che si incaricano di decodificare i dati passati da una chiamata nel codice
  Python.} Questi argomenti vengono passati direttamente alla chiamata a
\func{quotactl} (\texttt{\small 5}), a parte \texttt{who} che viene abbinato
con \macro{QCMD} al comando \const{Q\_GETQUOTA} per ottenere i dati.

La funzione viene eseguita all'interno di un condizionale (\texttt{\small
  5-16}) che in caso di successo provvede a costruire (\texttt{\small 6-12})
opportunamente una risposta restituendo tramite la opportuna funzione di
interfaccia un oggetto Python contenente i dati della struttura \struct{dqblk}
relativi a uso corrente e limiti sia per i blocchi che per gli
\textit{inode}. In caso di errore (\texttt{\small 13-15}) si usa un'altra
funzione dell'interfaccia per passare il valore di \var{errno} come eccezione.

\begin{figure}[!htbp]
  \footnotesize \centering
  \begin{minipage}[c]{\codesamplewidth}
    \includecodesample{listati/set_block_quota.c}
  \end{minipage}
  \caption{Esempio di codice per impostare i limiti sullo spazio disco.}
  \label{fig:set_block_quota}
\end{figure}

Per impostare i limiti sullo spazio disco si potrà usare una seconda funzione,
riportata in fig.~\ref{fig:set_block_quota}, che prende gli stessi argomenti
della precedente, con lo stesso significato, a cui si aggiungono i valori per
il \textit{soft limit} e l'\textit{hard limit}. In questo caso occorrerà,
prima di chiamare \func{quotactl}, inizializzare opportunamente
(\texttt{\small 5-7}) i campi della struttura \struct{dqblk} che si vogliono
utilizzare (quelli relativi ai limiti sui blocchi) e specificare gli stessi
con \const{QIF\_BLIMITS} in \var{dq.dqb\_valid}. 

Fatto questo la chiamata a \func{quotactl}, stavolta con il comando
\const{Q\_SETQUOTA}, viene eseguita come in precedenza all'interno di un
condizionale (\texttt{\small 9-14}). In questo caso non essendovi da
restituire nessun dato in caso di successo si usa (\texttt{\small 10}) una
apposita funzione di uscita, mentre si restituisce come prima una eccezione
con il valore di \var{errno} in caso di errore (\texttt{\small 12-13}).

\subsection{La gestione dei {chroot}}
\label{sec:file_chroot}

% TODO: valutare se introdurre una nuova sezione sulle funzionalità di
% sicurezza avanzate, con dentro chroot SELinux e AppArmor, Tomoyo, Smack,
% cgroup o altro

% TODO: spostare chroot e le funzioni affini relative ai container da qualche
% parte diversa se è il caso. 

% TODO Inheriting capabilities vedi http://lwn.net/Articles/632520/ eambient
% capabilities introdotte con il kernel 4.3, vedi 
% http://git.kernel.org/cgit/linux/kernel/git/torvalds/linux.git/commit/?id=58319057b7847667f0c9585b9de0e8932b0fdb08

Benché non abbia niente a che fare con permessi, utenti e gruppi, la funzione
\func{chroot} viene usata spesso per restringere le capacità di accesso di un
programma ad una sezione limitata del filesystem, per cui ne parleremo in
questa sezione.

Come accennato in sez.~\ref{sec:proc_fork} ogni processo oltre ad una
directory di lavoro, ha anche una directory \textsl{radice}\footnote{entrambe
  sono contenute in due campi (rispettivamente \var{pwd} e \var{root}) di
  \kstruct{fs\_struct}; vedi fig.~\ref{fig:proc_task_struct}.} che, pur
essendo di norma corrispondente alla radice dell'albero dei file dell'intero
sistema, ha per il processo il significato specifico di directory rispetto
alla quale vengono risolti i \textit{pathname} assoluti.\footnote{cioè quando
  un processo chiede la risoluzione di un \textit{pathname}, il kernel usa
  sempre questa directory come punto di partenza.} Il fatto che questo valore
sia specificato per ogni processo apre allora la possibilità di modificare le
modalità di risoluzione dei \textit{pathname} assoluti da parte di un processo
cambiando questa directory, così come si fa coi \textit{pathname} relativi
cambiando la directory di lavoro.

Normalmente la directory radice di un processo coincide con la radice generica
dell'albero dei file, che è la directory che viene montata direttamente dal
kernel all'avvio secondo quanto illustrato in sez.~\ref{sec:file_pathname}.
Questo avviene perché, come visto in sez.~\ref{cha:process_handling} la
directory radice di un processo viene ereditata dal padre attraverso una
\func{fork} e mantenuta attraverso una \func{exec}, e siccome tutti i processi
derivano da \cmd{init}, che ha come radice quella montata dal kernel, questa
verrà mantenuta.

In certe situazioni però è utile poter impedire che un processo possa accedere
a tutto l'albero dei file iniziale; per far questo si può cambiare la sua
directory radice con la funzione di sistema \funcd{chroot}, il cui prototipo
è:

\begin{funcproto}{
\fhead{unistd.h}
\fdecl{int chroot(const char *path)}
\fdesc{Cambia la directory radice del processo.} 
}

{La funzione ritorna $0$ in caso di successo e $-1$ per un errore, nel qual
  caso \var{errno} assumerà uno dei valori: 
  \begin{errlist}
  \item[\errcode{EPERM}] non si hanno i privilegi di amministratore.
  \end{errlist}
  ed inoltre \errval{EFAULT}, \errval{ENAMETOOLONG}, \errval{ENOENT},
  \errval{ENOMEM}, \errval{ENOTDIR}, \errval{EACCES}, \errval{ELOOP};
  \errval{EROFS} e \errval{EIO} nel loro significato generico.}
\end{funcproto}

La funzione imposta la directory radice del processo a quella specificata da
\param{path} (che ovviamente deve esistere) ed ogni \textit{pathname} assoluto
usato dalle funzioni chiamate nel processo sarà risolto a partire da essa,
rendendo impossibile accedere alla parte di albero sovrastante. Si ha così
quella che viene chiamata una \textit{chroot jail}, in quanto il processo non
può più accedere a file al di fuori della sezione di albero in cui è stato
\textsl{imprigionato}.

Solo un processo con i privilegi di amministratore può usare questa
funzione,\footnote{più precisamente se possiede la capacità
  \const{CAP\_SYS\_CHROOT}.} e la nuova radice, per quanto detto in
sez.~\ref{sec:proc_fork}, sarà ereditata da tutti i suoi processi figli. Si
tenga presente però che la funzione non cambia la directory di lavoro del
processo, che potrebbe restare fuori dalla \textit{chroot jail}.

Questo è il motivo per cui la funzione è efficace nel restringere un processo
ad un ramo di albero solo se dopo averla eseguita si cedono i privilegi di
amministratore. Infatti se per un qualunque motivo il processo resta con la
sua directory di lavoro al di fuori dalla \textit{chroot jail}, potrà accedere
a tutto il resto del filesystem usando dei \textit{pathname} relativi, dato
che in tal caso è possibile, grazie all'uso di ``\texttt{..}'', risalire
all'indietro fino alla radice effettiva dell'albero dei file.

Potrebbe sembrare che per risolvere il problema sia sufficiente ricordarsi di
eseguire preventivamente anche una \func{chdir} sulla directory su cui si
andrà ad eseguire \func{chroot}, così da assicurarsi che le directory di
lavoro sia all'interno della \textit{chroot jail}.  Ma se ad un processo
restano i privilegi di amministratore esso potrà comunque portare la sua
directory di lavoro fuori dalla \textit{chroot jail} in cui si trova. Basterà
infatti eseguire di nuovo \func{chroot} su una qualunque directory contenuta
nell'attuale directory di lavoro perché quest'ultima risulti al di fuori della
nuova \textit{chroot jail}.  Per questo motivo l'uso di questa funzione non ha
molto senso quando un processo di cui si vuole limitare l'accesso necessita
comunque dei privilegi di amministratore per le sue normali operazioni.

Nonostante queste limitazioni la funzione risulta utile qualora la si possa
applicare correttamente cedendo completamente i privilegi di amministratore
una volta eseguita.  Ed esempio caso tipico di uso di \func{chroot} è quello
di un server FTP anonimo in si vuole che il server veda solo i file che deve
trasferire. In tal caso si esegue una \func{chroot} sulla directory che
contiene i file, che il server dovrà in grado di leggere come utente
ordinario, e poi si cedono tutti i privilegi di amministratore.  Si tenga
presente però che in casi come questo occorrerà fornire all'interno della
\textit{chroot jail} un accesso anche a tutti i file (in genere programmi e
librerie) di cui il server potrebbe avere bisogno.


% LocalWords:  sez like filesystem unlink MacOS Windows VMS inode kernel unistd
% LocalWords:  int const char oldpath newpath errno EXDEV EPERM st Smack SysV
% LocalWords:  EEXIST EMLINK EACCES ENAMETOOLONG ENOTDIR EFAULT ENOMEM EROFS ls
% LocalWords:  ELOOP ENOSPC EIO pathname nlink stat vfat fsck EISDIR ENOENT cap
% LocalWords:  POSIX socket fifo sticky root system call count crash init linux
% LocalWords:  descriptor remove rename rmdir stdio glibc libc NFS DT obj dup
% LocalWords:  ENOTEMPTY EBUSY mount point EINVAL soft symbolic tab symlink fig
% LocalWords:  dangling access chdir chmod chown creat exec lchown lstat mkdir
% LocalWords:  mkfifo mknod opendir pathconf readlink truncate path buff size
% LocalWords:  grub bootloader grep MAXSYMLINKS cat VFS sys dirname fcntl tv Py
% LocalWords:  dev umask IFREG IFBLK IFCHR IFIFO SVr sgid BSD SVID NULL from to
% LocalWords:  stream dirent EMFILE ENFILE dirfd SOURCE fchdir readdir struct
% LocalWords:  EBADF namlen HAVE thread entry result value argument fileno ext
% LocalWords:  name TYPE OFF RECLEN UNKNOWN REG SOCK CHR BLK type IFTODT DTTOIF
% LocalWords:  DTYPE off reclen seekdir telldir void rewinddir closedir select
% LocalWords:  namelist compar malloc qsort alphasort versionsort strcoll myls
% LocalWords:  strcmp direntry while current working home shell pwd get stddef
% LocalWords:  getcwd ERANGE getwd change fd race condition tmpnam getfacl mark
% LocalWords:  string tmpdir TMP tempnam pfx TMPNAME suid tmp EXCL tmpfile pid
% LocalWords:  EINTR mktemp mkstemp stlib template filename XXXXXX OpenBSD buf
% LocalWords:  mkdtemp fstat filedes nell'header padding ISREG ISDIR ISCHR IFMT
% LocalWords:  ISBLK ISFIFO ISLNK ISSOCK IFSOCK IFLNK IFDIR ISUID UID ISGID GID
% LocalWords:  ISVTX IRUSR IWUSR IXUSR IRGRP IWGRP IXGRP IROTH IWOTH IXOTH  OLD
% LocalWords:  blocks blksize holes lseek TRUNC ftruncate ETXTBSY length QCMD
% LocalWords:  hole atime read utime mtime write ctime modification leafnode cp
% LocalWords:  make fchmod fchown utimbuf times actime modtime Mac owner uid fs
% LocalWords:  gid Control List patch mandatory control execute group other all
% LocalWords:  effective passwd IGID locking swap saved text IRWXU IRWXG subcmd
% LocalWords:  IRWXO capability FSETID mask capabilities chroot jail QUOTAOFF
% LocalWords:  FTP filter Attributes Solaris FreeBSD libacl hash at dqblk SYNC
% LocalWords:  XFS SELinux namespace attribute security trusted Draft Modules
% LocalWords:  attributes mime ADMIN FOWNER libattr lattr getxattr lgetxattr of
% LocalWords:  fgetxattr attr ssize ENOATTR ENOTSUP NUL setxattr lsetxattr list
% LocalWords:  fsetxattr flags XATTR REPLACE listxattr llistxattr flistxattr by
% LocalWords:  removexattr lremovexattr fremovexattr attributename acl GETINFO
% LocalWords:  OBJ setfacl len any prefix separator options NUMERIC IDS SMART
% LocalWords:  INDENT major number IDE Documentation makedev proc copy LNK long
% LocalWords:  euidaccess eaccess delete def tag qualifier permset calendar NOW
% LocalWords:  mutt noatime relatime strictatime atim nsec mtim ctim atimensec
% LocalWords:  mtimensec utimes timeval futimes lutimes ENOSYS futimens OMIT PR
% LocalWords:  utimensat timespec sec futimesat LIDS DAC OVERRIDE SEARCH chattr
% LocalWords:  Discrectionary KILL SETGID domain SETUID setuid setreuid SETPCAP
% LocalWords:  setresuid setfsuid IMMUTABLE immutable append only BIND SERVICE
% LocalWords:  BROADCAST broadcast multicast multicasting RAW PACKET IPC LOCK
% LocalWords:  memory mlock mlockall shmctl mmap MODULE RAWIO ioperm iopl PACCT
% LocalWords:  ptrace accounting NICE RESOURCE TTY CONFIG hangup vhangup dell'
% LocalWords:  LEASE lease SETFCAP AUDIT permitted inherited inheritable AND nn
% LocalWords:  bounding execve fork capget capset header hdrp datap ESRCH undef
% LocalWords:  version libcap lcap clear ncap caps pag capgetp CapInh CapPrm RT
% LocalWords:  fffffeff CapEff getcap scheduling lookup  dqinfo SETINFO GETFMT
% LocalWords:  NEWNS unshare nice NUMA ioctl journaling close XOPEN fdopendir
% LocalWords:  Btrfs mkostemp extN ReiserFS JFS Posix usrquota grpquota EDQUOT
% LocalWords:  aquota quotacheck limit grace period quotactl cmd caddr addr dqb
% LocalWords:  QUOTAON ENODEV ENOPKG ENOTBLK GETQUOTA SETQUOTA SETUSE SETQLIM
% LocalWords:  forced allowed sendmail SYSLOG WAKE ALARM CLOCK BOOTTIME dqstats
% LocalWords:  REALTIME securebits GETSTATS QFMT curspace curinodes btime itime
% LocalWords:  QIF BLIMITS bhardlimit bsoftlimit ILIMITS ihardlimit isoftlimit
% LocalWords:  INODES LIMITS USAGE valid dqi IIF BGRACE bgrace IGRACE igrace is
% LocalWords:  Python Truelite Srl quotamodule Repository who nell' dall' KEEP
% LocalWords:  SECURE KEEPCAPS prctl FIXUP NOROOT LOCKED dell'IPC dell'I IOPRIO
% LocalWords:  CAPBSET CLASS IDLE dcookie overflow DIFFERS Virtual everything
% LocalWords:  dentry register resolution cache dcache operation llseek poll ln
% LocalWords:  multiplexing fsync fasync seek block superblock gapil tex img du
% LocalWords:  second linked journaled source filesystemtype unsigned device
% LocalWords:  mountflags NODEV ENXIO dummy devfs magic MGC RDONLY NOSUID MOVE
% LocalWords:  NOEXEC SYNCHRONOUS REMOUNT MANDLOCK NODIRATIME umount MNT statfs
% LocalWords:  fstatfs fstab mntent ino bound orig new setpcap metadati sysfs
% LocalWords:  bind DIRSYNC lsattr Hierarchy FHS SHARED UNBINDABLE shared REC
% LocalWords:  subtree SILENT log unbindable BUSY EAGAIN EXPIRE DETACH NOFOLLOW
% LocalWords:  lazy encfs sshfs setfsent getfsent getfsfile getfsspec endfsent
% LocalWords:  setmntent getmntent addmntent endmntent hasmntopt such offsetof
% LocalWords:  member scan attack EOVERFLOW BITS blkcnt rdev FDCWD functions
% LocalWords:  faccessat grpid lacl AppArmor capsetp mygetfacl table Tb MSK
%  LocalWords:  LAZYTIME submount peer protected hardlink symlinks silly RDWR
%  LocalWords:  renames unreachable CLOEXEC mkstemps mkostemps suffixlen Aug
%  LocalWords:  prefissoXXXXXXsuffisso nell'I fstatat statx sull' drwxrwxrwt
%  LocalWords:  Disalloca

%%% Local Variables: 
%%% mode: latex
%%% TeX-master: "gapil"
%%% End: 
